% !TeX program = xelatex
% Standalone English volume. Use XeLaTeX until the references stabilize.
\documentclass[a4paper,11pt]{article}
\usepackage[T1]{fontenc}
\usepackage{lmodern}
\usepackage[english]{babel}
\usepackage[left=26mm,right=26mm,top=25mm,bottom=25mm,headheight=15pt,headsep=7mm,footskip=11mm]{geometry}
\usepackage{amsmath,amssymb,mathtools,amsthm}
\newtheorem{proposition}{Proposition}[section]
\newtheorem{lemma}[proposition]{Lemma}
\newtheorem{theorem}[proposition]{Theorem}
\newtheorem{corollary}[proposition]{Corollary}
\newtheorem{maintheorem}[proposition]{Main Theorem}
\theoremstyle{definition}
\newtheorem{definition}[proposition]{Definition}
\theoremstyle{remark}
\newtheorem{remark}[proposition]{Remark}
\makeatletter
\renewcommand*\l@section{\@dottedtocline{1}{0em}{2.4em}}
\renewcommand*\l@subsection{\@dottedtocline{2}{1.45em}{3.45em}}
\renewcommand*\l@subsubsection{\@dottedtocline{3}{3.8em}{4.55em}}
\makeatother
\usepackage{xcolor,fancyhdr,titlesec,needspace,xurl}
\usepackage{longtable,multicol}
\usepackage{hyperref}
\definecolor{ink}{HTML}{1A3542}
\definecolor{muted}{HTML}{4F606A}
\hypersetup{unicode=true,colorlinks=true,linkcolor=ink,filecolor=ink,urlcolor=ink,citecolor=ink,bookmarksnumbered=true,
 pdfdisplaydoctitle=true,pdflang={en},
 pdftitle={The generalized Riemann hypothesis for Dirichlet L-functions. Edition 930, Volume VIII},
 pdfsubject={A natural-language reconstruction of the proposed proof},
 pdfauthor={A mathematical reconstruction}}
\titleformat{\section}{\Large\bfseries\color{ink}}{\thesection}{.65em}{}
\titlespacing*{\section}{0pt}{.8em}{.65em}
\setlength{\parindent}{1.2em}
\linespread{1.045}
\setlength{\parskip}{.22em plus .05em minus .03em}
\setlength{\emergencystretch}{2em}
\setcounter{topnumber}{2}
\setcounter{tocdepth}{3}\setcounter{secnumdepth}{3}
\widowpenalty=4500
\clubpenalty=4500
\pagestyle{fancy}\fancyhf{}
\fancyhead[L]{\small\color{muted}Dirichlet GRH \textperiodcentered\ Volume I}
\fancyhead[R]{\small\color{muted}Dirichlet zeros and power intervals}
\fancyfoot[C]{\small\thepage}
\renewcommand{\headrulewidth}{.3pt}
\newcommand{\lead}[1]{\par\smallskip\noindent{\bfseries\color{ink}#1.}\enspace}
\newcommand{\sour}[2]{\unskip\penalty150\hspace{.35em}{\small\hyperlink{ref#1}{[#2]}}}
\titleformat{\subsection}{\large\bfseries}{\thesubsection}{.6em}{}
\titlespacing*{\subsection}{0pt}{1.15em plus .15em}{.4em}
\setlength{\abovedisplayskip}{9pt plus 2pt minus 3pt}
\setlength{\belowdisplayskip}{9pt plus 2pt minus 3pt}
\newcommand{\RR}{\mathbb R}\newcommand{\CC}{\mathbb C}\newcommand{\QQ}{\mathbb Q}
\newcommand{\ZZ}{\mathbb Z}\newcommand{\ord}{\operatorname{ord}}
\newcommand{\Frob}{\operatorname{Frob}}\newcommand{\Gal}{\operatorname{Gal}}
\newcommand{\len}{\operatorname{length}}\newcommand{\Res}{\operatorname*{Res}}
\newcommand{\wmean}{\operatorname{av}}\newcommand{\Ker}{\operatorname{ker}}
\newcommand{\Lh}{\mathcal L}\newcommand{\Gg}{\mathcal G}
\newcommand{\real}{\operatorname{Re}}\newcommand{\imag}{\operatorname{Im}}
\newcommand{\abs}[1]{\left|#1\right|}\newcommand{\eps}{\varepsilon}
\newcommand{\anchoritem}[2]{\par\Needspace{2.8\baselineskip}\hypertarget{ref#1}{}\noindent{\bfseries\color{ink}#2.}\enspace}
\newcommand{\srcloc}[2]{{\small\nolinkurl{#1.lean}}, line~#2}
\numberwithin{equation}{section}
\raggedbottom

\newcommand{\Sel}{\mathcal S}
\newcommand{\Cl}{\operatorname{Cl}}
\newcommand{\Aut}{\operatorname{Aut}}
\newcommand{\GL}{\operatorname{GL}}
\newcommand{\OO}{\mathcal O}
\hypersetup{bookmarksdepth=3}


\addto\captionsenglish{\renewcommand{\partname}{Volume}}


\makeatletter
\expandafter\def\csname r@cor:bfirst368\endcsname{{37.21}{55}{}{proposition.37.21}{Dirichlet-GRH-v930-Volume-VII.pdf}}
\expandafter\def\csname r@cor:componentdefect375\endcsname{{38.12}{78}{}{proposition.38.12}{Dirichlet-GRH-v930-Volume-VII.pdf}}
\expandafter\def\csname r@cor:corner120\endcsname{{33.9}{56}{}{proposition.33.9}{Dirichlet-GRH-v930-Volume-VI.pdf}}
\expandafter\def\csname r@cor:dual68\endcsname{{32.9}{33}{}{proposition.32.9}{Dirichlet-GRH-v930-Volume-VI.pdf}}
\expandafter\def\csname r@cor:levelcurve370\endcsname{{37.29}{62}{}{proposition.37.29}{Dirichlet-GRH-v930-Volume-VII.pdf}}
\expandafter\def\csname r@cor:present116\endcsname{{32.20}{45}{}{proposition.32.20}{Dirichlet-GRH-v930-Volume-VI.pdf}}
\expandafter\def\csname r@cor:tangent121\endcsname{{33.12}{58}{}{proposition.33.12}{Dirichlet-GRH-v930-Volume-VI.pdf}}
\expandafter\def\csname r@eq:word68\endcsname{{32.1}{31}{}{equation.32.1}{Dirichlet-GRH-v930-Volume-VI.pdf}}
\expandafter\def\csname r@lem:algebra116\endcsname{{32.17}{43}{}{proposition.32.17}{Dirichlet-GRH-v930-Volume-VI.pdf}}
\expandafter\def\csname r@lem:bcpoints371\endcsname{{37.30}{63}{}{proposition.37.30}{Dirichlet-GRH-v930-Volume-VII.pdf}}
\expandafter\def\csname r@lem:charideal373\endcsname{{38.1}{71}{}{proposition.38.1}{Dirichlet-GRH-v930-Volume-VII.pdf}}
\expandafter\def\csname r@lem:coh\endcsname{{30.9}{8}{}{proposition.30.9}{Dirichlet-GRH-v930-Volume-VI.pdf}}
\expandafter\def\csname r@lem:detangent372\endcsname{{37.35}{68}{}{proposition.37.35}{Dirichlet-GRH-v930-Volume-VII.pdf}}
\expandafter\def\csname r@lem:diagonalpatch359\endcsname{{36.10}{9}{}{proposition.36.10}{Dirichlet-GRH-v930-Volume-VII.pdf}}
\expandafter\def\csname r@lem:fit13\endcsname{{11.2}{3}{}{proposition.11.2}{Dirichlet-GRH-v930-Volume-II.pdf}}
\expandafter\def\csname r@lem:flag69\endcsname{{32.10}{34}{}{proposition.32.10}{Dirichlet-GRH-v930-Volume-VI.pdf}}
\expandafter\def\csname r@lem:frameeuler365\endcsname{{37.7}{40}{}{proposition.37.7}{Dirichlet-GRH-v930-Volume-VII.pdf}}
\expandafter\def\csname r@lem:int\endcsname{{30.10}{9}{}{proposition.30.10}{Dirichlet-GRH-v930-Volume-VI.pdf}}
\expandafter\def\csname r@lem:int13\endcsname{{11.3}{4}{}{proposition.11.3}{Dirichlet-GRH-v930-Volume-II.pdf}}
\expandafter\def\csname r@lem:kummer72\endcsname{{32.14}{39}{}{proposition.32.14}{Dirichlet-GRH-v930-Volume-VI.pdf}}
\expandafter\def\csname r@lem:matrixmodules376\endcsname{{38.13}{79}{}{proposition.38.13}{Dirichlet-GRH-v930-Volume-VII.pdf}}
\expandafter\def\csname r@lem:scale118\endcsname{{34.5}{68}{}{proposition.34.5}{Dirichlet-GRH-v930-Volume-VI.pdf}}
\expandafter\def\csname r@lem:spandual367\endcsname{{37.14}{49}{}{proposition.37.14}{Dirichlet-GRH-v930-Volume-VII.pdf}}
\expandafter\def\csname r@lem:strong119\endcsname{{33.4}{51}{}{proposition.33.4}{Dirichlet-GRH-v930-Volume-VI.pdf}}
\expandafter\def\csname r@lem:support35\endcsname{{31.4}{16}{}{proposition.31.4}{Dirichlet-GRH-v930-Volume-VI.pdf}}
\expandafter\def\csname r@nav:III-00\endcsname{{30.4}{8}{}{subsection.30.4}{Dirichlet-GRH-v930-Volume-VI.pdf}}
\expandafter\def\csname r@prop:auxfield361\endcsname{{36.15}{15}{}{proposition.36.15}{Dirichlet-GRH-v930-Volume-VII.pdf}}
\expandafter\def\csname r@prop:classrank365\endcsname{{37.6}{39}{}{proposition.37.6}{Dirichlet-GRH-v930-Volume-VII.pdf}}
\expandafter\def\csname r@prop:diamond370\endcsname{{37.26}{59}{}{proposition.37.26}{Dirichlet-GRH-v930-Volume-VII.pdf}}
\expandafter\def\csname r@prop:field12\endcsname{{6.4}{27}{}{proposition.6.4}{Dirichlet-GRH-v930-Volume-I.pdf}}
\expandafter\def\csname r@prop:finite15\endcsname{{11.4}{6}{}{proposition.11.4}{Dirichlet-GRH-v930-Volume-II.pdf}}
\expandafter\def\csname r@prop:fixed116\endcsname{{32.18}{43}{}{proposition.32.18}{Dirichlet-GRH-v930-Volume-VI.pdf}}
\expandafter\def\csname r@prop:odddescent376\endcsname{{38.14}{79}{}{proposition.38.14}{Dirichlet-GRH-v930-Volume-VII.pdf}}
\expandafter\def\csname r@prop:orighecke374\endcsname{{38.7}{74}{}{proposition.38.7}{Dirichlet-GRH-v930-Volume-VII.pdf}}
\expandafter\def\csname r@prop:origideal373\endcsname{{38.2}{71}{}{proposition.38.2}{Dirichlet-GRH-v930-Volume-VII.pdf}}
\expandafter\def\csname r@prop:quadratic345\endcsname{{36.3}{3}{}{proposition.36.3}{Dirichlet-GRH-v930-Volume-VII.pdf}}
\expandafter\def\csname r@prop:res23\endcsname{{6.11}{32}{}{proposition.6.11}{Dirichlet-GRH-v930-Volume-I.pdf}}
\expandafter\def\csname r@prop:strictclasses365\endcsname{{37.5}{38}{}{proposition.37.5}{Dirichlet-GRH-v930-Volume-VII.pdf}}
\expandafter\def\csname r@prop:tracecrystal375\endcsname{{38.10}{76}{}{proposition.38.10}{Dirichlet-GRH-v930-Volume-VII.pdf}}
\expandafter\def\csname r@prop:traceorder360\endcsname{{36.12}{13}{}{proposition.36.12}{Dirichlet-GRH-v930-Volume-VII.pdf}}
\expandafter\def\csname r@prop:upfamily372\endcsname{{37.34}{67}{}{proposition.37.34}{Dirichlet-GRH-v930-Volume-VII.pdf}}
\expandafter\def\csname r@sec:basesupport371\endcsname{{37.8}{62}{}{subsection.37.8}{Dirichlet-GRH-v930-Volume-VII.pdf}}
\expandafter\def\csname r@sec:classicaldefect360\endcsname{{36.6}{12}{}{subsection.36.6}{Dirichlet-GRH-v930-Volume-VII.pdf}}
\expandafter\def\csname r@sec:coho12\endcsname{{6.4}{28}{}{subsection.6.4}{Dirichlet-GRH-v930-Volume-I.pdf}}
\expandafter\def\csname r@sec:commonframe365\endcsname{{37.2}{38}{}{subsection.37.2}{Dirichlet-GRH-v930-Volume-VII.pdf}}
\expandafter\def\csname r@sec:levelmaps370\endcsname{{37.7}{59}{}{subsection.37.7}{Dirichlet-GRH-v930-Volume-VII.pdf}}
\expandafter\def\csname r@sec:lower36\endcsname{{31.5}{20}{}{subsection.31.5}{Dirichlet-GRH-v930-Volume-VI.pdf}}
\expandafter\def\csname r@sec:missing\endcsname{{6.6}{32}{}{subsection.6.6}{Dirichlet-GRH-v930-Volume-I.pdf}}
\expandafter\def\csname r@sec:ordinaryfamily363\endcsname{{36.9}{25}{}{subsection.36.9}{Dirichlet-GRH-v930-Volume-VII.pdf}}
\expandafter\def\csname r@sec:original68\endcsname{{32.4}{31}{}{subsection.32.4}{Dirichlet-GRH-v930-Volume-VI.pdf}}
\expandafter\def\csname r@sec:pairedtests367\endcsname{{37.4}{48}{}{subsection.37.4}{Dirichlet-GRH-v930-Volume-VII.pdf}}
\expandafter\def\csname r@sec:seed361\endcsname{{36.7}{15}{}{subsection.36.7}{Dirichlet-GRH-v930-Volume-VII.pdf}}
\expandafter\def\csname r@sec:selmerlength368\endcsname{{37.5}{52}{}{subsection.37.5}{Dirichlet-GRH-v930-Volume-VII.pdf}}
\expandafter\def\csname r@sec:selmerlift369\endcsname{{37.6}{55}{}{subsection.37.6}{Dirichlet-GRH-v930-Volume-VII.pdf}}
\expandafter\def\csname r@sec:upper35\endcsname{{31.3}{15}{}{subsection.31.3}{Dirichlet-GRH-v930-Volume-VI.pdf}}
\expandafter\def\csname r@seriesguide\endcsname{{1.1}{4}{}{subsection.1.1}{Dirichlet-GRH-v930-Volume-I.pdf}}
\expandafter\def\csname r@thm:adjointseed361\endcsname{{36.17}{17}{}{proposition.36.17}{Dirichlet-GRH-v930-Volume-VII.pdf}}
\expandafter\def\csname r@thm:alllabels375\endcsname{{38.11}{77}{}{proposition.38.11}{Dirichlet-GRH-v930-Volume-VII.pdf}}
\expandafter\def\csname r@thm:auxdescent361\endcsname{{36.18}{18}{}{proposition.36.18}{Dirichlet-GRH-v930-Volume-VII.pdf}}
\expandafter\def\csname r@thm:classicaldefect360\endcsname{{36.14}{14}{}{proposition.36.14}{Dirichlet-GRH-v930-Volume-VII.pdf}}
\expandafter\def\csname r@thm:commonframe365\endcsname{{37.8}{41}{}{proposition.37.8}{Dirichlet-GRH-v930-Volume-VII.pdf}}
\expandafter\def\csname r@thm:conductor232\endcsname{{19.15}{12}{}{proposition.19.15}{Dirichlet-GRH-v930-Volume-IV.pdf}}
\expandafter\def\csname r@thm:cot116\endcsname{{32.19}{44}{}{proposition.32.19}{Dirichlet-GRH-v930-Volume-VI.pdf}}
\expandafter\def\csname r@thm:crystalline121\endcsname{{33.11}{57}{}{proposition.33.11}{Dirichlet-GRH-v930-Volume-VI.pdf}}
\expandafter\def\csname r@thm:cyclo119\endcsname{{33.5}{52}{}{proposition.33.5}{Dirichlet-GRH-v930-Volume-VI.pdf}}
\expandafter\def\csname r@thm:evenroot358\endcsname{{36.6}{5}{}{proposition.36.6}{Dirichlet-GRH-v930-Volume-VII.pdf}}
\expandafter\def\csname r@thm:finite117\endcsname{{33.2}{49}{}{proposition.33.2}{Dirichlet-GRH-v930-Volume-VI.pdf}}
\expandafter\def\csname r@thm:first72\endcsname{{32.15}{40}{}{proposition.32.15}{Dirichlet-GRH-v930-Volume-VI.pdf}}
\expandafter\def\csname r@thm:form117\endcsname{{33.3}{50}{}{proposition.33.3}{Dirichlet-GRH-v930-Volume-VI.pdf}}
\expandafter\def\csname r@thm:full36\endcsname{{31.8}{20}{}{proposition.31.8}{Dirichlet-GRH-v930-Volume-VI.pdf}}
\expandafter\def\csname r@thm:integer\endcsname{{30.11}{10}{}{proposition.30.11}{Dirichlet-GRH-v930-Volume-VI.pdf}}
\expandafter\def\csname r@thm:iwahori370\endcsname{{37.28}{61}{}{proposition.37.28}{Dirichlet-GRH-v930-Volume-VII.pdf}}
\expandafter\def\csname r@thm:levelstrict345\endcsname{{22.17}{75}{}{proposition.22.17}{Dirichlet-GRH-v930-Volume-IV.pdf}}
\expandafter\def\csname r@thm:markedcurve366\endcsname{{37.13}{46}{}{proposition.37.13}{Dirichlet-GRH-v930-Volume-VII.pdf}}
\expandafter\def\csname r@thm:mixedtrace359\endcsname{{36.11}{10}{}{proposition.36.11}{Dirichlet-GRH-v930-Volume-VII.pdf}}
\expandafter\def\csname r@thm:oddtrace123\endcsname{{34.9}{71}{}{proposition.34.9}{Dirichlet-GRH-v930-Volume-VI.pdf}}
\expandafter\def\csname r@thm:pres70\endcsname{{32.13}{37}{}{proposition.32.13}{Dirichlet-GRH-v930-Volume-VI.pdf}}
\expandafter\def\csname r@thm:primes367\endcsname{{37.17}{51}{}{proposition.37.17}{Dirichlet-GRH-v930-Volume-VII.pdf}}
\expandafter\def\csname r@thm:reduciblefiber376\endcsname{{38.15}{80}{}{proposition.38.15}{Dirichlet-GRH-v930-Volume-VII.pdf}}
\expandafter\def\csname r@thm:rescurve364\endcsname{{37.4}{36}{}{proposition.37.4}{Dirichlet-GRH-v930-Volume-VII.pdf}}
\expandafter\def\csname r@thm:selmerlift369\endcsname{{37.25}{58}{}{proposition.37.25}{Dirichlet-GRH-v930-Volume-VII.pdf}}
\expandafter\def\csname r@thm:smooth69\endcsname{{32.11}{35}{}{proposition.32.11}{Dirichlet-GRH-v930-Volume-VI.pdf}}
\expandafter\def\csname r@thm:stableprime371\endcsname{{37.32}{65}{}{proposition.37.32}{Dirichlet-GRH-v930-Volume-VII.pdf}}
\expandafter\def\csname r@thm:strong118\endcsname{{34.6}{69}{}{proposition.34.6}{Dirichlet-GRH-v930-Volume-VI.pdf}}
\expandafter\def\csname r@thm:trace71\endcsname{{34.2}{65}{}{proposition.34.2}{Dirichlet-GRH-v930-Volume-VI.pdf}}
\expandafter\def\csname r@thm:tracebound118\endcsname{{34.7}{70}{}{proposition.34.7}{Dirichlet-GRH-v930-Volume-VI.pdf}}
\expandafter\def\csname r@thm:traces\endcsname{{30.8}{8}{}{proposition.30.8}{Dirichlet-GRH-v930-Volume-VI.pdf}}
\expandafter\def\csname r@thm:upper35\endcsname{{31.5}{17}{}{proposition.31.5}{Dirichlet-GRH-v930-Volume-VI.pdf}}
\expandafter\def\csname r@thm:word68\endcsname{{32.8}{32}{}{proposition.32.8}{Dirichlet-GRH-v930-Volume-VI.pdf}}
\makeatother
\begin{document}
\setcounter{part}{7}
\fancyhead[L]{\small\color{muted}Dirichlet GRH \textperiodcentered\ Volume VIII}
\fancyhead[R]{\small\color{muted}Integral trace quotients}
\thispagestyle{plain}\begin{center}
{\LARGE\bfseries The generalized Riemann hypothesis\\[2pt]for Dirichlet \(L\)-functions}\\[6pt]
{\small 27 November 2026 \quad\textperiodcentered\quad Edition 930}\end{center}
\part{Integral trace quotients and Hecke support}\label{part:VIII}
\begin{abstract}\noindent We study the part of the trace ring detected by reducible representations. Integral trace quotients and Frobenius determinants give bounds for strict classes. We then descend to auxiliary fields and examine the support of the Hecke module on ordinary curves.\end{abstract}
\setcounter{section}{39}
\noindent The \hyperref[seriesguide]{guide in Volume~I} lists all twenty-seven volumes.
\begingroup\small\setlength{\parskip}{0pt}\tableofcontents\endgroup
\clearpage

\section{Integral trace quotients and root modules}\label{nav:chapter-integral-trace-quotients}\label{nav:volume-VIII}
\subsection{The integral kernel after restriction to the cuspidal component}
\label{sec:integraldefect377}

Fix a large-prime triple for the original period field \(E_x\).
Keep \(T,\mathfrak a_G,\mathfrak d,\Phi,C\) over the unramified
coefficient ring \(O\) of Section~\ref{sec:classicaldefect360}.
Let \(Q_a\) be the finite algebra of
Proposition~\ref{prop:orighecke374}, allowing the coefficient
extension there. Set
\[
 B_a=\operatorname{im}(T\longrightarrow O\times Q_a),
 \qquad M_a=\operatorname{im}(\mathfrak d\longrightarrow B_a).
\]
The first coordinate is the split evaluation.
Both \(B_a\) and \(M_a\) are finite \(O\)-modules. Write
\(M_{a,\rm tor}\) for the \(O\)-torsion in \(M_a\) and put
\(J=\operatorname{Frac}O\). We first determine the rational
local rings on which this ideal can be nonzero.

\begin{lemma}[The cyclotomic image at an irreducible trace point]
\label{lem:cyclointegral377}
For an absolutely irreducible geometric integral trace
specialization \(x\) of \(T\), the image of
\(G_{K(\mu_{\ell^\infty})}\) under \(\rho_x\) is enormous
in the sense of Newton and Thorne, after finite extension
of its coefficient field.
\end{lemma}
\begin{proof}
Let \(D\) be the Zariski closure of the image over the
algebraic closure and \(D^0\) its identity component.
The unipotent radical of \(D^0\) has a nonzero fixed vector
by triangularization. Its fixed space is stable under all of
\(D\) by normality, so this radical acts trivially. Thus
\(D^0\) is reductive. Its isotypic summands on the original
three-dimensional space are permuted transitively. If it
acts reducibly, primeness of the dimension leaves a scalar
action or three distinct lines. A scalar action on an open
subgroup contradicts the three distinct Hodge--Tate weights
of Proposition~\ref{prop:tracecrystal375}, which persist
under finite restriction.

In the three-line case the permutation quotient kills the
matrix image of \(\ker\bar\Psi\). This image is pro-\(\ell\)
since reduction in a triangular residual lattice basis is
unipotent on the kernel. The quotient is therefore a
transitive cyclic subgroup of \(S_3\), of order three.
In a basis of the three lines, conjugation by the diagonal
matrix with entries \(1,\zeta_3,\zeta_3^2\) gives a nontrivial
cubic self-twist. Such a character remains of order three
on reduction, because \(\ell>3\). It would permute the
residual characters \(1,\bar\Psi,\bar\Psi^{-1}\) transitively
by multiplication. They would form the group generated by
that cubic character, whereas \(\bar\Psi\) has order greater
than four on inertia. This proves irreducibility of \(D^0\).

The center of \(D^0\) acts by scalars and the determinant is
one, so \(D^0\) is semisimple. The closure of the cyclotomic
kernel contains it. Indeed that kernel contains all commutators,
the commutator is an algebraic morphism, and a connected
semisimple group is its own algebraic derived group.
Inertia above \(\ell\) provides an element in \(D^0\) with
three distinct eigenvalues, by choosing in its open inverse
image an inertia element of infinite cyclotomic order.
The locus of such elements in the connected group is
nonempty open. For any nonzero endomorphism \(Z\), the
function \(h\mapsto\operatorname{tr}(hZ)\) does not vanish
identically on \(D^0\), by irreducibility, Burnside's theorem
and nondegeneracy of the matrix trace. The cyclotomic image
meets the intersection of the two open sets. At that element
some eigenprojector pairs nontrivially with \(Z\), by
expanding the element in its eigenprojectors. Applying this
to a nonzero vector of every simple submodule gives the
enormous-image test. The eigenvalues can be split by one
finite coefficient extension, since a local field has only
finitely many extensions of bounded degree. This is the
criterion used in Theorem~\ref{thm:cyclo119}.~\cite{VIII-VI-VI-nt119}
\end{proof}

\begin{theorem}[Reducedness at the irreducible classical points]
\label{thm:rigidtrace377}
Let \(x\) be an absolutely irreducible characteristic-zero
point of the finite quotient \(Q_a\), viewed as a trace
point on \(T\). The geometric local ring of \(T[1/\ell]\)
at this point is its residue field, after a finite splitting
extension of coefficients. Consequently the same holds
for its local factor in \(B_a[1/\ell]\).
\end{theorem}
\begin{proof}
The argument of Theorem~\ref{thm:alllabels375} applies
directly when absolute irreducibility is given. Hence
\(\rho=\rho_x\), twisted by \(\kappa_\ell^{-1}\), is
the representation of a regular algebraic, conjugate
self-dual cuspidal representation on \(\mathrm{GL}_3/K\).
The lemma proves the enormous-image hypothesis for its
cyclotomic restriction, where the twist disappears. Let
\(V=\operatorname{End}(\rho)\) with its polarized adjoint
action of \(G_\QQ\). The action on \(G_K\) is conjugation.
At complex conjugation it is the negative adjoint for the
polarization pairing. Indeed \(\rho^c\simeq\rho^\vee\) by
the trace identity and absolute irreducibility. The
intertwining form is unique up to scalar, and applying
\(c\) twice shows it is symmetric or alternating. It is
symmetric in odd dimension. Its signed action therefore
agrees with the adjoint in Newton and Thorne's Theorem~A,
unchanged by the common Tate twist. That theorem gives
\(H^1_f(\QQ,V)=0\).~\cite{VIII-VI-VI-nt119}

Make the coefficient extension so that the point is rational
over a finite extension \(J'/J\), with integer
ring \(O'\). To test the cotangent space take an algebraic
tangent map over \(J'\) into \(J'[\epsilon]/(\epsilon^2)\).
It is continuous on \(T\otimes_O O'\). Present the latter
as a quotient of a power series ring and translate the
variables by their integral values in the maximal ideal
of \(O'\). Modulo the square of the ideal of the point
a series then maps by its value and its first derivatives.
These evaluations are continuous for the adic topology.
This is the argument of Lemma~\ref{lem:detangent372}.

The resulting determinant is represented by a continuous
matrix lift of \(\rho\) over the dual numbers, unique up
to strict conjugation. Write it as
\((1+\epsilon z(h))\rho(h)\) on \(G_K\). Its cohomology
class belongs to \(H^1(K,V)\). It is invariant for the
polarized action of \(c\). The conjugate lift and its
contragredient have the same trace, and after identifying
their reductions by the polarization they are strictly
conjugate by uniqueness for the irreducible determinant.
Differentiating gives precisely invariance of the class
under negative adjoint and conjugation of the group
argument. Inflation--restriction, since two is invertible,
identifies it with a class in \(H^1(\QQ,V)\).
~\cite{VIII-VI-chenevier364}

At the places of \(K\) away from \(\ell\) the lift has only
the prescribed finite inertia, so the extension class on
inertia is zero by averaging. At the places over \(\ell\)
Proposition~\ref{prop:tracecrystal375} makes the underlying
self-extension crystalline. It splits over \(B_{\rm cris}\),
giving the finite condition for its adjoint class. These
conditions descend to \(\QQ\) by restriction and
corestriction on inertia and on \(B_{\rm cris}\)-cohomology
respectively. The degree at each step is at most two and
is nonzero in \(J'\). Thus the class vanishes by the
Selmer theorem above. The lift on \(G_K\) is strictly
conjugate to the constant lift and gives constant traces.
Since these traces topologically generate \(T\), continuity
makes the entire tangent map constant.

The local ring at the rational point is Noetherian. The
vanishing of its tangent space gives zero maximal ideal
by Nakayama's lemma. It is therefore the residue field.
Passing to its quotient in \(B_a[1/\ell]\) preserves this
conclusion.
\end{proof}

\begin{proposition}[The rational image and the exact torsion term]
\label{prop:torsionterm377}
Let \(X_a'\) be the distinct geometric points of
\(B_a[1/\ell]\) at which some inverse-trace difference is
nonzero. Over a splitting field \(M_a[1/\ell]\) is the direct
sum of one copy of the field for each \(\xi\in X_a'\), and
every \(\xi\) belongs to the original list \(X\) of
Section~\ref{sec:classicaldefect360}.
For a test prime \(u\in\mathcal P\), \(u\ne\ell\), write
\(\gamma_u=T_u-d_\alpha(u)\) as in that section. With
\(v_\ell(\ell)=1\) one has
\[
 \len_O(M_a/\gamma_u M_a)
  =\len_O M_{a,\rm tor}[\gamma_u]
       +\sum_{\xi\in X_a'}v_\ell\bigl(\xi(T_u)-d_\alpha(u)\bigr).
\]
\end{proposition}
\begin{proof}
The inclusion \(B_a\subset O\times Q_a\) is integral since
the right side is finite over \(O\). Thus each geometric
generic point of \(B_a\) lifts to one on the right. At a
reducible trace point the projected ideal vanishes throughout
the rational local factor, by Theorem~\ref{thm:reduciblefiber376}.
At an irreducible point the preceding theorem gives a field.
There the projected ideal is either zero or the full field.
Theorem~\ref{thm:alllabels375} gives the labels claimed.
This describes the entire ideal over the splitting field.

On the free \(O\)-module \(M_a/M_{a,\rm tor}\), multiplication
by \(\gamma_u\) is injective. Indeed at each contributing
point the factor \(\xi(T_u)-d_\alpha(u)\) is nonzero by
Proposition~\ref{prop:traceorder360}. Its determinant is
the product of these factors, each occurring once. Smith
normal form gives the sum in the formula, since the original
\(O\) is unramified. Apply the kernel and cokernel sequence
of multiplication by \(\gamma_u\) to the torsion submodule,
the full module and its free quotient. The kernel on the
free quotient vanishes. The endomorphism of the finite
torsion submodule has cokernel of the same length as its
kernel. Adding the two cokernel lengths proves the identity.
\end{proof}

\begin{corollary}[The two integral contributions to the remaining length]
\label{cor:twointegral377}
Let \(C_a\) be as in Corollary~\ref{cor:componentdefect375}
and put \(C_a^0=\Phi(\ker(\mathfrak d\to M_a))\).
Then \(C_a^0\subseteq C_a\), and the image of
\(M_{a,\rm tor}\) in \(C/C_a^0\) is exactly \(C_a/C_a^0\).
For every one of the above \(u\),
\[
 \len_O(C_a/C_a^0)\le \len_O M_{a,\rm tor}[\gamma_u],
 \qquad
 \len_O(C/C_a^0)\le\len_O(M_a/\gamma_u M_a).
\]
\end{corollary}
\begin{proof}
The map \(M_a\to C/C_a^0\) is defined by choosing a lift
to \(\mathfrak d\) and applying \(\Phi\). It is surjective.
On \(\mathfrak d\) the kernel of evaluation at all geometric
points of the finite quotient is exactly the inverse image
of \(M_{a,\rm tor}\). Indeed elements of \(M_a[1/\ell]\)
are detected by those evaluations by the proposition, and
the additional split coordinate of \(B_a\) is zero on
\(\mathfrak d\). This is precisely the intersection with
\(P_a\) in Corollary~\ref{cor:componentdefect375}, proving
the image formula. Multiplication by \(\gamma_u\) has zero
image under \(M_a\to C/C_a^0\), since
\(\Phi(\mathfrak a_G\mathfrak d)=0\). Both inequalities
now follow by taking quotients, and by equality of kernel
and cokernel lengths on finite modules.
\end{proof}

In particular the classical contribution on the common
interval has the polynomial bound of
Corollary~\ref{cor:componentdefect375}. The two remaining
integral quantities are the image of the kernel of
\(\mathfrak d\to M_a\) and the part of the torsion of \(M_a\)
that survives in \(C/C_a^0\). Bounds on the characteristic-zero
points above account for all their rational multiplicities.
The weighted bounds for these two integral quantities and
the projected theta-period comparison in
Section~\ref{sec:missing} are still required.


\subsection{A finite trace ring supported on reducible representations}
\label{sec:reducibleminors378}

Fix one of the large-prime triples on the common interval of
Section~\ref{sec:seed361}. Let \(G=G_K\) and keep the original
coefficient ring \(O\), with residue field \(k\). Write \(E=E_M\)
for the finite \(T\)-algebra of universal group matrices and
\(t\) for their trace. In this subsection \(E\) denotes this
algebra and \(E_x\) denotes the period field. Put \(I_m=I_m(G)\).
We construct a finite quotient directly from the pairing
\((a,b)\mapsto t(ab)\). Its defining ideal can be given
explicitly, even where the trace ring has torsion.

\begin{lemma}[All minors of the trace pairing]
\label{lem:allminors378}
Use the three rank-one projectors of
Lemma~\ref{lem:charideal373}. Put \(E_{ij}=e_iEe_j\), and let
\(B_{ij}\) be the ideal of \(T\) spanned by \(t(ab)\) for
\(a\in E_{ij}\), \(b\in E_{ji}\). Thus \(B_{ii}=T\) and
\(B_{ij}=B_{ji}\). For \(1\le j\le6\), \(I_{3+j}\) is the
sum of the products of \(j\) ideals chosen from the six
\emph{positions} in the list
\[
                 B_{12},B_{12}, B_{13},B_{13}, B_{23},B_{23}.
\]
Each position is used at most once. In particular,
\[
 I_5=\mathfrak a_G^2,\qquad I_9=(B_{12}B_{13}B_{23})^2,
 \qquad I_{10}=0.
\]
\end{lemma}
\begin{proof}
The minors have the same ideal if rows and columns run
through \(E\), since the group matrices span it over \(T\).
Split each argument into its nine corners. The trace pairs
a corner \((i,j)\) only with the opposite corner \((j,i)\),
by orthogonality of the idempotents and cyclicity.
If the multiplicities of a row corner and its opposite
column corner differ, no permutation in the determinant
can contribute. For matching multiplicities the determinant
factors up to sign into the corresponding block determinants.
In a basis of the three rank-one summands over the local
marked coefficient ring \(A\), each block has entries of
the form \(a_u b_v\) in \(A\). Every block of size at least
two has determinant zero. The inclusion \(T\subseteq A\)
makes the same assertion true in \(T\). Thus only distinct
row corners contribute. Conversely, choosing one row and
one opposite column in each of any set of distinct corners
gives their product of pairings up to sign. Arbitrary generators
of each ideal product occur this way. There are three diagonal
positions, with pairing one on \((e_i,e_i)\). Every product
of \(3+j\) corners therefore belongs to one of the stated
products of \(j\) off-diagonal ideals, and adjoining the
three diagonal positions realizes each such product.
For \(j=1\) the sum is \(I_4=\mathfrak a_G\) by
Proposition~\ref{prop:origideal373}. For \(j=2\) every square
and every pairwise product in its expansion is allowed.
For \(j=6\) all positions are forced. Ten distinct positions
are impossible.
\end{proof}

\begin{lemma}[Reducible traces over a field of Laurent series]
\label{lem:curvedescent378}
Let \(k'/k\) be finite and \(T\to k'[[z]]\) a continuous
local \(O\)-algebra map. If its geometric semisimple
representation of \(G\) is reducible over
\(\overline{k'((z))}\), every trace belongs to the constant field.
\end{lemma}
\begin{proof}
Put \(F=k'((z))\). The determinant over a field has a
semisimple matrix realization over its algebraic closure
by Chenevier's theorem on determinants over algebraically
closed fields. This realization factors through
\(E\otimes_T\overline F\). In fact its trace there agrees
on group elements, and for any relation among the group
matrices in the latter algebra the trace of the relation
times every group element vanishes. On the semisimple
image the pairing is nondegenerate: it is the ordinary
matrix trace on each simple factor multiplied by its
multiplicity, an integer from one to three, invertible in
characteristic \(\ell>3\). This proves the factorization.
A finite basis of the finite-dimensional algebra and the
entries of its image matrices show that a finite field
extension of \(F\) suffices. The realization is continuous
there, since \(E\) is finite over \(T\) with its module
topology. The compact matrix image preserves a lattice:
the span of all translates of any fixed lattice is bounded,
finitely generated over the valuation ring, and stable.
Reduction of such a lattice has semisimplification
\(\bar\Psi\oplus1\oplus\bar\Psi^{-1}\), by the residual
determinant. The three characters are distinct on inertia.
~\cite{VIII-VI-chenevier364}

If there are three lines, \(I_4\) vanishes at the point and
Proposition~\ref{prop:origideal373} makes the trace constant.
Otherwise write the representation as \(\theta\oplus P\),
with \(P\) absolutely irreducible of dimension two.
The relation between conjugation by \(c\) and duality for
the trace gives \(\theta^c=\theta^{-1}\) by uniqueness of
the line. Its residual character is fixed by \(c\), since
all three residual characters extend to \(G_\QQ\). Only
the middle character has square one. Thus \(\bar\theta=1\).

At any place over \(\ell\) the inertia matrices are upper
triangular in the order \(\bar\Psi,1,\bar\Psi^{-1}\).
Here these are the original fixed characters specialized
in equal characteristic. To see the assertion on this
realization, the ordered product in \eqref{eq:word68}
vanishes already in \(E\). Choose one separating inertia
element \(g_u\). In its eigenbasis replace just the \(j\)-th
letter by an arbitrary inertia element \(h\), leaving all
other letters equal to \(g_u\). Left multiplication by the
preceding factors is invertible on rows of index at least
\(j\). Right multiplication by the following factors is
invertible on columns of index at most \(j\). It follows
that in this rectangle the matrix of \(h\) agrees with
the \(j\)-th prescribed character times the identity.
All three rectangles give the assertion. In particular
\(\theta\) is trivial on inertia by its residual character.
Elsewhere inertia factors through a group of order prime
to \(\ell\) and \(\theta\) takes values in principal units,
so it is trivial there too. The class number of \(K\) is
prime to \(\ell\) on the interval. Class field theory gives
\(\theta=1\), and then \(\det P=1\).

We spell out descent of the plane in this characteristic.
The marked involution \(a\mapsto J_3a^{\mathsf T}J_3\)
preserves the trace and descends through the radical of
the specialized trace pairing to the semisimple image
\(\overline F\times M_2(\overline F)\). It fixes both
factors. On the plane it is adjunction by a symmetric
form. Indeed the identity \(2XYX=t(XY)X\) for skew elements
in Proposition~\ref{prop:odddescent376} is an integral
polynomial identity. Skew elements lift through the quotient
by taking half the difference with their adjoints. For
alternating adjunction the traceless matrices
\(X=\operatorname{diag}(1,-1)\), \(Y=E_{12}+E_{21}\)
in the plane violate the identity since \(XYX=-Y\).
The symmetric form with matrix \(S\), together with the
alternating determinant form \(J_0\), now gives the
intertwiner \(S^{-1}J_0\) for conjugation by \(c\).
It has trace zero and nonzero scalar square. Rescale it
to square one and assign it to \(c\). This yields a
continuous representation \(W\) of \(G_\QQ\) with
\(W|_G=P\) and \(\det W=\bar\eta\), after finite
coefficient extension.

Choose a stable lattice for \(W\). Over a finite residue
extension its residual semisimplification is a sum of
two characters. Indeed an irreducible residual plane would
restrict semisimply to the normal subgroup of index two:
the translates of a simple submodule span the plane.
The two distinct residual character lines of the restriction
would then each be preserved by \(c\), since their characters
extend to \(G_\QQ\). This contradicts irreducibility.
Twisting \(W\) by \(\bar\eta\) if necessary, the constituent
extending \(\bar\Psi|_G\) is \(\bar\Psi\) itself. The
determinant makes the other \(\bar\eta\bar\Psi^{-1}\).

Over \(K_u\) the first inertia eigenline just constructed
belongs to the plane: its eigenvalue at \(g_u\) differs from
that of the trivial line. It carries \(\bar\Psi\) on inertia.
Since \(K_u/\QQ_\ell\) is unramified, the inertia groups
agree. A conjugate of this line by \(G_{\QQ_\ell}\) again
carries \(\bar\Psi\) on inertia, since \(\bar\Psi\) is a global
character. Separation at \(g_u\) gives uniqueness of the line.
Thus \(W\bar\Psi^{-1}\) has reducible local trace and the
constituent lifting the trivial residual character is
unramified at \(\ell\). Its determinant is the specialization
of \(\chi=\eta\Psi^{-2}\), an odd characteristic-zero character
de Rham at \(\ell\), with nontrivial reduction there because
\(0<2r<\ell-1\). The representation is unramified outside
the places dividing \(sN_*\ell\), where \(K=\QQ(\sqrt{-s})\).

Apply Pan's Theorem~5.1.2(1) over \(\QQ\) with this finite
ramification set and determinant \(\chi\). Its hypotheses
are listed in Section~5.1.1 of that paper. The ordinary
pseudodeformation ring is finite over the completed group
algebra of \(1+\ell\ZZ_\ell\). The preceding unramified
character gives the augmentation of this group algebra.
Using the ring of integers of our Laurent series field
extension as coefficient target, the map therefore factors
through a finite \(O\)-module, as in Theorem~\ref{thm:finite117}.
One may enlarge \(O\) unramified to match the residue field.
In characteristic \(\ell\) the image is a finite integral
domain, hence a finite field. All traces of \(P\) lie in
the constant field after restoring the finite-valued twist.
Intersecting that field with \(k'[[z]]\) gives \(k'\).
~\cite{VIII-VI-VI-pan40}
\end{proof}

\begin{theorem}[Finiteness on the full reducible locus]
\label{thm:reducibleminors378}
For \(6\le m\le9\), the original trace algebra \(T/I_m\) is
finite over \(O\). At any of its geometric field points
the semisimple representation of \(G\) has a one-dimensional
summand. In particular the conclusion holds for the single
product ideal \((B_{12}B_{13}B_{23})^2\).
\end{theorem}
\begin{proof}
Over an algebraically closed field of the given residue
characteristic or characteristic zero, the group span on
a semisimple representation is the product of the full
matrix algebras of its distinct simple constituents.
This follows from the semisimple quotient and the matrix
algebra theorem over an algebraically closed field.
Its trace pairing is nondegenerate for the multiplicity
reason in the lemma. At an irreducible point it therefore
has rank nine. Some nine group matrices give a basis,
and pairings with group matrices span the full dual.
There is a nonzero minor of each size through nine:
choose that many independent rows and then independent
restrictions of column functionals to their span.
At a reducible point the sum of the squares of the
constituent dimensions is at most five, so every minor
of size at least six vanishes. This proves the point
assertion in both characteristics, and in particular
on the generic point of any complete curve in the
special fiber.

The complete Noetherian local ring \(T/I_m\) has finite
residue field. If its special fiber had positive dimension,
it would have a one-dimensional complete local domain
quotient. Normalize this quotient. Normalization of a
one-dimensional complete local domain is finite, and a
local factor is an equal-characteristic complete DVR with
finite residue field, of the form \(k'[[z]]\). This gives
a continuous local map with nonconstant image. Since
traces and coefficients topologically generate \(T\),
some trace would be nonconstant, contrary to the lemma.
Hence the special fiber is Artinian. Lift a finite
\(k\)-basis to the ring. Successive approximation modulo
powers of the uniformizer expresses each element as an
\(O\)-linear combination of the lifts. The errors tend
to zero in the maximal ideal topology, and the coefficients
converge in \(O\). This proves finiteness. The product
formula is the lemma on minors.
\end{proof}

\begin{corollary}[Integral torsion surjecting onto the strict group]
\label{cor:torsionimage378}
Let \(B=T/I_9\) and \(N\) be the image in \(B\) of the ideal
\(\mathfrak d\) of inverse-trace differences. Then \(N\) is
a finite-length \(O\)-module. There is an \(O\)-linear
surjection \(N\to S(\Psi^{-1})^\vee\) for every large-prime
triple on the original interval. The same assertions hold
with \(I_9\) replaced by \(I_m\), \(5\le m\le8\).
\end{corollary}
\begin{proof}
At \(m=5\) finiteness follows from
\(T/I_5=O\oplus\mathfrak a_G/\mathfrak a_G^2\) and
Noetherianity. All generic field points there are split.
For the other sizes use the finiteness theorem. Split the
finite generic algebra over a finite extension of
\(\operatorname{Frac}O\) into local Artin factors with
rational residue points. These points are integral and
continuous on \(T\), since \(B\) is finite and local over
the complete DVR and all its special points coincide.
Their representations are reducible, so each factor kills
\(\mathfrak d\) by Theorem~\ref{thm:reduciblefiber376}.
Faithful field extension gives \(N[1/\ell]=0\). Thus \(N\)
is a torsion submodule of a finite module over a DVR and
has finite length. Finally \(I_m\subseteq\mathfrak a_G^2\)
for every indicated size. The map \(\Phi\) to the square-zero
extension in Theorem~\ref{thm:mixedtrace359} factors through
this quotient and its restriction to \(\mathfrak d\) has
image the full dual strict group. This proves the surjection.
\end{proof}

Each ideal and each module above is taken at the character
of the given \(E_x\). Finiteness gives an exponent annihilating
\(N\) for each triple. The weighted sum required in
Section~\ref{sec:missing} still requires control of these
orders as \(x,\ell,r,\alpha\) vary.


\subsection{The order of reducible traces on the common interval}
\label{sec:reducibleorder379}

Continue with \(B=T/I_9\) on \(G_K\). Write \(J=\operatorname{Frac}O\),
and let \(A=B/B[\ell^\infty]\), where \(B[\ell^\infty]\) denotes
the ideal of elements annihilated by some power of \(\ell\).
This is a finite free \(O\)-module. Its augmentation ideal
is denoted by \(\mathfrak n_A\). We give a uniform estimate for
\(\mathfrak n_A/\mathfrak n_A^2\) on the original intervals.

\begin{lemma}[The matrix algebra at a line and a plane]
\label{lem:planeunits379}
Let \(D\) be a finite local algebra over a finite extension
\(J'/J\), with residue field \(J'\), and let \(T\to D\) be as
in Lemma~\ref{lem:matrixmodules376}. Suppose the residue
representation is \(1\oplus P\) with absolutely irreducible
plane \(P\), realized over \(J'\). Then the group acts in
\[
              E_M\otimes_T D\simeq D\times M_2(D)
\]
as \(1\oplus P_D\), with \(P_D\) of determinant one.
Its trace is \(1+\operatorname{tr}P_D\), and the underlying
\(J'\)-representation on \(D^2\) is crystalline above \(\ell\).
\end{lemma}
\begin{proof}
Put \(R=E_M\otimes_T D\) and use the inertia projectors
\(e_1,e_2,e_3\) of Lemma~\ref{lem:charideal373}. The proof of
Theorem~\ref{thm:reduciblefiber376} gives a central \(e_2\)
with \(e_2R=De_2\), and the group acts as one on it.
For every \(i\) we have \(e_i R e_i=De_i\). Indeed sandwiching
any marked matrix between the two copies of the rank-one
projector gives its scalar trace on that line times \(e_i\).
The identity passes to \(R\) by linear extension and
\(t(e_i)=1\) identifies the scalars faithfully.

The semisimple quotient on the plane is \(M_2(J')\), with
the two remaining projectors as its diagonal units. Lift
opposite matrix units to \(x\in e_1Re_3\) and
\(y\in e_3Re_1\) by projecting arbitrary lifts. We have
\(xy=u e_1\), \(yx=u e_3\), since the two traces agree by
cyclicity. The scalar \(u\) reduces to one, hence is a unit.
Replace \(y\) by \(u^{-1}y\). Then \(x,y,e_1,e_3\) are
matrix units. Every \(z\in e_1Re_3\) satisfies
\(z=(zy)x\) with \(zy\in De_1\), and the opposite corner
is handled by interchanging \(x,y\). This proves the
matrix-algebra assertion, and the trace on it is the
ordinary trace because it is one on each diagonal unit and
zero on both off-diagonal units. Newton's identities applied
to the traces of the first three powers of each group
element give its original characteristic polynomial as
\((X-1)\det(X-P_D(h))\). Its determinant was one in the
original trace ring, so \(\det P_D=1\). The column module
\(D^2\) is a finite \(R\)-module. Its local properties now
follow from Lemma~\ref{lem:matrixmodules376}.
\end{proof}

\begin{theorem}[Rigidity at every reducible trace point]
\label{thm:reduciblerigid379}
At every geometric integral characteristic-zero trace point
of \(T\) with reducible representation, the geometric local
ring of \(T[1/\ell]\) is its residue field, after extension
to a finite field of definition. In particular \(B[1/\ell]\)
is finite \'etale over \(J\).
\end{theorem}
\begin{proof}
Choose a finite extension \(J'\) where the point and its
constituents are rational and split. Any relative tangent
map at this point can be tested in
\(D=J'[\epsilon]/(\epsilon^2)\). Its restriction to the
original complete local algebra is continuous. Indeed after
presenting the latter by power series, translation by the
tuple of integral point values in the maximal ideal of the
integer ring gives the value and first derivatives formula
of Lemma~\ref{lem:detangent372}. These are continuous
in the power series topology. For a split triple
Theorem~\ref{thm:reduciblefiber376} already makes this
tangent map constant.

In the other case take the preceding lift \(P_D\).
Choose the odd representation \(W\) extending \(P\) from
Proposition~\ref{prop:odddescent376} and choose compatible
bases in its residue representation. Write
\(P_D(h)=(1+\epsilon z(h))P(h)\). It defines a continuous
class in \(H^1(K,\operatorname{ad}^0W)\) by the cocycle
identity and the determinant. This class is invariant under
\(\operatorname{Gal}(K/\QQ)\) with its usual action through
\(W\). Indeed the marked trace satisfies
\(t(chc^{-1})=t(h^{-1})\), and \(P_D\) has determinant one.
For a two-by-two matrix of determinant one the traces of
the matrix and its inverse coincide. Subtracting the
constant line gives equal traces for \(P_D\) and the lift
\(h\mapsto W(c)^{-1}P_D(chc^{-1})W(c)\).
Their reductions agree and are absolutely irreducible.
The uniqueness of matrix realizations of an absolutely
irreducible residual determinant over a local Artin ring
gives strict conjugacy, hence equality of the two cocycle
classes.~\cite{VIII-VI-chenevier364}

Restriction from \(\QQ\) therefore gives this class by
inflation--restriction for the group of order two in
characteristic zero. The finite local conditions hold on
\(K\). Away from \(\ell\) the extension splits over inertia
since the image factors through the specified finite quotient.
At \(\ell\) the lemma makes the self-extension crystalline,
hence it splits after tensoring with \(B_{\rm cris}\). These
conditions descend to \(\QQ\): the restricted classes in
inertial cohomology away from \(\ell\), and in
\(B_{\rm cris}\)-cohomology at \(\ell\), vanish. Local
corestriction composes with restriction to multiplication
by the relevant index, at most two. It is invertible on
these cohomology groups.

By Proposition~\ref{prop:odddescent376}, a scalar twist of
\(W\) comes from a non-CM cuspidal newform of weight \(2r+1\).
The twist leaves its adjoint unchanged. Newton and Thorne's
Theorem~B(1) gives \(H_f^1(\QQ,\operatorname{ad}W)=0\).
Its trace-zero summand has zero finite Selmer group as well.
Consequently \(z\) is a coboundary. Strict conjugacy to
the constant lift makes all group traces constant. Their
topological generation of \(T\), together with continuity,
makes the original tangent map constant. The maximal ideal
of the Noetherian local ring has zero cotangent space and
is zero by Nakayama's lemma. This proves the local assertion.
Theorem~\ref{thm:reducibleminors378} gives finiteness of
\(B\) and reducibility at all its generic geometric points.
Each generic local factor is a quotient of the corresponding
local ring just determined. It is a field, which proves
the assertion about \'etaleness in characteristic zero.
~\cite{VIII-VI-VI-nt119}
\end{proof}

\begin{proposition}[The primitive forms on the original field]
\label{prop:planelabels379}
Write \(K=\QQ(\sqrt{-s})\), with discriminant \(-s\), and fix
a good rational prime \(u\) split in \(K\). Each nonsplit
geometric point of \(B[1/\ell]\) can be assigned a normalized
primitive cuspidal form \(f\) of weight \(2r+1\), nebentypus
\(\eta\alpha^{-2}\), and level the conductor of that
nebentypus. The level divides \(sN_*\). At either place over
\(u\nmid sN_*\ell\) its trace is
\[
                       1+u^{-r}\alpha(u)a_u(f).
\]
With a fixed identification of algebraic coefficient embeddings,
distinct trace points have distinct assigned forms.
\end{proposition}
\begin{proof}
Use \(W\) from Proposition~\ref{prop:odddescent376} at
\(1\oplus P\). Twisting its modular representation by the
finite character \(\alpha^{-1}\) makes
\(W\kappa_\ell^r\alpha^{-1}\) the arithmetic representation
of \(f\), with the stated determinant character after
removing \(\kappa_\ell^{2r}\). At \(p\ne s,\ell\) the
inertia groups of \(K\) and \(\QQ\) agree. On the plane
the two characters there are \(\alpha,\alpha^{-1}\):
the marked finite inertia has these two characters together
with the trivial line, with multiplicities read from its
traces. The twisted characters are \(1,\alpha^{-2}\).
At \(s\) the character \(\alpha\) is unramified. The plane
kills the inertia of \(K\), so \(W\) on rational inertia
factors through a group of order two. Its determinant
is the ramified quadratic character \(\eta\), so both
characters \(1,\eta\) occur. At all these finite places
monodromy vanishes, and the Artin conductor exponent is
that of \(\eta\alpha^{-2}\). Crystallinity at \(\ell\)
gives an unramified Weil--Deligne parameter with zero
monodromy there. Local--global compatibility and the
newform conductor formula give the primitive level.
Restriction to \(G_K\) removes the choice between the
two possible extensions of \(P\). Removing the scalar twist
there gives the displayed Frobenius formula. Conversely
the embedded form and the chosen twist determine \(P\)
on \(G_K\), hence determine the trace point since group
traces generate \(T\).~\cite{VIII-VI-VI-carayol117,VIII-VI-VI-saito40}
\end{proof}

\begin{theorem}[An integer bounding the reducible order on every field]
\label{thm:reducibleorder379}
Choose once the split prime \(u=u_1\) of
Proposition~\ref{prop:quadratic345}. For all sufficiently
distant common intervals there is a positive integer
\(D_2(V)\) such that for every \(x\in[V,V^{3/2}]\)
\[
 \begin{split}
 \sum_{\ell>P(x)}(\log\ell)\sum_{1\le r<t(x)}
       \sum_{[\alpha]} b_\alpha\,
           \operatorname{length}_O(\mathfrak n_A/\mathfrak n_A^2)
   &\le \log D_2(V),\\
 \log D_2(V)&\ll_u t_*^3 N_* (sN_*)^2.
 \end{split}
\]
Here \([\alpha]\) and \(b_\alpha\) are the residue orbits
of complementary sign on \(E_x\) and their degrees as in
Theorem~\ref{thm:classicaldefect360}. For each large prime
\(\ell\nmid D_2(V)\) one has \(A=O\) for every such triple.
Writing \(U=B[\ell^\infty]\) and \(\mathfrak n_B=\ker(B\to O)\),
one has an exact sequence
\[
 0\longrightarrow U/(U\cap\mathfrak n_B^2)
 \longrightarrow \mathfrak a_G/\mathfrak a_G^2
 \longrightarrow \mathfrak n_A/\mathfrak n_A^2
 \longrightarrow0 .
\]
The left term maps onto \(S(\Psi^{-1})^\vee\) under \(\Phi\).
\end{theorem}
\begin{proof}
Use all complex characters \(\alpha\) of \(E_*\) of sign
\((-1)^r\), all \(1\le r<t_*\), and for each such pair all
normalized primitive forms of the weight, character and level
in the proposition, including their coefficient embeddings.
Let \(D_2(V)\) be the absolute value of the product of
\[
     \alpha(u)a_u(f)-u^{2r}\alpha(u)-\alpha(u)^{-1}
\]
over this list, with empty product one. These algebraic
integers are permuted with their indices by absolute
Galois conjugation. The product is rational integral.
It is nonzero: at any complex embedding \(|a_u(f)|\le2u^r\)
by Deligne, whereas \(u^{2r}-1>2u^r\) for \(u\ge7\).
Each factor has absolute value at most \((u^r+1)^2\).
For fixed \(r,\alpha\) the number of normalized forms with
their embeddings is at most the dimension of
\(S_{2r+1}(\Gamma_1(sN_*))\), since distinct eigensystems
give linearly independent old or new vectors at this level.
The modular curve dimension bound is \(O((2r+1)(sN_*)^2)\)
by the usual index bound for \(\Gamma_1\), as in the proof
of Theorem~\ref{thm:tracebound118}. There are at most
\(N_*\) characters. Taking logarithms and summing over
\(r\) proves the upper estimate.~\cite{VIII-VI-VI-deligne40}

Fix a triple and set \(a_u=T_u-d_\alpha(u)\in\mathfrak n_A\).
The finite free module \(\mathfrak n_A\) has generic fiber
with one geometric line per nonsplit point, since the
preceding rigidity theorem makes every local factor a field.
Multiplication by \(a_u\) is injective on it by the nonvanishing
of the factors just proved. Its determinant is the product
over these points of their trace differences. The split
value is \(1+u^r\alpha(u)+u^{-r}\alpha(u)^{-1}\).
The displayed integer factors equal the trace differences
times \(u^r\), an \(\ell\)-adic unit throughout this range.
Smith normal form over the unramified \(O\) therefore gives
\[
 \operatorname{length}_O(\mathfrak n_A/\mathfrak n_A^2)
 \le \operatorname{length}_O(\mathfrak n_A/a_u\mathfrak n_A)
 =\sum_{\xi\ne\mathrm{split}}
      v_\ell\bigl(\xi(T_u)-d_\alpha(u)\bigr).
\]
The assigned forms are distinct for this character. For
each of its \(b_\alpha\) residue conjugates one may apply
the corresponding local embedding to the same factors:
valuations agree and the character labels are distinct.
Every character of \(E_x\) inflates to \(E_*\). Thus after
summing these bounds over twists and residue orbits for
fixed \(\ell\), the sum is at most \(v_\ell(D_2(V))\).
All omitted factors have nonnegative valuation. Summation
over primes with weight \(\log\ell\) proves the first bound.
If \(\ell\nmid D_2(V)\), the cotangent quotient vanishes.
The augmentation ideal then equals its square and lies in
the maximal ideal of this finite local ring. Nakayama's
lemma makes it zero, so \(A=O\).

Every torsion element of \(B\) has split value zero.
Taking augmentation ideals in \(B\to A\) and dividing
by their squares gives the displayed kernel and cokernel.
Here \(\mathfrak n_B/\mathfrak n_B^2=\mathfrak a_G/
\mathfrak a_G^2\) since \(I_9\subseteq\mathfrak a_G^2\).
The image of \(\mathfrak d\) in \(B\) lies in \(U\) by
Corollary~\ref{cor:torsionimage378}. Its image under
\(\Phi\) is the full dual strict group, and \(\Phi\) kills
the square of the augmentation ideal. This proves the
last assertion.
\end{proof}

The bound holds simultaneously at all the period fields in
the interval. It controls the quotient on the right of the
exact sequence. A uniform upper bound for the image of the
left term is still required for the large primes of
Section~\ref{sec:missing}.


\subsection{Two integral quotients of the three-character trace ring}
\label{sec:planedesc381}

Fix \(x\in[V,V^{3/2}]\) on one of the intervals of
Section~\ref{sec:missing}, and use a triple
\((\ell,r,\alpha)\) from Section~\ref{sec:seed361}.
In particular \(\ell>P(x)\), \(0<r<t(x)\),
\(\Psi=\alpha\kappa_\ell^r\) is even, and
\(K=\QQ(\sqrt{-s})\) has quadratic character \(\eta\).
Write \(G=G_K\), let \(c\) be the marked complex conjugation
and put \(D=\operatorname{Frac}(O)/O\).
Let \(A\) be the original fixed-weight marked matrix ring and
\(T\subset A\) its closed trace algebra on \(G\), with
augmentation ideal \(\mathfrak a\). We use one-based matrix
indices. At the split augmentation the matrices are
\(d(h)=\operatorname{diag}(\Psi(h),1,\Psi(h)^{-1})\).
Let \(J_3\) reverse the three coordinates. For \(h\in G\)
the polarization gives
\[
 M(ch^{-1}c)=J_3M(h)^{\mathsf T}J_3.
\]
The marking at \(c\) is the identity. The other marks set the
entries in positions \(12,21\) of
\(Z=M(g_*)-M(cg_*^{-1}c)\) to zero, for a separating inertia
element \(g_*\). Since \(Z=-J_3 Z^{\mathsf T}J_3\), it is
\(\operatorname{diag}(z,0,-z)\) with \(z\) a unit.
Indeed its first residue is \(\bar\Psi(g_*)-\bar\Psi(g_*)^{-1}\ne0\).
The identity \(z^2=\operatorname{tr}(Z^2)/2\) puts \(z^2\) in
\(T\). The square root with the specified residue lies in \(T\)
by Hensel's lemma. Hence the three standard diagonal projectors
\(e_i\) belong to the \(T\)-span \(\mathcal E\) of \(M(G)\),
by interpolation on \(z,0,-z\).

In \(A\) let \(J_1\) be generated by the entries \(M(h)_{ij}\)
with \(|i-j|=1\), for all \(h\in G\), and \(J_2\) by those with
\(|i-j|=2\). In \(T\) let \(X\) be generated by the traces of
\(e_1\mathcal E e_2\mathcal E e_1\), and \(Y\) by those of
\(e_1\mathcal E e_3\mathcal E e_1\). The reflection identity
identifies the ideal for the pair \(2,3\) with \(X\).
Use weight \(j-i\) for the \(ij\)-entry under conjugation by
\(\operatorname{diag}(a,1,a^{-1})\), with \(a\) an indeterminate.
Thus weight zero means invariance for the split torus on each
finite truncation.

\begin{lemma}[Contraction of the two matrix-entry ideals]
\label{lem:pairedcorners381}
One has
\[
 T\cap J_1=X,\qquad T\cap J_2=Y,\qquad
 X+Y=\mathfrak a,\qquad X^2\subseteq Y .
\]
In particular \(T/X\) and \(T/Y\) are the full invariant
subrings of \(A/J_1\) and \(A/J_2\), respectively.
\end{lemma}
\begin{proof}
The entries with their split constants subtracted generate \(A\)
topologically. Indeed the marked ring is generated by the
entries of the universal group matrices, as in the presentation
of Theorem~\ref{thm:trace71}. The ramification conditions give
a finitely generated pro-\(\ell\) quotient of the kernel of the
residual representation, by class field theory and the pro-\(\ell\)
generator criterion. Finitely many matrix entries therefore
suffice after lifting the finite residual image. The entries
on the other coset are obtained using the marked \(c\).
The defining conditions are preserved by the displayed torus.
On an Artinian truncation the action separates the integer
weights by the Laurent monomials in \(a\). Taking the degree-zero
coefficient of a polynomial lift thus lifts an invariant,
also after imposing either homogeneous entry ideal.

Every monomial of weight zero factors into diagonal entries
or their differences from split values, and blocks of weights
\((1,-1),(2,-2),(1,1,-2),(-1,-1,2)\). To see this, remove
all opposite pairs. In what remains the positive and negative
weights must have distinct absolute values and group into
the indicated triples. By reflection an entry of weight one
can be put in either position \(12,23\), and one of weight
minus one in either \(21,32\). Each block is therefore a
directed cycle. Its value is the trace of the corresponding
matrix product with the diagonal projectors inserted, so it
lies in \(T\). This proves the assertion about invariants
before imposing an ideal, since \(T\) is closed. After imposing
an ideal, lifting invariants on each truncation and using the
compactness of \(T\) shows that all invariants still come from \(T\).

In the image of the ideal \(J_j\) at a truncation, take a
polynomial expression with one specified generating entry of
that ideal in each term, and extract weight zero. The primitive
block containing that entry passes through the corresponding
edge length. If it contains an adjacent entry, its trace is
in \(X\): in a triple one groups the two edges returning
to the starting vertex into one matrix corner, opposite the
chosen adjacent edge. The same reasoning with an outer entry
puts the trace in \(Y\). The remaining blocks have values
in \(T\). Conversely the generators of \(X,Y\) lie in the
respective entry ideals by expanding in group matrices.
An approximation in every truncation gives membership in
\(X\) or \(Y\) itself. These are closed in \(A\), since they
are finitely generated ideals of the compact Noetherian ring
\(T\). This proves the contractions.

Modulo \(X+Y\) the trace pairs only the three diagonal corners.
All four-by-four trace minors therefore vanish. Each generator
of \(X\) or \(Y\) is itself a four-by-four minor using the
three projectors and the two opposite off-diagonal entries
as its last row and column. Proposition~\ref{prop:origideal373}
gives \(X+Y=\mathfrak a\). Finally write a product from
\(X^2\) using the pair \(12,21\) for the first factor and
\(23,32\) for the second. The entries are scalars in \(A\).
Their product can be ordered as \((12)(23)(32)(21)\).
Grouping the first two and last two gives opposite outer
corners. This proves the last containment.
\end{proof}

\begin{theorem}[Odd descent on the integral trace quotient]
\label{thm:planedesc381}
The ring \(T/X\) carries the trace of an odd two-dimensional
representation of \(G_\QQ\) realized over \(A/J_1\), of
determinant \(\eta\). Its trace on \(G\) is \(t(h)-1\), where
\(t(h)=\operatorname{tr}M(h)\) is taken modulo \(X\).
After twisting by \(\Psi^{-1}\) it induces a surjection onto
\(T/X\) from the weight-augmentation quotient of the ordinary
two-dimensional pseudodeformation ring over \(\QQ\) with
residual trace \(1+\bar\eta\bar\Psi^{-2}\) and determinant
\(\eta\Psi^{-2}\). These statements hold over \(O\), including
on every nonreduced specialization of the trace quotient.
Every inverse-trace difference \(t(h)-t(h^{-1})\) belongs to
\(X\) in \(T\).
\end{theorem}
\begin{proof}
Over \(A/J_1\) the middle entry gives a character \(m\) of \(G\).
It is unramified above \(\ell\). In fact the fixed inertia
polynomial has roots \(\Psi(u),1,\Psi(u)^{-1}\) at every
inertia element \(u\), while \(m(u)\) is a root congruent
to one. If \(\bar\Psi(u)\ne1\), the other two factors evaluated
at \(m(u)\) are units, hence \(m(u)=1\). If \(\bar\Psi(u)=1\),
apply this argument to \(g\) and \(gu\) for an inertia element
\(g\) with \(\bar\Psi(g)\ne1\) at that place. At all other
finite places inertia acts through a group of order prime to
\(\ell\), whereas \(m\) takes values in principal units of a
complete local ring with finite residue field. Their finite
quotients have \(\ell\)-power order. Class field theory and
\(\ell\nmid |\operatorname{Cl}(K)|\) now give \(m=1\).
The outer matrix \(U(h)\) consequently has determinant one.
Writing \(J_2'\) for the two-by-two reversal matrix and
\(C_0=\operatorname{diag}(1,-1)\), direct multiplication gives
\[
 U(chc)=J_2'U(h)^{-\mathsf T}J_2'=C_0U(h)C_0 .
\]
Since \(G_\QQ=G\rtimes\langle c\rangle\), assigning \(C_0\)
to \(c\) extends \(U\) to \(W\) with determinant \(\eta\).
Traces on the other coset are differences of the outer
diagonal entries. Each such entry belongs to the image of
\(T\) by insertion of \(e_i\), hence all traces take values
in the faithfully embedded ring \(T/X\). This proves the
existence of its two-dimensional pseudocharacter there by
the matrix trace identities. On \(G\), a determinant-one
two-by-two matrix and its inverse have the same trace.
This proves the last assertion integrally.

We verify the local condition over the full coefficient ring.
At a place over \(\ell\) choose \(g\) in inertia with the
prescribed three roots distinct modulo the maximal ideal.
The outer \(\Psi(g)\)-eigenprojector is a direct summand of
rank one by interpolation and reduction modulo the maximal
ideal. In the ordinary word identity of
Theorem~\ref{thm:word68} take the last two letters equal
to \(g\). On this eigenline these last two factors are units,
so \(U(u)\) acts by \(\Psi(u)\) there for all inertia \(u\).
The inertia groups of \(K\) and \(\QQ\) at this prime agree
since \(K\) is unramified. Every local element of \(G_\QQ\)
preserves the line: conjugating it gives an inertia eigenline
of the same global character, which must lie in the same
rank-one eigensummand at \(g\); applying the inverse gives
equality. The local line character has values in \(T/X\):
its value at \(v\) is the trace of its eigenprojector times
\(W(v)\), and the projector is a polynomial in \(W(g)\)
over \(O\). The other local character is obtained from
the determinant. Thus the local trace of \(W\Psi^{-1}\)
splits over the trace quotient itself, and the character
lifting one is unramified.

The twist has odd, de Rham determinant \(\eta\Psi^{-2}\)
and residual inertia character \(\bar\kappa_\ell^{-2r}\ne1\).
The residual outer lines have characters \(\bar\Psi\) and
\(\bar\eta\bar\Psi^{-1}\) before twisting, by their values on
\(G\) and on \(c\). Ramification is confined to the original
finite set together with the primes of \(K\)'s discriminant.
Take this set in Pan's ordinary ring over the totally real
abelian field \(\QQ\). Its universal trace maps to the twisted
trace just constructed. Local reciprocity and the unramified
line make the action of the weight algebra
\(O[[1+\ell\ZZ_\ell]]\) factor through augmentation.
Pan's Theorem~5.1.2(1) makes this specialized source finite
over \(O\). The image contains all \(t(h)\) on \(G\), since
one restores their finite or cyclotomic scalar twist and adds
one. It is closed since it is finite over the compact ring
\(O\). The traces generate \(T/X\), which proves surjectivity.
~\cite{VIII-VI-VI-pan40}
\end{proof}

For the second quotient put \(A'=A/J_2\).
Let \(V_+,V_-\) denote its submodules of weight \(1,-1\),
respectively. Here and below tensor products of these modules
are over \(O\).

\begin{proposition}[Multiplication in the quotient with vanishing outer entries]
\label{prop:ordinaryplane381}
The ring \(A'\) has an \(O\)-module decomposition
\[
       A'=O\oplus V_+\oplus V_-\oplus L,\qquad
       T/Y=O\oplus L,\qquad L=V_+V_- .
\]
Each of \(V_+^2,V_-^2\) is zero. The product of \(L\) with the
entire augmentation ideal is zero. In particular \(L\) is a
quotient of \(V_+\otimes V_-\). The full negative module is
\[
 V_-\ \simeq\ S(\Psi^{-1})^\vee\oplus S(\eta\Psi^{-1})^\vee,
 \qquad U^\vee=\operatorname{Hom}_O(U,D).
\]
All these modules are finite over \(O\), and \(L\) has finite
length. The two rings over the common augmentation satisfy
\[
        T/(X\cap Y)\ \simeq\ (T/X)\times_O(O\oplus L).
\]
\end{proposition}
\begin{proof}
In \(A'\), the \(13\)-entry of \(M(hk)\) is zero for all
\(h,k\in G\). Multiplying matrices leaves only
\(M(h)_{12}M(k)_{23}\) in that entry. Reflection allows any
entry of weight one to be used in either position. Hence every
product of two such entries vanishes. The \(31\)-entry proves
the same for weight minus one. Let \(I\) be the ideal generated
by the adjacent entries. It satisfies \(I^3=0\).
Modulo \(I\) the representation is diagonal. Its three
characters agree with the prescribed ones on inertia at
\(\ell\): apply the separating inertia test in the proof
of Proposition~\ref{prop:origideal373}, which works for
diagonal lifts over a ring with nilpotents. Explicitly, the
traces at \(g^j u\), \(0\le j\le2\), give an invertible
Vandermonde matrix in the three inertia eigenvalues of \(g\).
The tame inertia and class group argument in the preceding
proof then fixes the three characters globally. Thus
\(A'/I=O\), by generation by entries, and \(I\) is the
augmentation ideal.

The module \(I/I^2\) is generated over \(O\) in degrees
\(1,-1\). The module \(I^2\) is generated in degree zero by the mixed
products since \(I^3=0\). These modules are finite over \(O\)
by Noetherianity. The weight decomposition therefore has
just the stated degrees, and \(L=I^2\) is the entire
degree-zero augmentation. Lemma~\ref{lem:pairedcorners381}
identifies that invariant ring with \(T/Y\).
Passing from \(A\) to \(A'\) leaves its degree minus one
augmentation cotangent module intact: the removed generators
have weights \(2,-2\), and modulo the augmentation square
their multipliers reduce to \(O\). Since \(I^2\) has weight
zero, the surviving negative cotangent module is \(V_-\).

We recall how its dual identification keeps track of its two
summands. A square-zero lower deformation is a cocycle in
the two adjacent lower entries after dividing out the split
diagonal. The difference and the sum of the two lower matrix
units have coefficient characters \(\Psi^{-1}\) and
\(\eta\Psi^{-1}\). The mark at \(c\) fixes the odd coboundary,
and the remaining marking equation fixes the even coboundary
since \((\Psi(g_*)+1)(\Psi(g_*)^{-1}-1)\) is a unit.
The local word condition gives strictness at \(\ell\). At
the other places averaging on the prime-to-\(\ell\) inertia
quotient gives the unramified condition, as proved for the
marked cocycles in Corollary~\ref{cor:dual68}. Its calculation
applies with the sign of the present \(K\), since the extension is
unramified at \(\ell\) and the same kernels are imposed
away from \(\ell\). Taking divisible coefficients and dualizing
gives exactly the two finite summands shown. In particular a
tensor product of \(V_-\) with the finite module \(V_+\) has
finite length, which proves the assertion for \(L\).
Finally the image of \(T\to T/X\times T/Y\) consists of pairs
with the same residue modulo \(X+Y=\mathfrak a\): adjust lifts
whose difference lies in \(X+Y\) by its two summands.
Its kernel is \(X\cap Y\), proving the identity.
\end{proof}

\begin{theorem}[The complementary strict group as an integral summand]
\label{thm:adjtorsion381}
Let \(C=S(\Psi^{-1})^\vee\), \(C_1=S(\eta\Psi^{-1})^\vee\)
and \(h=\dim_k(V_+/\mathfrak m_O V_+)\). The mixed trace map
of Theorem~\ref{thm:mixedtrace359} factors through
\(T/Y=O\oplus L\), and its map \(L\to C\) has an
\(O\)-linear section. More precisely,
\[
  \len_O C\ \le\ \len_O L\ \le\
          h\bigl(\len_O C+\len_O C_1\bigr).
\]
If \(h=1\) and \(C_1=0\), it identifies \(L\) with \(C\).
Both the summand and the upper bound refer to the character
of the given \(E_x\) on the original interval.
\end{theorem}
\begin{proof}
For \(C=0\) the section is the zero map and the length estimate
follows from the preceding proposition. Suppose \(C\ne0\).
Use the actual mixed lift from Theorem~\ref{thm:mixedtrace359}
over \(A_m=(O\oplus C)[y]/(y^2)\), where \(C^2=0\).
Give \(y\) degree one and \(C\) degree minus one. Before
the final change of frame its matrix on \(h_0\in G\) is
\[
 \exp\bigl(z(h_0)X_-+y b(h_0)X_++y w(h_0)H_0\bigr)d(h_0),
\]
where \(X_-=E_{21}-E_{32}\), \(X_+=E_{12}+E_{23}\),
\(H_0=[X_-,X_+]\) is diagonal, and \(z,w\) take values in
\(C\). In this formula \(E_{ij}\) are the matrix units.
The exponential stops after degree two in the nilpotent ideal
\((C,y)\). Two lower steps vanish by \(C^2=0\), and two
upper steps vanish by \(y^2=0\). A mixed pair of steps is
diagonal. Thus every outer entry is zero. The remaining marking
conjugates by \(1+y t(E_{12}-E_{23})\) for a scalar \(t\).
It respects the indicated weights. Its mixed terms with a
lower step are diagonal and its terms with a second upper step
vanish. Consequently the marked matrix map factors through
\(A'\) as a weight-preserving map. It sends
\(V_-\) into \(C\), \(V_+\) into \(Oy\), and \(L\) into \(yC\).
On traces the mixed construction has image \(O\oplus yC\),
so the last map is onto.

The map \(V_-\to C\) is projection onto the first of the two
summands in the preceding proposition, under its natural
cocycle identification. In fact setting \(y=0\) leaves just
\((1+z(h_0)X_-)d(h_0)\). The cocycle \(z\) used in the mixed
construction is the universal normalized cocycle for the
\(\Psi^{-1}\) strict classes, as constructed in
Theorem~\ref{thm:evenroot358}. For each functional \(C\to D\)
its specialization is precisely that class on the difference
of the two lower matrix units and has zero component on
their sum. This identifies the dual map as inclusion of the
first summand, hence identifies the original map by finite
double duality. Write \(j:C\to V_-\) for the resulting section.

The image of \(V_+\) is an \(O\)-submodule of \(Oy\).
If it were contained in \(\mathfrak m_Oy\), the equality
\(L=V_+V_-\) would put its image in
\(y\mathfrak m_O C\), strictly smaller than \(yC\) by
Nakayama's lemma. Choose \(u\in V_+\) mapping to \(y\).
Multiplication gives the required section
\(a\mapsto u\,j(a)\) from \(C\) to \(L\), after identifying
\(yC\) with \(C\). The section may depend on \(u\).
Finally \(h\) elements generate \(V_+\) over the discrete
valuation ring. Their products with \(V_-\) give a surjection
from \(h\) copies of \(C\oplus C_1\) onto \(L\), proving the
upper bound and the last assertion.
\end{proof}

Thus the odd two-dimensional quotient kills all inverse-trace
differences, even integrally, while on the other quotient their
image contains the full complementary group under the mixed
specialization. The splitting computes its position inside a
finite module. A uniform upper bound for its order in
Proposition~\ref{prop:res23} still requires arithmetic
information about that module or its summand.


\subsection{The integral product of the two adjacent root modules}
\label{sec:tensortrace382}

Keep the ring \(A'=A/J_2\) and the notation of
Section~\ref{sec:planedesc381}. The multiplication
\(V_+\otimes_O V_-\to L\) there is surjective. We compute
its kernel and the elementary divisors of both factors.
All calculations use the same \(x,\ell,r,\alpha,K\) as the
mixed trace construction. Write \(k=O/\ell O\). The coefficient
ring \(O\) is unramified. For any character \(\xi\) write
\(S(\xi)\) for the strict group over \(\QQ\) with coefficients
\(D(\xi)\), unramified away from \(\ell\) and locally zero at
\(\ell\), where \(D=\operatorname{Frac}(O)/O\).
Let \(F=E_*K\), with \(E_*\) the common period overfield of
conductor dividing \(N_*\) on the chosen interval. Thus
\([F:\QQ]\) is prime to \(\ell\), and \(F\) is unramified at
\(\ell\). For an \(O[\operatorname{Gal}(F/\QQ)]\)-module \(M\)
write \(M_\theta\) for the summand with action by the
specified character \(\theta\).

\begin{proposition}[Elementary divisors of the positive root module]
\label{prop:upperorders382}
Write \(P\) for the torsion submodule of \(V_+\). The module
\(V_+\) has free rank one. For \(r\ge2\), its torsion has
the elementary divisors of
\[
 S(\kappa_\ell\Psi^{-1})\ \oplus\ S(\eta\kappa_\ell\Psi^{-1}).
\]
For \(r=1\), put \(B_F=\operatorname{Cl}(F)\otimes_\ZZ O\).
The elementary divisors are those of
\[
       (B_F)_{\alpha^{-1}}\oplus (B_F)_{(\eta\alpha)^{-1}}.
\]
In particular the same class number \(h_F\) annihilates this
torsion for all the first twists on the interval, and
\(\log h_F\ll [F:\QQ]\log(2sN_*)\).
\end{proposition}
\begin{proof}
Put \(O_j=O/\ell^j O\). A map \(V_+\to O_j\) is a
square-zero deformation at the augmentation supported on the
upper adjacent entries. Passage to \(A'\) leaves the positive
cotangent module unchanged for the same reason as the negative
one in Proposition~\ref{prop:ordinaryplane381}.
After dividing out the split diagonal, the two upper coordinates
are one-cocycles. In the directions
\(E_{12}-E_{23},E_{12}+E_{23}\) their coefficient characters
are \(\Psi,\eta\Psi\), since the sign acts on the Lie algebra
by minus the reflected transpose. Global cohomology classes
determine their marked representatives uniquely. The mark at
\(c\) determines the odd coboundary, as two is a unit. An even
coboundary with scalar parameter \(v\) changes the upper
mark at \(g_*\) by \(\pm(1-\lambda)(1+\lambda^{-1})v\),
where \(\lambda=\Psi(g_*)\). The first and second upper
entries of the even root times \(d(g_*)\) are \(1\) and
\(-\lambda^{-1}\), which gives the displayed coefficient.
It is a unit. These are all the marks in this weight.

Away from \(\ell\) the cocycles must give unramified classes.
Indeed vanishing on the original inertia kernel is equivalent
to being a coboundary on inertia, by averaging over its
prime-to-\(\ell\) image. At \(\ell\) the upper triangular
matrix with either such cocycle satisfies the ordinary word
and the fixed inertia polynomial. For \(r\ge2\) this gives
the selected local condition by Theorem~\ref{thm:smooth69}, since
the Tate duals of the positive root entries have inertia
characters \(\bar\kappa_\ell^{1-r},\bar\kappa_\ell^{1-2r}\),
which are nontrivial in this range. Thus there is no further
local condition on the two upper cohomology classes at \(\ell\).

For \(\gamma=\Psi,\eta\Psi\), give \(O_j(\gamma)\) these local
conditions and call the resulting Selmer group \(H_j(\gamma)\).
The orthogonal conditions on \(O_j(\kappa_\ell\gamma^{-1})\)
are unramified away from \(\ell\) and zero at \(\ell\).
At a finite place away from \(\ell\), two unramified classes
pair to zero since their cup product comes from the residue
group of cohomological dimension one. The size of each
unramified group equals the size of its local invariants
by the kernel and cokernel formula for Frobenius minus one.
Local duality and the local Euler characteristic therefore
give equality with the orthogonal complement, including at
ramified places. The trace pairing of the unramified coefficient
ring with \(\ZZ_\ell\) is perfect, identifying the indicated
coefficient twist with the Tate dual for this calculation.

Subtract the lengths for these two Selmer conditions in
the Poitou--Tate size formula. It gives the difference of
the global invariant lengths plus the sum over places of
the length of the imposed local condition minus the local
invariant length. This formula follows by taking orders in
the Poitou--Tate sequence and applying the global and local
Euler characteristics as in Proposition~\ref{prop:strictclasses365}.
Both global invariant groups vanish by their inertia powers
at \(\ell\). The terms at all finite places away from \(\ell\)
vanish by the Frobenius calculation. At \(\ell\) the term is
\(j\) by the local Euler characteristic, since the local
dual invariants vanish. At the real place the term is \(-j\)
for an even \(\gamma\) and zero for an odd one. The dual
finite-coefficient Selmer group is
\(S(\kappa_\ell\gamma^{-1})[\ell^j]\). In fact the multiplication
sequence on \(D(\kappa_\ell\gamma^{-1})\) identifies the global
first cohomology at finite coefficients with the indicated
torsion because there are no global invariants. It does
the same at \(\ell\). At other places the inertia invariants
of the divisible module are divisible by prime-to-\(\ell\)
averaging, giving injectivity on inertia in the same sequence.
This proves, for \(r\ge2\),
\[
 \len_O H_j(\gamma)
    = j\,\mathbf1_{\gamma(c)=-1}
      +\len_O S(\kappa_\ell\gamma^{-1})[\ell^j].
\]
The strict groups here are finite by the negative-twist
finiteness in Proposition~\ref{prop:finite15} after descent
to the abelian field \(F\).~\cite{VIII-VI-V-milne98}

When \(r=1\), write \(\gamma=\kappa_\ell\theta\) with
\(\theta=\alpha\) or \(\eta\alpha\). The local condition at
\(\ell\) in the upper entries is now the unit Kummer condition.
Indeed Theorem~\ref{thm:first72} imposes precisely the vanishing
of the valuation coordinates. After the unramified splitting
of the finite characters, the two coordinates are the ordinary
Kummer valuations on the adjacent lines. In the nonsplit
case the two lines are interchanged and the same description
follows over the unramified quadratic extension by Shapiro's
lemma, as in that theorem. The changes between the root
bases use only units, including two. Restriction to the
completion of \(F\) detects these conditions, since the degree
is prime to \(\ell\) and valuations are multiplied by the
ramification index under extension. Finite group averaging
also identifies the global cohomology with the corresponding
character part over \(F\).

Kummer theory now identifies the Selmer group with the
\(\theta^{-1}\)-part of the global elements whose valuations
at every finite place are zero modulo \(\ell^j\), taken
modulo \(\ell^j\)-th powers and extended to \(O\).
Away from \(\ell\), extracting roots of units is unramified,
while an element whose valuation is not a multiple of \(\ell^j\)
cannot have an \(\ell^j\)-th root in an unramified extension.
The ideal of a permissible element is \(\mathfrak b^{\ell^j}\).
Sending the element to the class of \(\mathfrak b\) gives
the exact sequence with kernel
\((\mathcal O_F^\times\otimes_\ZZ O)/\ell^j\) and quotient
\(B_F[\ell^j]\), before taking the character summand.
Surjectivity follows by taking a generator of any principal
\(\mathfrak b^{\ell^j}\). The character projection is exact.
There is no \(\ell\)-power torsion in the units, since
\(F\) is unramified at \(\ell>2\).
Dirichlet's unit theorem identifies their rational character
with the permutation character on complex-conjugate pairs
of embeddings minus the trivial character. This gives
multiplicity zero for the odd \(\alpha^{-1}\) and one for
the even \((\eta\alpha)^{-1}\), which is nontrivial since
\(K\) is disjoint from \(E_*\). Hence the lengths have the
same free contribution \(j\) in total as above and the
displayed class group torsion contributions at each \(j\).

In either case the sum of the two Selmer lengths equals
\(\len_O\operatorname{Hom}_O(V_+,O_j)\) by the marked
description. For a finitely generated module over this
discrete valuation ring, successive differences of these
lengths count the free rank plus the number of torsion
elementary divisors with exponent at least \(j\).
Their limiting value is the rank. This proves the assertions
about \(V_+\). Finally Minkowski's bound represents every
ideal class by an integral ideal of norm at most
\(C^{[F:\QQ]}|\operatorname{disc}F|^{1/2}\) for an absolute
\(C\). The number of ideals up to \(X\) is at most
\(X^2\zeta_F(2)\le X^2\zeta(2)^{[F:\QQ]}\). The abelian
conductor-discriminant formula bounds the discriminant by
\((sN_*)^{[F:\QQ]}\), giving the asserted uniform estimate.
\end{proof}

\begin{lemma}[Diagonal corrections for both signs]
\label{lem:diagonals382}
Let \(Z_0\) be a finite \(O\)-module with trivial Galois
action, and let \(S\) contain infinity, \(\ell\) and all
primes of the prescribed ramification. For each of the
modules \(B=Z_0,Z_0(\eta)\) the map on second cohomology
from \(G_{\QQ,S}\) to the local groups is injective.
The map on first cohomology is onto the direct sum at
the finite places in \(S\) of local first cohomology
modulo its unramified subgroup.
\end{lemma}
\begin{proof}
For the sign summand these are exactly the two assertions of
Lemma~\ref{lem:diagonalpatch359}, which uses the class number
of the same \(K\). For the trivial action decompose the
underlying finite abelian group into cyclic summands of
orders \(\ell^j\). The Tate dual of one such summand is
\(\mu_{\ell^j}\). A global dual class orthogonal to every
local unramified subgroup at \(S\) is a rational Kummer
element whose valuation is divisible by \(\ell^j\) at each
finite place of \(S\), by local reciprocity. The same
holds outside \(S\) by unramifiedness away from \(\ell\).
Its absolute value is therefore an \(\ell^j\)-th power in
\(\QQ^\times\), as is its sign. The dual class is zero.
Poitou--Tate duality now gives the first-cohomology
surjectivity because the orthogonal complement of the global image
is supplied by the global dual classes just tested.
An everywhere locally trivial dual class belongs to the same
zero group, giving injectivity in degree two. The real place
has zero modified cohomology at odd \(\ell\).~\cite{VIII-VI-V-milne98}
\end{proof}

\begin{theorem}[The full integral multiplication map]
\label{thm:tensortrace382}
On every large-prime triple under consideration, multiplication
is an isomorphism
\[
             V_+\otimes_O V_-\ \xrightarrow{\ \sim\ }\ L .
\]
It identifies \(A'\) with the tensor product over \(O\) of the
two square-zero algebras \(O\oplus V_+\) and \(O\oplus V_-\)
as augmented graded algebras. It also identifies the invariant
trace ring \(T/Y\) with \(O\oplus(V_+\otimes_O V_-)\).
\end{theorem}
\begin{proof}
Set \(Z_0=V_+\otimes_O V_-\), a finite-length module, and
form the tensor product algebra \(B\) in the statement.
Its augmentation ideal has cube zero. Products of two
positive entries or two negative entries are zero. Multiplying
one of each gives its tensor in \(Z_0\), which is killed by
the augmentation ideal. Over each \(O\oplus V_\pm\) take
the representation supplied by the respective projection
from \(A'\), killing the other weight and \(L\).
Its matrix part is \((1+p(g))d(g)\), respectively
\((1+m(g))d(g)\), for \(g\in G_{\QQ,S}\), with the original
sign in the semidirect product. Here \(p,m\) are one-cocycles
in the upper and lower adjacent matrices, with their split
adjoint actions. We shall realize both projections over \(B\).

Seek a matrix part \(\exp(p(g)+m(g)+w(g))d(g)\), with \(w\)
in trace-zero diagonal matrices of coefficient \(Z_0\).
On this latter module the split action is trivial on the
diagonal vector \((1,0,-1)\) and has character \(\eta\)
on \((1,-2,1)\). These form a basis since two is invertible.
Products of an upper and a lower step are diagonal, so all
commutators of the two directions lie in this module.
The central extension in the unipotent matrices has the multiplication rule
\(\exp U=1+U+U^2/2\) and
\(\exp U\exp V=\exp(U+V+[U,V]/2)\), by multiplication.
Thus the defect for \(w=0\) is a continuous two-cocycle
with the indicated diagonal coefficients. Associativity gives
the cocycle identity, and adding a diagonal one-cochain changes
the defect by its coboundary. All resulting matrices have
determinant one, since the exponent has trace zero. This
also follows by expanding the determinant through degree two.

At \(v\ne\ell\) in \(S\) both cocycles factor through the
local quotient by the prescribed inertia kernel. Its finite
inertia has order prime to \(\ell\) and its residue group is
procyclic. The inflation sequence therefore gives zero
second cohomology there for the diagonal coefficients.
Choose a local solution on this quotient. At \(\ell\) the
universal negative cocycle is a coboundary, by strictness
in the finite module \(V_-\). The marked classification in
Corollary~\ref{cor:dual68} applies to every finite coefficient
module. Conjugate the pure upper lift over \(B\) by a
matrix \(1+v_-\) with lower adjacent entries in \(V_-\)
giving that coboundary. This gives a local solution with
the same two first-order parts. The difference in exponential
coordinates is a trace-zero diagonal with coefficients in
\(Z_0\). The change of basis is taken in the semidirect
product, so the assertion also holds at a nonsplit place.

The two-cocycle is now locally trivial everywhere. Apply
Lemma~\ref{lem:diagonals382} to both signs to obtain a
global solution. Differences from the local solutions are
one-cocycles. Subtract a global one-cocycle so that these
differences have unramified cohomology class at every finite
place of \(S\). At \(v\ne\ell\) the remaining difference is
a coboundary on inertia and vanishes on the prescribed
inertia kernel, on which its coefficient action is trivial.
At \(\ell\) the diagonal action is unramified, hence the
remaining difference vanishes on inertia itself. At \(c\)
the pure cocycles are zero by the marks. The relation
\(c^2=1\) makes \(w(c)\) anti-invariant, and a diagonal
coboundary kills it since two is a unit. This conjugation
preserves the conditions just imposed.

To check the entire local component at \(\ell\), undo the
lower change of basis used above. The resulting upper
triangular lift differs from the pure upper lift only by
the diagonal cocycle with unramified class. It has the same
inertial diagonal, since the coboundary at inertia is zero.
For \(r\ge2\) the ordinary word already gives the component
by Theorem~\ref{thm:smooth69}. For \(r=1\), use the valuation
description of Theorem~\ref{thm:first72}. The adjacent
upper cocycle takes its values in the submodule \(V_+\subset B\)
of the additive coefficients. The new diagonal acts on
this submodule exactly as the original diagonal, since
\(Z_0 V_+=0\). Its cocycle there is the same as in the pure
upper lift, whose valuation coordinates are zero. Mapping
that coefficient module into the full adjacent module
commutes with valuation by Lemma~\ref{lem:kummer72}.
In the nonsplit case perform this check on the first line
over the unramified quadratic extension as in the proof of
Theorem~\ref{thm:first72}. These assertions can equally be
tested modulo all powers of the maximal ideal of \(B\).
They prove membership in the selected component over \(B\)
for the first twist as well.

All outer matrix entries vanish, since a path of length two
to an outer entry uses two steps of the same sign. The upper
and lower entries of the marking equation at \(g_*\) remain
those of the original pure lifts because every mixed correction
is diagonal. The constructed lift thus determines a map
\(A'\to B\). Its coefficients have the indicated torus
weights, and it induces the identity on \(V_+\) and \(V_-\)
after projecting to the pure square-zero algebras. These
are also the whole degree-one and degree-minus-one parts
of \(B\), so the map on \(L\) sends the product of any two
such elements to their tensor. This is a left inverse to
the multiplication surjection of
Proposition~\ref{prop:ordinaryplane381}. Hence multiplication
is bijective. The decompositions and the products already
computed identify the two rings as asserted.
\end{proof}

\begin{corollary}[Lengths and exponents of the additional torsion]
\label{cor:tensorexcess382}
Put \(C=S(\Psi^{-1})^\vee\), \(C_1=S(\eta\Psi^{-1})^\vee\)
and \(M=C\oplus C_1\). Write
\(d_j(H)=\dim_k(\ell^{j-1}H/\ell^j H)\) for a finite
\(O\)-module \(H\). The kernel of the mixed map \(L\to C\)
has, as an \(O\)-module, the elementary divisors of
\(C_1\oplus(P\otimes M)\). In particular
\[
 \len_O L-\len_O C
    =\len_O C_1+\sum_{j\ge1}d_j(P)d_j(M).
\]
Any power of \(\ell\) killing both \(C_1\) and \(P\) kills
this entire kernel. At the first twist, if \(\ell\nmid h_F\),
the only additional summand is \(C_1\).
\end{corollary}
\begin{proof}
If \(C\ne0\), the positive map in the proof of
Theorem~\ref{thm:adjtorsion381} is a surjection onto \(Oy\).
By Proposition~\ref{prop:upperorders382} it has kernel \(P\),
and a preimage of \(y\) splits off the free summand.
On the negative factor the map is projection of \(M\)
onto \(C\). Theorem~\ref{thm:tensortrace382} shows that
on the entire tensor it is the product of these two maps.
Its kernel has the asserted summands. If \(C=0\), choose
any free generator modulo \(P\). The entire tensor then
has the asserted elementary divisors. Two cyclic modules
of exponents \(a,b\) have tensor length \(\min(a,b)\).
Summing this minimum as a sum of indicators over \(j\)
gives the formula. The annihilation statements follow from
the tensor description and Proposition~\ref{prop:upperorders382}.
\end{proof}

The length formula distinguishes the orders at weight \(r\)
from the additional contributions detected at weight \(r-1\).
Even when both additional groups vanish, the summand
\(S(\Psi^{-1})^\vee\) at weight \(r\) is still present at
its full order. An independent upper estimate for the weighted sum
in Proposition~\ref{prop:res23} could be obtained by bounding this same
trace quotient. The length identity leaves the size of the surviving
summand undetermined.


\subsection{Counting lower triangular lifts on the period interval}
\label{sec:lowerfibers384}

Keep the interval \([V,V^{3/2}]\), period fields \(E_x\subset E_*\)
and fixed \(K=\QQ(\sqrt{-s})\) of
Proposition~\ref{prop:quadratic345}. For a triple on this interval
write \(\Psi=\alpha\kappa_\ell^r\), with \(\Psi(c)=1\),
\(1\le r<t(x)\), and \(\ell>\exp(x^\varepsilon/12)\). Let \(\eta\)
be the character of \(K\) and let \(O=\ZZ_\ell[\alpha]\) have
fraction field \(J\) and residue field \(k\). In particular \(O\)
is unramified. The columns pair each character with its root matrix.
\[
\begin{array}{c|c|c}
 a=\Psi^{-1} & b=\eta\Psi^{-1} & d_2=\eta\Psi^{-2}=ab\\
 X=E_{21}-E_{32} & Y=E_{21}+E_{32} & Z=E_{31}.
\end{array}
\]
For any finite coefficient module \(W\), denote by \(S(W)\) the
kernel of restriction in \(H^1(\QQ,W)\) to \(G_{\QQ_\ell}\)
and to inertia at every other finite prime. Write \(S(\xi)\) for
the same kernel with divisible coefficient \((J/O)(\xi)\).
We use \(B_j=O/\ell^{j+1}O\), \(I_j=\ell B_j\) and
\(S_j(\xi)=S(I_j(\xi))\), for \(j\ge1\).

Let \(B\) be a finite local Artinian \(O\)-algebra with residue
field \(k\), and put \(I=\mathfrak m_B\). Denote by
\(\mathcal L(B)\) the lifts in the original marked ordinary
matrix problem of Section~\ref{sec:planedesc381} that are lower
triangular on \(G_K\). Denote by \(\mathcal O(B)\) those whose
matrices on \(G_K\) preserve the fixed antidiagonal form \(J_3\)
in the marked basis. The sign acts on matrices by the involution
\(U\mapsto J_3 U^{-\mathsf T}J_3\). Both conditions consequently
make sense on the full polarized representation. We determine the
fibers of \(\mathcal L(B)\) over the strict classes for \(a\).

\begin{lemma}[Local splitting and the two marks]\label{lem:triangularmark384}
Put \(d(g)=\operatorname{diag}(\Psi(g),1,\Psi(g)^{-1})\), and
write \(\delta(g)=(d(g),\eta(g))\) in the polarized group. Let
\(U(I)\) consist of lower triangular unipotent matrices whose
off-diagonal entries belong to \(I\), with \(G_\QQ\) acting by
conjugation through \(\delta\). Then \(\mathcal L(B)\) identifies
with the \(U(I)\)-conjugacy classes of lifts \(v(g)\delta(g)\),
where \(v\) is a \(U(I)\)-valued cocycle, locally a coboundary
at \(\ell\) and on inertia at every \(p\ne\ell\). Each such class
has a unique marked representative.
\end{lemma}
\begin{proof}
Over \(G_K\) the three diagonals of any triangular lift are
characters. Locally above \(\ell\) the prescribed ordinary flag
and the standard lower triangular flag are opposite on reduction.
They remain opposite over \(B\). Indeed the sums of complementary
steps are isomorphisms because their determinants are units.
Intersecting the ordinary step of rank \(i\) with the lower step
spanned by \(e_i,\ldots,e_3\) gives a stable line projecting
isomorphically to \(B e_i\) along the smaller lower step.
Its generator with \(i\)-th coordinate one reduces to \(e_i\).
These three generators make a lower unipotent change of basis.
The ordinary flag used here exists over \(B\) itself by the
separating inertia projectors and the word identity of
Lemma~\ref{lem:flag69}. In this basis the inertia characters on
the three lines are respectively \(\Psi,1,\Psi^{-1}\).
A triangular conjugation preserves the diagonal. At the other
finite places the prescribed inertia quotient has order prime
to \(\ell\). The ratio of each diagonal to its indicated
character takes inertia there into \(1+I\), an \(\ell\)-group,
and is trivial. Thus each ratio gives an everywhere unramified
\(\ell\)-power character of \(G_K\). It is trivial by class
field theory and \(\ell\nmid|\operatorname{Cl}(K)|\).

The local change of basis at \(\ell\) now conjugates the lift
to \(\delta\) on \(G_{K,w}\) for the chosen place \(w\).
The stabilizer of \(\delta|_{G_{K,w}}\) in \(U(I)\) is trivial.
The inertia element \(g_*\) has diagonal eigenvalues with pairwise
unit differences. A matrix commuting with it is diagonal.
For an element of \(G_{\QQ_\ell}\) outside \(G_{K,w}\), if
there is one, the cocycle identity and normality show that its
unipotent factor after this conjugation belongs to that
stabilizer. This proves the full local splitting. At \(p\ne\ell\)
the cocycle is trivial on the prescribed inertia kernel.
Any cocycle on the remaining prime-to-\(\ell\) quotient
is a coboundary. To see this, project to the two adjacent
entries and average, then remove the residual central cocycle
by the same argument. Conversely an inertia coboundary vanishes
on the prescribed kernel, where the action is trivial. At \(\ell\)
a conjugate of the split lift belongs to the selected ordinary
component, also at \(r=1\). Indeed conjugate the split lift by
the universal lower unipotent matrix over \(O[[t_1,t_2,t_3]]\).
Its evaluations at \(t_i\in\ell O\) are conjugates of the same
crystalline lift. A power series vanishing at all these points
is zero, by the one-variable identity theorem in each variable.
The family therefore factors through the reduced crystalline
closure with the ordinary equations. Its kernel is prime and
is contained in the kernel of split evaluation. It contains
a minimal prime through that evaluation, and hence the
intersection defining the selected local quotient in
Section~\ref{sec:original68}. Every lower unipotent matrix
congruent to one over \(B\) is obtained by specialization.

At a marked lift the matrix at \(c\) is one. As \(G_K\) and \(c\)
generate \(G_\QQ\), the fixed diagonal on \(G_K\) gives the expression
\(v\delta\) on the whole group. Conversely take a cocycle
satisfying the local conditions. First conjugate the lift's value at
\(c\) to \(\delta(c)\). This is possible by the same averaging
argument on the adjacent quotient and the center, since two is
a unit. The remaining allowable conjugators in \(U(I)\) are
\(\exp(tX)\), \(t\in I\). In fact logarithm on a strictly lower
three-by-three matrix is \(N\mapsto N-N^2/2\) for \(1+N\),
and the differential of the involution fixes \(X\) and negates
\(Y,Z\). Write the adjacent coordinates of \(\log v(g_*)\) as
\(uX+wY\), with \(\lambda=\Psi(g_*)\). The remaining nonautomatic
mark equates the \(21\) and \(32\) entries of the matrix
\(v(g_*)d(g_*)\). Their difference is
\((\lambda+1)u+(\lambda-1)w\). Conjugation by \(\exp(tX)\)
adds \((1-\lambda^{-1})t\) to \(u\) and fixes \(w\).
The coefficient \((\lambda+1)(1-\lambda^{-1})\) is a unit,
so exactly one \(t\) gives the mark. The upper mark already
vanishes. This proves the identification, including uniqueness.
\end{proof}

\begin{theorem}[The fiber over each complementary class]
\label{thm:lowerfibers384}
Projection to the \(X\)-coordinate gives a surjection
\(\mathcal L(B)\to S(I(a))\) with a section taking values in
\(\mathcal O(B)\). For \(u\in S(I(a))\) choose its cocycle
representative \(x\) zero at \(g_*\). Define \(W_x=I\oplus I\)
with action
\[
 g(y,z)=\bigl(b(g)y,\ d_2(g)z-2x(g)b(g)y\bigr).
\]
The fiber above \(u\) is in bijection with \(S(W_x)\). It fits
into an exact sequence of finite \(O\)-modules
\[
 0\longrightarrow S(I(d_2))\longrightarrow S(W_x)
       \longrightarrow D_x\longrightarrow 0,\qquad
 D_x\subseteq S(I(b)).
\]
For \(B=B_j\) this gives, at every \(j\) and every \(u\),
\[
 |S_j(d_2)|\ \le\ \#\mathcal L(B_j)_u
       \ \le\ |S_j(d_2)|\,|S_j(b)|.
\]
The upper bound is an equality whenever
\(\ell^e u=0\), \(\ell^h S_j(b)=0\), and \(e+h\le j+1\),
with \(0\le e,h\le j\). In particular
\[
 \lim_{j\to\infty}\#\mathcal L(B_j)
       = |S(a)|\,|S(b)|\,|S(d_2)|.
\]
\end{theorem}
\begin{proof}
Every element of \(U(I)\) factors uniquely as
\(\exp(yY+zZ)\exp(xX)\). This follows by reading the two
adjacent coefficients and then the central coefficient. The
first factor is an abelian normal subgroup because \([Y,Z]=0\)
and \([X,Y]=-2Z\). Under conjugation by \(\delta(g)\), \(X,Y,Z\)
have the three characters in the table. The projection to \(x\)
is therefore an additive cocycle with character \(a\).
Any class with the stated local conditions has exactly one
representative zero at \(g_*\), since \(a(g_*)-1\) is a unit.
That representative vanishes on \(G_{\QQ_\ell}\) as its local
class is zero. The lift \(\exp(xX)\delta\) preserves \(J_3\),
and satisfies the local conditions of the lemma. It gives the
claimed section.

Fix this representative \(x\). In the factorization the other
two coordinates are a cocycle with values in the abelian normal
subgroup, acted upon through \(\exp(xX)\delta\). Conjugation sends
\(Y\) to \(Y-2xZ\), which gives exactly \(W_x\). Once \(x\)
is fixed, a conjugation preserving it has zero \(X\)-coordinate
because \(I(a)^{G_\QQ}=0\). Thus the equivalence on the other
two coordinates is precisely coboundary equivalence in \(W_x\).
At \(\ell\) this argument works on the decomposition group
as well, since \(x=0\) there and \(a(g_*)-1\) is a unit.
At another inertia group both the action through \(\delta\) and
the base lift through \(x\) factor through the prescribed
prime-to-\(\ell\) quotient. A cocycle in \(W_x\) is unramified
exactly when it vanishes on the kernel of that quotient. This
follows by averaging after passage to the quotient. Combined
with the lemma, these facts identify the fiber with \(S(W_x)\).

Use the coefficient sequence \(0\to I(d_2)\to W_x\to I(b)\to0\).
Invariants of \(I(b)\) vanish globally and at \(\ell\) by its
nontrivial residual inertia character. Consequently degree one
on the central module injects both globally and at \(\ell\).
At any other inertia group, every invariant in \(I(b)\) lifts
to an invariant in \(W_x\) by averaging a lift over the finite
prime-to-\(\ell\) quotient. This gives the same injectivity
on inertia. A central global class mapping into \(S(W_x)\)
therefore has every required local vanishing. Exactness in
global cohomology proves the sequence with \(D_x\) the image
of the last map. Finiteness follows at finite coefficients
from the finiteness of global cohomology with restricted
ramification.

For each of the three characters \(\xi\), identify \(I_j\)
with \((J/O)[\ell^j]\) by dividing its representatives by
\(\ell^{j+1}\). The coefficient sequence then identifies
\(S_j(\xi)=S(\xi)[\ell^j]\). Indeed global and \(\ell\)-local
invariants of \((J/O)(\xi)\) vanish by \(g_*\). At every other
inertia group its invariants are divisible, since they are the
image of a prime-to-\(\ell\) averaging idempotent. The same
argument applies to each subideal of \(I_j\). In particular if
\(\ell^eu=0\), the class has a representative in
\(\ell^{j+1-e}B_j\). Normalizing it at \(g_*\) keeps its values
in that ideal. Every class in \(S_j(b)\) comes from
\(\ell^{j+1-h}B_j\) when \(\ell^h S_j(b)=0\).
If \(e+h\le j+1\) the product of the two ideals is zero. Then
\(y\mapsto(y,0)\) embeds this last coefficient ideal with
character \(b\) into \(W_x\) equivariantly. It carries strict
classes to strict classes and yields \(D_x=S_j(b)\).
All three divisible strict groups are finite by restriction
to \(E_*K\), of degree prime to \(\ell\), and
Proposition~\ref{prop:finite15}. Taking \(j\) larger than their
exponents measured in powers of \(\ell\), and large enough for
the sum condition, gives the asserted eventual equality.
\end{proof}

\begin{proposition}[The total contribution of the fibers on one interval]
\label{prop:odddivisor384}
Set \(F=E_*K\), \(g=[F:\QQ]\), \(N=sN_*\), and
\(T=\lfloor V^{(3/2)401}\rfloor\). For the cyclotomic modulus
\(N\) of \(F\) use the nonzero critical determinants of
Lemma~\ref{lem:int13}, and put
\[
 H=\prod_{n=2}^{2T+1}|D_n(F)|,\qquad H'=h_F H,
 \qquad h_F=|\operatorname{Cl}(F)|.
\]
Then \(\log H'\ll_{\psi,\varepsilon,s}V^{1220}\).
Take all distinct triples \((\ell,r,[\alpha])\) from the
chosen interval, counting each residue Frobenius orbit once
after inflation to \(E_*\). At every triple take the supremum
over all depths \(j\ge1\) and all \(u\in S_j(a)\). One has
\[
 \sum_{\ell,r,[\alpha]}\ \sup_{j,u}
        \log\#\mathcal L(B_j)_u\ \le\ 17\log H .
\]
In particular the same bound holds for the sum of the logarithms
of \(\#\mathcal L(B_j)/|S_j(a)|\), with depths chosen independently
at each triple.
\end{proposition}
\begin{proof}
For clarity, the determinant at weight \(n\) is the product of
\(A_n L^{(N)}(1-n,\chi^{-1})\) over the complex characters of
\(\operatorname{Gal}(F/\QQ)\) of sign \((-1)^n\). Here \(L^{(N)}\)
omits Euler factors at primes dividing \(N\). Using two distinct
fixed primes \(P_1,P_2\) unramified in \(F\), prime to \(N\),
one may take
\[
 A_n=(P_1^{gn}-1)(P_2^{gn}-1)\,nN(n+1)! .
\]
Such primes can be fixed over all sufficiently distant intervals,
since the two moving prime factors of \(N_*\) tend to infinity.
The Bernoulli formula, nonvanishing by the functional equation,
and the size estimate of Lemma~\ref{lem:int13} apply to this
abelian CM field. They give
\(\log H\ll gT^2\{g+\log(2NT)\}\). The Minkowski ideal count
and the conductor--discriminant formula give
\(\log h_F\ll g\log(2N)\). The estimates
\(g\ll_\psi V^{3+12\varepsilon}\),
\(\log(2N_*)\ll_{\psi,\varepsilon}V^{3\varepsilon/8}\) and
\(T^2\le V^{1203}\) from Theorem~\ref{thm:conductor232}
give the asserted bound for \(0<\varepsilon<1/52\).

All primes \(\ell\) occurring in the proposition are prime to
\(Ng\) on a sufficiently distant interval. For a character
\(\beta\) of \(F\), restriction identifies
\(S(\beta^{-1}\kappa_\ell^{-m})\) with the \(\beta\)-part of the
strict group of \(F\) at weight \(m\) after extension of
coefficients. This follows by finite group averaging in global
and local restriction, on inertia as well as on the group at
\(\ell\). For \(m\ge1\), the sum of the \(O\)-lengths of those
parts with \(\beta(c)=(-1)^{m+1}\), weighted by their minimal
residue degrees and taken over residue orbits, is bounded by
\(v_\ell(D_{m+1}(F))\). Indeed Poitou--Tate duality embeds
the dual of the strict part into degree two. The refined
Coates--Sinnott theorem puts the critical special value times
the indicated integer in its zeroth Fitting ideal. The required
Iwasawa \(\mu\)-invariant vanishes since \(F(\mu_\ell)\) is
abelian over \(\QQ\). Each factor of the critical integral group
ring is an unramified discrete valuation ring because
\(\ell\nmid g\). Its Fitting valuation bounds the strict
length, and weighting by residue degrees gives the determinant
valuation. Lemma~\ref{lem:fit13} and its proof give these
statements with the imprimitive modulus and the precise
coefficient duality. The primary special-value theorem is
Theorem~6.11 of Greither and Popescu.
\href{https://arxiv.org/abs/1103.3069}{Authors' paper}.

For \(b,d_2\) the corresponding pairs \((m,\beta)\) are
\((r,\eta\alpha)\) and \((2r,\eta\alpha^2)\). Both signs are
critical. Counting geometric characters counts each residue
orbit with its residue degree, because \(\ell\nmid g\).
Passing from the minimal ring of \(\eta\alpha^2\) to that of
\(\alpha\) leaves its length unchanged, by finite free
unramified base change of the defining cochains and local maps.
The first character map is injective. The second has at most
sixteen preimages. Indeed restriction embeds the Galois group
of \(E_*\) into a product of two cyclic groups and two groups
of order two. Its subgroup killed by two has order at most
sixteen. The character group has the same two-torsion order by
the structure theorem for finite abelian groups. At any given
\(m\), the first list uses only \(r=m\), the second only
\(r=m/2\). The sum of \(\log|S(b)|+\log|S(d_2)|\) over the
triples is therefore at most \(17\log H\), since
\(\sum_\ell v_\ell(D)\log\ell=\log|D|\) for a nonzero integer.
Theorem~\ref{thm:lowerfibers384} bounds each fiber by these two
orders. The ratio in the last assertion is the arithmetic
mean of the fiber cardinalities.
\end{proof}

\begin{theorem}[Equality with the marked even orthogonal locus]
\label{thm:orthcollapse384}
For a triple in Proposition~\ref{prop:odddivisor384} with
\(\ell\nmid H'\), the two closed subfunctors \(\mathcal L\)
and \(\mathcal O\) coincide on all finite local Artinian
\(O\)-algebras. Their representing ring, including its
\(\ell\)-torsion, is
\[
  O[[T_1,\ldots,T_q]]/(\ell^{a_1}T_1,\ldots,\ell^{a_q}T_q),
 \qquad
 S(\Psi^{-1})\simeq\bigoplus_{i=1}^q O/\ell^{a_i}O .
\]
On this whole ring every trace on \(G_K\) is the split trace.
\end{theorem}
\begin{proof}
Proposition~\ref{prop:odddivisor384}, or its determinant argument
at this prime, gives \(S(b)=S(d_2)=0\). For clarity, each
local length in its proof is bounded by a valuation of a factor
of \(H\), and every such valuation is zero. By finite-coefficient
comparison the same vanishing holds with coefficients \(O/\ell^jO\).
Every finite \(O\)-module \(I\) is a sum of cyclic modules.
Thus \(S(I(b))=S(I(d_2))=0\). The fiber theorem shows that
every lower triangular lift equals the marked representative
of its orthogonal section. This proves \(\mathcal L\subseteq
\mathcal O\).

For the reverse inclusion let \(R\) be the quotient of the
original marked ring by the equations preserving \(J_3\).
It is local, complete and Noetherian and is augmented at \(d\).
Conjugation by \(\operatorname{diag}(t,1,t^{-1})\) preserves
the form and both marks. Use the integer weights of matrix
entries under this split torus, as in
Lemma~\ref{lem:pairedcorners381}. We show that the relative
cotangent space \(\mathfrak m_R/(\ell R+\mathfrak m_R^2)\)
is generated by entries of weight \(-1\).

A first-order perturbation has values, after division by
the split representation, in \(\mathfrak{so}(J_3,k)\).
This has basis \(X,E_{12}-E_{23},E_{11}-E_{33}\) because
two is a unit. There is no outer root at first order.
The sign involution acts identically on this Lie algebra.
The diagonal coordinate is a cocycle with trivial coefficients.
The prescribed ordinary inertia characters at \(\ell\) force
it to vanish there on inertia. Infinitesimal conjugation of a
diagonal matrix has zero diagonal variation. At every other
finite place the same vanishing follows by the prime-to-\(\ell\)
inertia quotient. An everywhere unramified additive character
of \(G_\QQ\) with coefficients in \(k\) is zero by class field
theory. Thus the diagonal contributes no generators.

The upper coordinate has character \(\Psi\). At the marks its
cohomology class determines its representative uniquely, as in
Proposition~\ref{prop:upperorders382}. For \(r\ge2\) that
proposition identifies the length of the finite-coefficient
upper Selmer group for this even character at depth one with
the length of \(S(\kappa_\ell\Psi^{-1})[\ell]\). Its pair
\((m,\beta)=(r-1,\alpha)\) is critical and \(m\ge1\).
The determinant \(D_r(F)\) in \(H\) thus gives zero length.
For \(r=1\) the character \(\alpha^{-1}\) is odd. Its unit
part in the Kummer calculation of that proposition vanishes,
and its only possible contribution comes from
\(\operatorname{Cl}(F)\otimes O\), which is zero since
\(\ell\nmid h_F\). Hence no upper tangent coordinate survives.
Together with the Lie algebra calculation, this proves the
claim about weights of the relative cotangent space. Indeed
the entries generate the original matrix ring and its quotient,
and maps to the dual numbers separate the cotangent space.

Choose entries of weight \(-1\) lifting a basis of that space.
Together with \(\ell\) they generate \(\mathfrak m_R\) by
Nakayama's lemma. Successive approximation in its powers and
completeness give a surjection from a power series ring over
\(O\) in these entries onto \(R\). On each Artinian truncation
their monomials have nonpositive weights. Extracting the
coefficient of \(t^w\) in the torus action kills their span
for each \(w>0\). As they span the truncation over \(O\), every
positive-weight entry is zero there. Separation gives the same
conclusion in \(R\). In particular all upper triangular
off-diagonal entries vanish in the universal orthogonal lift.
This proves \(\mathcal O\subseteq\mathcal L\) at every depth.

The section of Theorem~\ref{thm:lowerfibers384} now represents
the whole functor by \(S(I(a))\). Decomposing the finite module
\(S(a)^\vee\), its normalized cocycle presentation in
Theorem~\ref{thm:evenroot358} identifies this functor with
\(\operatorname{Hom}_O(S(a)^\vee,I)\). The presentation depends
only on \(a\) and its local conditions, so it applies to this
choice of \(K\). Each cyclic generator maps to an arbitrary
element of \(I\) killed by its stated power of \(\ell\).
These equations give the displayed ring. Its universal
matrices are lower triangular with fixed diagonal, proving
the assertion about traces.
\end{proof}

The complementary order \(|S(\Psi^{-1})|\) occurs with its full
exponent in the limiting count in Theorem~\ref{thm:lowerfibers384}.
Proposition~\ref{prop:odddivisor384} controls the remaining
factors on the very same interval. An upper bound for the count
itself would therefore give information on the complementary
order. The rank-three comparison in
\srcloc{Development/AutomorphicArithmetic/AutomorphicVarieties/ModularityArgument/ThresholdArguments/DenseThresholdCases/FOneCharU}{33}
uses the finiteness of the full invariant algebra and the
automorphic description of its contributing points. The
orthogonal trace calculation above supplies neither an upper
bound for this count nor that description of the full algebra.


\subsection{Invariant summands and the largest trace minors}
\label{sec:flagtraceideals385}

Use the marked ring \(A\), its trace algebra \(T\) on \(G=G_K\)
and its trace augmentation \(\mathfrak a=\ker(T\to O)\) of
Section~\ref{sec:planedesc381}. In particular the coefficient
prime, the characters and the quadratic field are those of
Section~\ref{sec:lowerfibers384} on the original interval.
Write \(M(h)\) for the universal matrix of \(h\in G\).
Let \(J_\uparrow,J_\downarrow\) be the ideals of \(A\) cutting
out upper and lower triangular matrices, respectively.
Recall that \(J_1\) kills every adjacent entry and \(J_2\)
kills both outer entries. Put
\[
 J_{\rm f}=J_\uparrow\cap J_\downarrow\cap J_1,\qquad
 X=T\cap J_1,\qquad Y=T\cap J_2 .
\]
For a local Artinian coefficient algebra \(D\), an invariant
summand means a \(G\)-stable direct summand of \(D^3\), of rank
one or two as a \(D\)-module. The complementary summand need
not be stable.

\begin{proposition}[The trace image of families with an invariant summand]
\label{prop:flagimage385}
Every invariant summand of a marked lift over \(D\) is a coordinate
summand in the marked basis. A lift admits such a summand if and
only if its map from \(A\) factors through one of
\(A/J_\uparrow,A/J_\downarrow,A/J_1\).
The common kernel of the traces on all these families is \(X\).
Thus their scheme-theoretic image in the formal trace space has
coordinate ring \(T/X\).
\end{proposition}
\begin{proof}
The three coordinate projectors belong to the algebra spanned
by the matrices on \(G\). Indeed the difference of the matrices
at \(g_*\) and \(cg_*^{-1}c\) is diagonal with entries \(z,0,-z\)
whose differences are units. Interpolation gives the projectors,
as in Section~\ref{sec:planedesc381}. A stable direct summand
\(P\) is consequently the sum of three ideals in the three
coordinate lines. Its quotient \(D^3/P\) is free. Each of the
three corresponding cyclic quotients is a direct summand of
this free module and hence projective. A quotient \(D/I\)
which is projective splits off \(D\), so \(I\) is generated by
an idempotent. A local ring has just the idempotents zero and
one. This proves the assertion about coordinates.

If the subset of indices \(S\) spans an invariant summand,
then the complement of \(\{4-i:i\in S\}\) does as well.
Indeed the marked polarization on \(G\) is
\(M(ch^{-1}c)=J_3M(h)^{\mathsf T}J_3\). It reflects each
vanishing entry through the antidiagonal. Thus stability of
the first line forces stability of the first two coordinates.
Together these conditions give upper triangular matrices.
Stability of the last line gives lower triangular matrices.
Stability of the middle line gives stability of the outer
plane, which kills precisely the four adjacent entries.
Starting with a two-element subset gives one of the same
three cases by reflection. Each stated quotient visibly
preserves at least one of these summands.

In either triangular quotient every trace is the split trace.
For completeness, the ordered ordinary inertia characters give
the prescribed diagonal characters on inertia above \(\ell\).
One can see this without choosing an ordinary basis compatible
with the given triangularization. For a separating inertia
element \(g\), each diagonal entry at \(g\) is its prescribed
root by reduction and the pairwise unit differences. For any
inertia element \(h\), compute the trace at \(g^j h\) for
\(j=0,1,2\) in both triangular bases. The Vandermonde
matrix in the three roots at \(g\) is invertible, giving the
three diagonal values at \(h\). Away from \(\ell\) the ratios
to the split characters are trivial on inertia by the
prime-to-\(\ell\) inertia quotient. The class number of \(K\)
is prime to \(\ell\), so the three ratios are trivial globally.
Apply this argument in every Artinian truncation of each
triangular quotient to obtain the trace identity in the
complete rings. Hence
\[
 T\cap J_\uparrow=T\cap J_\downarrow=\mathfrak a .
\]
The split point belongs also to the outer-block quotient
\(A/J_1\), so \(X\subseteq\mathfrak a\).
It follows that \(T\cap J_{\rm f}=X\). Taking the intersection
of the three kernels computes the image ideal, including for
the disjoint union of families with their choice of summand.
This use of the intersection imposes no assertion that every
Artinian point of a scheme-theoretic union lifts to one of
its components.
\end{proof}

\begin{theorem}[A nilpotent extension of the flag trace ring]
\label{thm:flagthick385}
Let \(I_9\) be the ideal generated by the nine-by-nine determinants
of \((\operatorname{tr}M(h_i k_j))\), for \(h_i,k_j\in G\).
Then
\[
          I_9=X^4Y^2,\qquad X^8\subseteq I_9\subseteq X^4 .
\]
In particular \(T/I_9\) is finite over \(O\). Its surjection
onto \(T/X\) has kernel with eighth power zero. This kernel
is killed by a power of \(\ell\).

Put \(A'=A/J_2\). The ideal defining the image of the
invariant-summand families in \(A'\) is \(J_{\rm f}A'\),
and there is an exact sequence of \(O\)-modules
\[
 0\longrightarrow V_+\otimes_O V_-\longrightarrow A'
 \longrightarrow (O\oplus V_+)\times_O(O\oplus V_-)
 \longrightarrow0 .
\]
Its first term is precisely \(J_{\rm f}A'\). It is also the
image in \(A'\) of \(X/I_9\).
\end{theorem}
\begin{proof}
Let \(\mathcal E\) be the \(T\)-span of the matrices on \(G\)
and use the coordinate projectors \(e_i\) from the proposition.
The trace pairs \(e_i\mathcal E e_j\) only with
\(e_j\mathcal E e_i\). A matrix of pairings for one pair of
corners has entries \(a_u b_v\) over \(A\), since both
coordinate lines have rank one. Any two-by-two minor within
that block vanishes, also in \(T\subseteq A\). By multilinearity
the rows and columns of the nine-by-nine determinants can
range over the corners of \(\mathcal E\). A nonzero term uses
all nine different corners with the opposite column corners,
and its value is the product of their nine trace pairings,
up to sign. Conversely any such product is realized by taking
one row and the opposite column from each corner. The three
diagonal corners contribute unit ideals since \(e_i^2=e_i\)
has trace one. The four adjacent positions each contribute
\(X\), and the two outer positions contribute \(Y\), by
Lemma~\ref{lem:pairedcorners381}. This gives the formula for
\(I_9\) with the very projectors used for the invariant summands.
The same lemma gives \(X^2\subseteq Y\) by composing adjacent
corners. Squaring this inclusion and multiplying by \(X^4\)
gives \(X^8\subseteq I_9\). The other inclusion follows from
the formula.

The flag trace ring \(T/X\) is finite over \(O\) by the odd
two-dimensional descent of Theorem~\ref{thm:planedesc381}.
There the residual ratio on inertia is nontrivial, the character
lifting one is ordinary with unramified weight, and the
weight-augmentation quotient of Pan's ordinary pseudodeformation
ring is finite. The descent gives a surjection to this precise
trace ring. The terms in the filtration by powers of \(X/I_9\)
are finite modules over \(T/X\) because \(T\) is Noetherian.
The filtration has at most eight terms including its degree
zero quotient. This proves finiteness of \(T/I_9\) directly.
For the statement about torsion use
Theorem~\ref{thm:reduciblerigid379}, which gives a reduced
generic fiber for \(T/I_9\). Its nilpotent ideal \(X/I_9\)
becomes zero after inversion of \(\ell\). Finiteness over
\(O\) now gives a common power of \(\ell\) killing that ideal.

In \(A'\), Theorem~\ref{thm:tensortrace382}
gives a direct sum of \(O,V_+,V_-\)
and \(L=V_+\otimes_O V_-\). Both pure modules have square zero,
and the augmentation ideal annihilates \(L\).
The pure modules are generated over \(O\) by their upper or
lower adjacent matrix entries. Indeed products of two
augmentation elements have weight zero in \(A'\), and entries
generate the matrix algebra of coefficients. Imposing upper
triangularity in \(A'\) kills \(V_-\) and \(L\). Imposing
lower triangularity kills \(V_+\) and \(L\). Imposing the
outer block kills both modules, leaving just \(O\). The
intersection of these three ideals in \(A'\) is exactly \(L\).
The image of \(J_{\rm f}\) belongs to that intersection.
Conversely \(X\subset J_{\rm f}\) by the proposition, and
the image of \(X\) in \(T/Y\) is its entire augmentation \(L\)
because \(X+Y=\mathfrak a\). Thus the intersection is
\(J_{\rm f}A'\), also the image of \(X/I_9\). The direct sum
decomposition gives the asserted exact sequence.
\end{proof}

\begin{theorem}[A critical exponent for traces on invariant summands]
\label{thm:flagvaluation385}
Write \(\mathfrak b=\mathfrak a/X\) in \(T/X\). Let
\(e=e_{\ell,r,[\alpha]}\) be the least nonnegative integer such
that \(\ell^e S(\eta\Psi^{-2})=0\). For all \(m\ge1\),
\[
 \ell^e\mathfrak b^m\subseteq\mathfrak b^{m+1},\qquad
 \ell^{me}\mathfrak b\subseteq\mathfrak b^{m+1}.
\]
For the distinct triples on the interval and the integer \(H\)
of Proposition~\ref{prop:odddivisor384},
\[
 \sum_{\ell,r,[\alpha]} [O:\ZZ_\ell]\,
                e_{\ell,r,[\alpha]}\log\ell\ \le\ 16\log H .
\]
At \(\ell\nmid H\), every marked lift admitting an invariant
summand on \(G\) is upper or lower triangular. The trace ring
of these families is \(O\), and the augmentation ideal of
\(T/I_9\) has eighth power zero.
\end{theorem}
\begin{proof}
Use \(A_{\rm p}=A/J_1\), whose middle coordinate and outer plane
are stable. Its weight-zero subring is \(T_{\rm p}=T/X\) by
Lemma~\ref{lem:pairedcorners381}. Write \(P_-,P_+\) for its
submodules of weights \(-2,2\), respectively. They are finite
\(T_{\rm p}\)-modules generated by the outer matrix entries
of the indicated weights. Indeed a finite set of group
entries generates \(A_{\rm p}\) topologically. Among the
surviving entries the weights are \(0,2,-2\). Any monomial of
weight \(-2\) or \(2\) factors as a single outer entry of that
weight times a weight-zero monomial. The conclusion follows
on finite truncations by extracting the indicated weight. It
passes to the completion because the span of finitely many
entries over the compact \(T_{\rm p}\) is closed. The same
argument with the diagonal split values subtracted works for
the augmentation filtration.

In particular \(P_-/\mathfrak bP_-\) is the weight \(-2\)
augmentation cotangent module of \(A_{\rm p}\) over \(O\).
A monomial of weight \(-2\) of augmentation degree at least
two contains either a diagonal augmentation factor or an
opposite pair of outer entries. Each of these factors belongs
to \(\mathfrak b\), proving the assertion. The adjacent entries
removed in passing from \(A\) to \(A_{\rm p}\) have weights
\(1,-1\). Modulo the square of the augmentation ideal their
multipliers reduce to scalars in \(O\). Removing them therefore
leaves the weight \(-2\) cotangent module intact. By
Theorem~\ref{thm:cot116} this module is \(S(\eta\Psi^{-2})^\vee\).
To recall the sign and the local conditions, a square-zero
perturbation in \(E_{31}\) has coefficient character
\(\eta\Psi^{-2}\) since inverse transpose followed by reversal
negates \(E_{31}\). The mark at \(c\) fixes its coboundary
representative, the ordinary flag gives a locally split class
at \(\ell\), and the specified inertia kernels elsewhere give
the unramified condition by finite averaging. These arguments
apply with the character of the present quadratic field.
It follows that \(\ell^e P_-\subseteq\mathfrak bP_-\).

The ideal \(\mathfrak b\) is \(P_+P_-\), where a product of
modules means the span of all products. Indeed it is generated
by traces of opposite outer corners since \(\mathfrak a=X+Y\).
In the coordinate basis these traces are exactly the products
of the corresponding scalar entries. Conversely every such
product is invariant and has zero augmentation. Hence
\(\ell^e\mathfrak b\subseteq\mathfrak b^2\). Multiplication
by \(\mathfrak b^{m-1}\) gives the first inclusion. Applying
it at \(m\) successive filtration steps gives the second.

The exponent \(e\) is bounded by the \(O\)-length of the outer
strict group. In the determinant calculation of
Proposition~\ref{prop:odddivisor384} the map at weight \(2r\)
is \(\alpha\mapsto\eta\alpha^2\), of multiplicity at most
sixteen. Its terms alone give the stated sum. In particular
at \(\ell\nmid H\) the exponent is zero. The equality
\(\mathfrak b=\mathfrak b^2\) then forces \(\mathfrak b=0\)
by Nakayama's lemma. Also \(P_-=\mathfrak bP_-=0\), so the
entire outer-block family is upper triangular. The proposition
on invariant summands now gives the claimed classification
and \(\mathfrak a=X\). The eighth-power assertion follows
from Theorem~\ref{thm:flagthick385}.
\end{proof}

\begin{corollary}[Finite coefficient algebras detecting the missing lengths]
\label{cor:fourcoeff385}
Let \(C=S(\Psi^{-1})^\vee\), let \(O_j=O/\ell^jO\), and put
\[
         D_j=O_j[u,v]/(u^2,v^2).
\]
There are \(|S(\Psi^{-1})[\ell^j]|\) distinct trace maps from
\(T\) to \(D_j\) indexed by \(f\in\operatorname{Hom}_O(C,O_j)\).
They are realized by marked matrix lifts and factor on traces
through \(T/I_9\). For each \(f\) the image of \(X\) is
\(uv\, f(C)\). If \(f\ne0\), this trace admits no realization
by a marked lift with an invariant direct summand on \(G\), even
over a faithfully flat local Artinian extension of \(D_j\).

If \(C\) has a cyclic factor of length at least \(j\), this
construction detects an element of additive order \(\ell^j\)
in the mixed trace image. Among finite free \(O_j\)-algebras
testing \(A/J_2\) for such an element the minimum rank is four.
Every map from \(A/J_2\) into a local Artinian principal ideal
ring has split trace.
\end{corollary}
\begin{proof}
Choose a surjection \(p:V_+\to O\) splitting off the free
summand of Proposition~\ref{prop:upperorders382}. Use the
projection \(q:V_-\to C\) of Proposition~\ref{prop:ordinaryplane381}.
Map the two square-zero algebras \(O\oplus V_+\) and
\(O\oplus V_-\) to \(D_j\) by their usual scalar map and
by \(w\mapsto u\,p(w)\), \(w'\mapsto v\,f(q(w'))\) on the
augmentations. Theorem~\ref{thm:tensortrace382} identifies
their tensor product with \(A/J_2\), so these maps specialize
the universal marked matrices. The image of \(L\) is exactly
\(uv\, f(C)\), by surjectivity of \(p,q\). The image of \(X\)
is \(L\) before that specialization and \(I_9\) already maps
to zero. If two \(f\)'s differ on \(c_0\in C\), take a lift
of \(c_0\) under \(q\) and an element sent to one by \(p\).
Their tensor has a preimage in \(X\). Its two images in the
free \(O_j\)-summand \(O_j uv\) differ. This gives the count
by finite duality, identifying \(O_j\) with \((J/O)[\ell^j]\).
An invariant summand forces \(X\) to have zero image by
Proposition~\ref{prop:flagimage385}. The coordinate projector
argument works over every local Artinian coefficient extension,
and faithful flatness preserves any nonzero image. This proves
the assertion about realizations of these traces.

If the stated factor exists, choose \(f\) surjective to \(O_j\).
The image contains \(uv\), of order \(\ell^j\). Conversely let
a map from \(A/J_2\) to an \(O_j\)-algebra detect that order
in \(L\). As \(L\) is spanned by products from its two pure
modules, one such product \(ab\) must have order \(\ell^j\).
Otherwise \(\ell^{j-1}\) kills the whole image. Its factors
satisfy \(a^2=b^2=0\). A relation
\(s_0+s_1a+s_2b+s_3ab=0\) over \(O_j\), multiplied first
by \(ab\), gives \(s_0=0\) since the annihilator of \(ab\)
in \(O_j\) is zero. Multiplying next by \(b\) and by \(a\)
gives \(s_1=s_2=0\), and then \(s_3=0\). Thus these four
elements span a free submodule of rank four. Length over
\(O\) gives the lower bound on the rank of a finite free
target. The displayed \(D_j\) attains it.

Finally, in a local Artinian principal ideal ring whose maximal
ideal has nilpotence index \(n\), each nonzero element is a unit
times a power of a generator of that ideal. Every square-zero
element belongs to its \(\lceil n/2\rceil\)-th power. Any
product of two such elements is zero, including in the field
case. Thus \(L\) maps to zero and the trace on this quotient
is the split trace.
\end{proof}

The eighth-power inclusion is a bound on products in an ideal.
The estimate in Theorem~\ref{thm:flagvaluation385} controls
the coefficient exponents on the invariant-summand quotient.
On the kernel between the two trace conditions, the preceding
corollary exhibits the entire complementary group at every
coefficient depth. A uniform upper bound for the weighted
orders of the resulting complementary summands on the retained
triples is the large-prime estimate still required in
Section~\ref{sec:missing}.


\subsection{A cyclic complementary group in the fixed-weight deformation}
\label{sec:cyclicstrict397}

Return to the large-prime components on the intervals of
Section~\ref{sec:missing}. We use the marked quadratic
problem with \(K=\QQ(i)\) of Section~\ref{sec:original68}.
Write \(\eta\) for its character, \(O=\ZZ_\ell[\alpha]\),
\(k=O/\ell O\), \(J=\operatorname{Frac}(O)\) and
\(\Psi=\alpha\kappa_\ell^r\). Take \(2\le r<t(x)\),
\(\alpha(c)=(-1)^r\) and \(\ell>\exp(x^\varepsilon/12)\).
Write \(S(\chi)\) for the strict group with coefficients
\((J/O)(\chi)\) on \(\QQ\), as in Section~\ref{sec:lowerfibers384}.
This is the scalar comparison for the actual field \(E_x\)
in Theorem~\ref{thm:levelstrict345}.

There is a single nonzero integer \(H_0\) on the interval which
may be used to ensure
\[
      S(\eta\Psi^{-1})=S(\eta\Psi^{-2})=0
                           \qquad(\ell\nmid H_0).
\]
Indeed take \(F=E_*(i)\) on the interval, and the product of the
critical determinants \( |D_n(F)|\) of
Lemma~\ref{lem:int13} for \(2\le n\le 2t_*+1\), where
\(t_*=\lfloor V^{(3/2)401}\rfloor\). The characters
\(\eta\alpha^{-1}\) and \(\eta\alpha^{-2}\) have precisely
the critical signs for \(n=r+1\) and \(n=2r+1\).
Restriction to \(F\) is injective on the indicated classes
by corestriction, since \(\ell>[F:\QQ]\) for large \(V\).
The critical order estimate in the proof of
Proposition~\ref{prop:odddivisor384} bounds each relevant
length by the valuation of its determinant. The calculation
there uses an abelian field with the modulus enlarged by the
conductor of the fixed quadratic field. Here \(4N_*\) is such
a modulus. Its estimates on the degree, modulus and twist
range give \(\log H_0\ll_{\psi,\varepsilon} V^{1220}\) as there.

Assume in this subsection that \(\ell\nmid H_0\) and that
\(S(\Psi^{-1})\) is cyclic of positive length \(a\).
Let \(A\) be the complete fixed-weight marked matrix ring,
\(T\) its trace ring on \(G_K\), and \(\mathfrak a\) the
augmentation ideal in \(T\). The letter \(T\) in this
subsection denotes the ring. We shall construct a quotient
of \(T\) whose fiber at the split trace retains \(O/\ell^a O\).

\begin{proposition}[A presentation of the trace ring for a cyclic strict group]
\label{prop:cyclictrace397}
There are integers \(p,q\ge1\), a power series ring
\[
 S=O[[Z_1,\ldots,Z_p,U_1,\ldots,U_q]],\qquad
 \mathfrak b=(Z_1,\ldots,Z_p,U_1,\ldots,U_q),
\]
two elements \(F_1,F_2\in S\), and an ideal
\(J_+\subset\mathfrak b\) generated by at most \(p+q-2\)
elements, such that
\[
 \begin{split}
       T&\simeq S/(J_++\mathfrak b(F_1,F_2)),\\
       D&=S/(J_+,F_1,F_2),\qquad
       D/\mathfrak bD\simeq O/\ell^aO .
 \end{split}
\]
The trace augmentation corresponds to \(Z_i=U_j=0\).
The quotient \(T\to D\) is an isomorphism wherever some
element of \(\mathfrak a\) is invertible.
\end{proposition}
\begin{proof}
Use the integer weights for the marking torus
\(\operatorname{diag}(s,1,s^{-1})\). By
Theorem~\ref{thm:cot116}, the negative cotangent modules at the
split augmentation, over \(O\), are the duals of
\(S(\Psi^{-1})\oplus S(\eta\Psi^{-1})\) in weight \(-1\)
and of \(S(\eta\Psi^{-2})\) in weight \(-2\).
The weight-zero module vanishes by
Proposition~\ref{prop:fixed116}. In the present case the
three modules are \(O/\ell^aO,0,0\), respectively, with the
nonzero one in weight \(-1\).

In the homogeneous presentation of Theorem~\ref{thm:pres70},
eliminate each variable which is a residual linear term of
a relation with unit coefficient. To justify elimination,
apply a linear change among variables of the same weight
over \(O\), and then replace the variable by the entire
homogeneous relation. This is an automorphism of the complete
power series ring because its linear change on the residue
cotangent space is invertible. Its inverse can also be
constructed successively modulo the powers of the maximal
ideal. All relations have zero value at the split augmentation.
Quotient by the new variable and continue. Every elimination
removes one variable and one relation of the same weight.
Their classes give a \(k\)-basis of the relative cotangent
space at the closed point. The difference bounds between
the two rows of Theorem~\ref{thm:pres70} still hold weight by
weight, since a redundant relation can only decrease their
number.

There is therefore a single variable \(y\) of weight \(-1\),
and all the other variables have positive weights. Write
\(x_1,\ldots,x_p\) for those of weight \(1\) and
\(w_1,\ldots,w_q\) for those of weight \(2\).
The same table gives at most two relations of weight \(-1\),
at most \(p-1\) of weight \(1\) and at most \(q-1\) of
weight \(2\), with no others. In particular \(p,q\ge1\).
Pad the two negative relations with zeros if necessary.

In the ambient power series ring set \(Z_i=x_i y\) and
\(U_j=w_jy^2\). An invariant monomial contains precisely as
many powers of \(y\) as the sum of its positive weights.
Thus invariant series form exactly \(S\), with unique
expansions. A monomial of degree \(d\) in \(Z,U\) has
ordinary degree between \(2d\) and \(3d\) after substitution.
Hence both directions respect completion. Every series of weight \(-s\) for \(s>0\)
is \(y^s\) times an invariant series. Write the negative
relations as \(yF_1,yF_2\). A positive relation \(g\) of
weight \(s\) gives the invariant \(y^s g\in\mathfrak b\).
Use these invariants to generate \(J_+\).

To contract the relation ideal to the invariants, take
weight zero in each expression of an element as a finite sum
of relations times power series. For \(g\) positive of
weight \(s\) the allowed multiplier is \(y^s S\).
For \(yF_i\) the allowed multipliers \(h\) have weight one.
The elements \(yh\) fill \(\mathfrak b\), since they contain
\(y x_i=Z_i\) and \(y(yw_j)=U_j\), and they have zero
constant term. This gives the claimed ideal of relations
in \(S\). To identify the invariants with the trace ring, use
the coordinate projectors. They belong to the span of the matrices
on \(G_K\), and invariant products of their entries factor
into traced cycles, as in Lemma~\ref{lem:pairedcorners381}.
This calculation uses the two marks and the distinct residual
characters at the selected inertia element, all present for
the chosen quadratic problem. Alternatively the passage to
the invariant part of the completed presentation is precisely
the weight-zero calculation of Lemma~\ref{lem:algebra116}.

The integral cotangent module in weight \(-1\) has generator
\(y\) and relations \(F_1(0)y,F_2(0)y\), by taking linear
terms at the augmentation over \(O\). It was identified above
with \(O/\ell^aO\). Consequently
\((F_1(0),F_2(0))=\ell^aO\). This proves the fiber formula.
Inverting any member of \(\mathfrak b\) makes the ideals
\(\mathfrak b(F_1,F_2)\) and \((F_1,F_2)\) equal, proving
the final assertion.
\end{proof}

\begin{theorem}[A residual trace curve or a determinant bounding the full length]
\label{thm:cyclicstrict397}
For the cyclic group and ring \(D\) of the proposition the
following alternatives are exhaustive.

If \(D/\ell D\) has positive dimension, there is a continuous
trace specialization \(T\to k'[[s]]\), with \(k'/k\) finite,
which is split at the closed point and nonconstant. Its geometric
generic semisimple representation on \(G_K\) is absolutely
irreducible.

If \(D/\ell D\) is zero-dimensional, \(D\) is finite free
over \(O\) and is a complete intersection. Every geometric
point of \(D[1/\ell]\) has absolutely irreducible semisimple
representation on \(G_K\). For every \(d\in\mathfrak a\)
which is nonzero at each of these points,
\[
 a\ \le\ \operatorname{length}_O(D/dD)
     \ =\sum_\rho m_\rho\,v_\ell(d(\rho)),\qquad v_\ell(\ell)=1 ,
\]
where \(\rho\) ranges over the geometric trace points and
\(m_\rho\) is the dimension of the corresponding local
algebra of \(D\otimes_O\overline J\) over \(\overline J\).
In particular this applies to a Frobenius trace minus its
split value whenever it separates the points from that value.
\end{theorem}
\begin{proof}
First consider the reducible locus in the trace space. Let
\(J_1\) in \(A\) be generated by the adjacent entries and
\(X=T\cap J_1\), as in Section~\ref{sec:flagtraceideals385}.
The weight \(-1\) tangent variable can be taken from the
adjacent matrix entries. Indeed these entries span the
cotangent space in this weight, since the matrix entries
generate \(A\) and the diagonal projectors belong to their
span on \(G_K\). Choosing a homogeneous adjacent entry
with nonzero residual linear term as \(y\) gives another
coordinate by the same change-of-variables argument.
Killing all adjacent entries therefore kills \(y\).
The resulting ring is generated by positive weights over
\(O\) and has invariant quotient exactly \(O\), using its
split augmentation. Thus \(X=\mathfrak a\).

The ideal \(I_9\) of the nine-rowed minors of the trace
pairing contains \(X^8\) by Theorem~\ref{thm:flagthick385}.
For clarity, the matrix projectors separate three diagonal
corners and the six off-diagonal corners. Pairing each
corner with its opposite gives one factor of \(X\) for
each adjacent corner and one factor of the outer pairing
ideal \(Y\) for each outer corner. The largest minors
generate \(X^4Y^2\), and multiplication through the middle
coordinate gives \(X^2\subset Y\). These identities use
only the paired corners of the marked matrices and hold
with the present \(K\). At a geometric field point with
reducible semisimple representation the trace pairing
has rank at most \(2^2+1^2=5\), by taking traces in the
simple constituent blocks. Here the semisimple representation
associated with the degree-three determinant exists over the
algebraically closed field, in residue characteristic greater
than three as in Theorem~\ref{thm:reducibleminors378}.
Thus any geometric point at which some element of
\(\mathfrak a\) is nonzero is absolutely irreducible.

If \(D/\ell D\) has positive dimension, choose a prime
giving a one-dimensional complete local domain quotient.
The normalization is finite, since this quotient of a
power series ring over a finite field is excellent. Completing
at a maximal ideal of the normalization gives a complete
equal-characteristic discrete valuation ring \(k'[[s]]\)
with finite residue extension. The map from \(D\) into it
has one-dimensional image. Since \(D/(\ell,\mathfrak b)=k\),
some element of \(\mathfrak b\) has nonzero image in the
maximal ideal. Some trace must then be nonconstant, since
the ring is topologically generated over \(O\) by traces.
The closed and generic traces have the properties asserted.

In the other case put \(n=p+q\). The \(n\)-variable ring \(S\)
is regular local of dimension \(n+1\). The at most \(n\)
relations defining \(D\), together with \(\ell\), generate a
maximal-ideal-primary ideal, because \(D/\ell D\) is Artinian.
The fiber calculation gives \(D/\ell D\ne0\).
The height theorem forces exactly \(n\) relations in such
a list. These relations followed by \(\ell\) form a system of
parameters in \(S\), hence a regular sequence since \(S\) is
Cohen--Macaulay. In particular \(\ell\) is regular on \(D\).
A power of the maximal ideal of \(D\) is contained in \(\ell D\).
Completeness and successive lifting of a \(k\)-basis of
\(D/\ell D\) show that \(D\) is finitely generated over \(O\).
Being torsion-free over a discrete valuation ring, it is free.
The quotient \(D/\mathfrak bD=O/\ell^aO\) disappears after
inverting \(\ell\). Hence every point in characteristic zero
has some nonzero element of \(\mathfrak a\), proving the
irreducibility assertion there.

Finally \(D/dD\) surjects onto \(D/\mathfrak bD\). Multiplication
by \(d\) on \(D\) has nonzero determinant by the hypothesis
on its values. The valuation of this determinant is the length
of the cokernel by diagonalization over \(O\). After extension
to \(\overline J\), each local factor has a filtration with
\(m_\rho\) copies of its residue field. Multiplication acts on
each quotient by \(d(\rho)\). Multiplication of these determinants
proves the equality and the required inequality.
\end{proof}

The contribution of this residue orbit to the large-prime
sum is \([O:\ZZ_\ell]a\log\ell\), with the full exponent \(a\). To obtain the uniform sum in
Section~\ref{sec:missing} by this argument one would still
need to handle the residual curves, control the specializations
and their multiplicities by the common Hecke data, and treat
the complementary groups having more than one generator.

\section{Strict orders and descent to auxiliary fields}\label{nav:chapter-descent-strict-orders}
\subsection{The full strict order in the fiber of a trace quotient}
\label{sec:negativecovariants398}

Keep the marked problem over \(K=\QQ(i)\), the field \(E_x\),
\(2\le r<t(x)\), and the large prime \(\ell\) of
Section~\ref{sec:cyclicstrict397}. In this subsection the number of
generators of the complementary group is arbitrary. Let \(A_j\) be
the weight-\(j\) part of the matrix deformation algebra. Products
of homogeneous elements can be taken in \(A=\bigoplus_j A_j\),
as in Lemma~\ref{lem:algebra116}. Its invariant ring \(T=A_0\)
is the trace algebra on \(G_K\), by the projector calculation of
Lemma~\ref{lem:pairedcorners381}. Write \(\mathfrak n\) for the
augmentation ideal of \(A\) and \(\mathfrak a=T\cap\mathfrak n\).
Set \(M=A_{-1}\) and \(C_j=(\mathfrak n/\mathfrak n^2)_j\).
Each \(A_j\) is finite over the complete Noetherian ring \(T\).
We write \(N^\vee=\operatorname{Hom}_O(N,J/O)\) for a finite
\(O\)-module. All lengths below are over \(O\).

\begin{proposition}[The fiber of the negative weight module]\leavevmode
\label{prop:negfiber398}
At every such prime there is an exact sequence of \(O\)-modules
\[
 C_{-2}\otimes_O C_1\ \longrightarrow\ M/\mathfrak aM
     \ \longrightarrow\ C_{-1}\ \longrightarrow 0 .
\]
The first map is multiplication of lifts. In particular the middle
module has finite length. If \(\ell\nmid H_0\), where \(H_0\) is
the common integer of Section~\ref{sec:cyclicstrict397}, then
\[
              M/\mathfrak aM\simeq S(\Psi^{-1})^\vee .
\]
\end{proposition}
\begin{proof}
Use split-augmented variables of weights \(0,\pm1,\pm2\) in the
presentation of Theorem~\ref{thm:pres70}. A monomial of weight
\(-1\) which contains a variable of weight zero, an opposite pair,
or a triple of weights \((1,1,-2)\) or \((-1,-1,2)\) belongs
to \(\mathfrak aM\). Its indicated subproduct is a nonconstant
invariant, and the remaining factors still have weight \(-1\).
Without such a subproduct, the only possibilities are a single
variable of weight \(-1\) and a pair of weights \((-2,1)\).
Indeed after removing opposite pairs, mixed signs require distinct
absolute weights. If the positive weight is two, reaching total
weight minus one requires at least two negative units as soon as
a positive factor occurs. If the positive weight is one, two
negative factors of weight minus two would require at least three
positive units. This proves the assertion about monomials. It
applies to series by choosing a subproduct in each discarded
monomial and collecting terms by the finitely many possible
subproducts. The remaining series in each group has weight minus
one and converges in the completed presentation.

The cotangent projection kills \(\mathfrak aM\). Its kernel in the
fiber consists of classes represented by terms of ordinary degree
at least two, since the cotangent of the quotient is obtained
by taking the linear terms of the defining equations. Equivalently
one can lift a sum of products from \((\mathfrak n^2)_{-1}\)
to products of augmented series. Only the pairs \((-2,1)\)
can survive. Changing either factor by an element of the square
of the augmentation introduces terms of degree at least three,
all of which vanish in the fiber by the monomial calculation.
Multiplication therefore gives the asserted map on cotangents
and fills the kernel. Theorem~\ref{thm:cot116} identifies
\(C_{-1}\) with the dual of
\(S(\Psi^{-1})\oplus S(\eta\Psi^{-1})\) and \(C_{-2}\) with the
dual of \(S(\eta\Psi^{-2})\). Both have finite length and \(C_1\)
is finitely generated over \(O\). At \(\ell\nmid H_0\) the two
strict groups involving \(\eta\) vanish by the critical estimate
in Section~\ref{sec:cyclicstrict397}. This gives the last claim.
\end{proof}

\begin{theorem}[Trace quotients detecting every elementary divisor]\leavevmode
\label{thm:negfit398}
Assume \(\ell\nmid H_0\). Write \(C=S(\Psi^{-1})^\vee\) and
list the exponents of its elementary divisors in increasing
order as \(e_1\le\cdots\le e_s\), each positive. For \(j\ge0\) let
\(D_j=T/\operatorname{Fitt}_j^T(M)\). Then
\[
  D_j/\mathfrak aD_j \simeq O/\ell^{b_j}O,\qquad
  b_j=\sum_{i=1}^{\max(s-j,0)}e_i .
\]
At every prime of \(T\) outside \(V(\mathfrak a)\), the natural
map to \(D_0\) is an isomorphism after localization. For
\(j\ge1\) the ring \(D_j\) has finite length over \(O\).
\end{theorem}
\begin{proof}
In the minimal homogeneous presentation, the only possible
negative variables are \(y_1,\ldots,y_s\) of weight minus one.
Indeed \(C_{-2}=C_0=0\). Eliminate each residual linear relation
with a unit coefficient by taking that relation as a new variable
of its weight, as in Proposition~\ref{prop:cyclictrace397}.
The number of remaining weight-minus-one variables is the residue
dimension of \(C_{-1}\). Write \(x_i\) and \(w_h\) for the variables
of weights one and two. By pairing every \(x_i\) with one \(y_a\)
and every \(w_h\) with two of them, every invariant monomial is a
product of
\[
                 x_i y_a,\qquad w_h y_a y_b .
\]
These elements generate \(\mathfrak a\) as an ideal. In fact the
products have degrees two or three, so that grouping invariant
series by a first factor gives series in the ring of invariants.
Likewise every monomial of weight minus one has a factor \(y_a\)
with invariant complement. The \(y_a\) thus generate \(M\) over
\(T\), and these statements descend through the homogeneous relations.

Each displayed invariant generator has the form \(g=f y_a\)
with \(f\in A_1\). Inverting \(g\) in the graded algebra makes
\(y_a\) invertible, with inverse \(f/g\). Multiplication by
\(y_a\) identifies \(T[1/g]\) with \(M[1/g]\), with inverse
multiplication by \(f/g\). These open sets cover the complement
of \(V(\mathfrak a)\). There the zeroth Fitting ideal is zero and
the higher ones are the unit ideal, as is seen from a free
presentation on one generator.

For a finite presentation on \(n\) generators the \(j\)th Fitting
ideal is generated by the minors of size \(n-j\) of the relation
matrix, with the unit ideal for size zero or less. Base change of
the matrix presents the base change of the module, and taking its
minors commutes with base change. Apply this to \(T\to O\).
The proposition gives \(C\) as the resulting module. Its diagonal
presentation over \(O\) has entries \(\ell^{e_i}\), so the ideal
of the minors is generated by \(\ell^{b_j}\). This also covers
the zero module and empty sums.

For \(j\ge1\) the preceding localization shows that every prime
of \(D_j\) contains the image of \(\mathfrak a\). By the fiber
formula the only possible such prime is the closed point. The
Noetherian local quotient is therefore zero or Artinian. In the
latter case its simple factors are copies of \(k=O/\ell O\),
proving the finite-length assertion.
\end{proof}

\begin{corollary}[The exterior powers on the split trace]\leavevmode
\label{cor:negwedge398}
Under the hypotheses of Theorem~\ref{thm:negfit398}, for every
\(h\ge2\) reduction at the split trace gives an isomorphism
\[
             \bigwedge_T^h M\ \simeq\ \bigwedge_O^h C .
\]
For \(2\le h\le s\) its length is
\(\sum_{i=1}^{s-h+1}\binom{s-i}{h-1}e_i\). It is zero for
\(h>s\).
\end{corollary}
\begin{proof}
Use a generator \(g=fy\) of \(\mathfrak a\) from the preceding
proof, with \(f\in A_1\) and \(y\in M\). For \(m,n\in M\),
the elements \(fm\) and \(fn\) belong to \(T\). Commutativity of
the matrix coefficient algebra gives
\[
 \begin{aligned}
  g(m\wedge n)&=(fm)y\wedge n
     =y\wedge (fn)m\\
     &=-(fn)m\wedge y=-g(m\wedge n).
 \end{aligned}
\]
Here \((fm)n=(fn)m\) and \((fn)y=gn\). Since two is a unit,
\(g\) annihilates \(\bigwedge_T^2 M\), and hence also every
higher exterior power. Tensoring such a power with \(T/\mathfrak a\)
is thus an isomorphism. Exterior powers commute with this base
change by their alternating multilinear presentation. The
proposition identifies the result with \(\bigwedge_O^h C\).
In a direct sum of cyclic modules each increasing \(h\)-tuple
of summands contributes their tensor product, a cyclic module of
length the smallest exponent in the tuple. Exactly
\(\binom{s-i}{h-1}\) tuples have smallest index \(i\). Summing
proves the formula.
\end{proof}

In particular \(b_0=\operatorname{length}_O S(\Psi^{-1})\)
records the exponent of the entire group order in residue degree
one. In residue degree \(f\) its contribution to the logarithmic
sum is \(f b_0\log\ell\). The fiber formula gives this length
even when several generators occur. Bounding it on the interval
still requires an upper estimate for the integral trace quotient.


\subsection{A Frobenius determinant for a strict group of arbitrary rank}
\label{sec:horizontalstrict399}

The cyclic case used a finite flat complete intersection. For a
strict group with several generators the quotient in
Theorem~\ref{thm:negfit398} can have integral torsion. We obtain a
bound which uses the multiplicities in its generic fiber. We first
give the modification of the intersection calculation that accounts
for those multiplicities.

\begin{proposition}[The intersection inequality with local multiplicities]\leavevmode
\label{prop:localmult399}
Use the augmented presentation and rational formal chart of
Section~\ref{sec:lower36}, with its variable and equation counts.
Write \(B\) for the graded solution algebra, \(T=B_0\),
and \(k_0=\overline J\). Suppose \(T\) is finite over \(O\), the
negative equations span the two rational negative cotangents,
and no fixed point other than the distinguished point has both
a purely positive and a purely negative point over its invariant
character. Let \(\Sigma\) be the geometric characters of
\(T\otimes_O k_0\) nonzero on the ideal \(B_{-2}B_2\). For each
\(\sigma\in\Sigma\) put
\[
 m_\sigma=\max_{j=0,1}\dim_{k_0}
                  (B_j\otimes_O k_0)_\sigma ,
\]
where the subscript means localization at the character of the
invariant algebra. Let \(d\in T\) vanish at the distinguished
point and be nonzero at every \(\sigma\in\Sigma\). For the quotient
\(C_F\) of the integral negative cotangent lattice in
Theorem~\ref{thm:full36}, one has
\[
       \operatorname{length}_O C_F
          \le 2\sum_{\sigma\in\Sigma}m_\sigma v_\ell(d(\sigma)).
\]
Here \(O\) is unramified and \(v_\ell(\ell)=1\).
\end{proposition}
\begin{proof}
Lift \(d\) to the ambient invariant ring. In the three-color
algebra of Section~\ref{sec:upper35}, Lemma~\ref{lem:support35}
puts the high multidegrees on a finite \(O\)-module with generic
support in \(\Sigma\). Its proof uses finiteness, the condition
on fixed points and rational cotangent spanning. Indeed a relevant
prime over a character which kills the product of the two even
weights would give a purely positive and a purely negative point
over one fixed point. At the distinguished point rational
cotangent spanning excludes the purely negative point. None of
these arguments use a bound on a tangent dimension at a mixed
point.

For \(\sigma\in\Sigma\), some product from weights \(-2,2\)
is a unit after localization. Both factors are homogeneous units
in the localized graded algebra. This holds also in the ambient
algebra localized at the inverse image of \(\sigma\). In the
solution algebra multiplication by the weight-two unit identifies
any two pieces whose weights differ by two. Each piece therefore
has dimension at most \(m_\sigma\). These dimensions are finite
because the individual weight modules are finite over \(T\).

Write \(R,R'\) for the numbers denoted by \(r,r'\) in
Section~\ref{sec:upper35}. In its localized calculation of the
high pieces, one squares \(R+1+R'\) equations of odd weight.
If \(f\) has weight \(b\) and \(S\) is the algebra with the other
current equations imposed, multiplication by \(f\) gives the
exact sequence in each weight
\[
 (S/(f))_{j-b}\longrightarrow (S/(f^2))_j
                      \longrightarrow(S/(f))_j\longrightarrow0 .
\]
Consequently each squaring at most doubles the maximum dimension
over all weights. The unit of weight two makes each of the three
localized generating modules in Section~\ref{sec:upper35} cyclic.
As in the proof of Theorem~\ref{thm:upper35}, a high piece is
therefore a quotient of the weight-zero module with the squared
equations imposed. Its dimension at \(\sigma\) is at most
\(2^{R+1+R'}m_\sigma\). Multiplication by \(d\) on each residue
factor has eigenvalue \(d(\sigma)\). On the integral torsion the
kernel and cokernel have equal lengths. Thus the upper bound for
the difference of lengths in that theorem becomes
\(2^{R+1+R'}\sum_\sigma m_\sigma v_\ell(d(\sigma))\).

For the lower bound use the same formal chart and Koszul complex
as in Theorem~\ref{thm:full36}. Its Smith basis gives factors
\(\ell^{2c_i}\) from the squared negative equations, with
\(\sum c_i=\operatorname{length}_O C_F\). The finite difference of
the two monomial counts in that proof is \(2^{R-1+R'}\). Removing
the scalar factors in the Koszul complex therefore bounds its
Euler characteristic below by \(2^{R+R'}\sum c_i\). The lower
calculation takes place in the chart at the distinguished point.
Its use of the other characters is only through admissibility
and the upper estimate, both just proved. Dividing the new upper
bound by \(2^{R+R'}\) proves the inequality.
\end{proof}

\begin{theorem}[A residual curve or a horizontal determinant bound]\leavevmode
\label{thm:horizontalstrict399}
Use the original period field \(E_x\) and the marked quadratic
problem over \(K=\QQ(i)\). Let \(1\le r<t(x)\), \(\ell>P(x)\),
\(\ell\nmid H_0\), and suppose \(C=S(\Psi^{-1})^\vee\ne0\).
Let \(T\) be the full invariant trace algebra, \(M=B_{-1}\) its
negative weight module, and \(D=T/\operatorname{Fitt}_0^T(M)\).
On sufficiently distant intervals precisely one of the following
cases occurs.

If \(T/\ell T\) has positive dimension, there is a continuous
nonconstant trace \(T\to k'[[z]]\), split at the closed point,
with \(k'/k\) finite. Its generic representation on \(G_K\)
remains absolutely irreducible over every finite extension of \(K\).

If \(T/\ell T\) has dimension zero, set
\(H=D/D[\ell^\infty]\). Then \(H\) is finite free and nonzero
over \(O\). The geometric generic points of \(H\) are exactly
the nonsplit trace points of \(T\), and all are absolutely
irreducible. For every \(d\in\mathfrak a\) nonzero at these points,
\[
 \operatorname{length}_O C
       \ \le\ 2v_\ell\bigl(\det(d\mid H)\bigr)
       \ =\ 2\sum_{\sigma} n_\sigma v_\ell(d(\sigma)),
 \qquad n_\sigma=\dim_{\overline J}(H\otimes_O\overline J)_\sigma .
\]
Here \(d\) acts by multiplication. The determinant counts the
lengths of the local algebras even at nonreduced points.
\end{theorem}
\begin{proof}
The negative variables and invariant generators in the proof of
Theorem~\ref{thm:negfit398} apply throughout this range. For the
first twist the two extra strict groups still have the critical
signs, now for determinant indices two and three in \(H_0\).
The integral cotangent formula is Theorem~\ref{thm:cot116} also
at this twist. The exceptional character in the first-twist
equation count has zero complementary group at every sufficiently
large prime, by Corollary~\ref{cor:present116}. Thus a contributing
first twist has the same required chart and counts. In particular
\(D/\mathfrak aD=O/\ell^{\operatorname{length}_O C}O\), and
\(T\to D\) is an isomorphism locally off \(V(\mathfrak a)\).

To check irreducibility, choose the minimal variables of weight
minus one from the lower adjacent entries. This is possible since
the matrix entries generate the ring and the classes of the
entries of this weight span that residual cotangent. Removing
all adjacent entries now leaves only positive variables and the
split augmentation, so that the invariant quotient is \(O\).
In the notation of Section~\ref{sec:flagtraceideals385}, this
says \(X=\mathfrak a\). The paired-corner computation in the proof
of Theorem~\ref{thm:cyclicstrict397} gives \(X^8\subseteq I_9\)
for the trace pairing. Its rank is at most five on a reducible
semisimple representation and nine on an absolutely irreducible
one. Thus every nonsplit geometric trace is absolutely irreducible.

If the special fiber has positive dimension, take a one-dimensional
complete local domain quotient of that fiber. Its normalization
in its fraction field is finite since complete Noetherian local
rings are excellent. Completing at a maximal ideal of the
normalization gives \(k'[[z]]\) with finite residue field. The
trace ring has one-dimensional image, so some trace series is
nonconstant. The closed point is the split residual trace by
the defining augmentation modulo \(\ell\). The asserted strong
irreducibility follows from Theorem~\ref{thm:strong118} for
this same marked quadratic problem.

Otherwise the special fiber is Artinian. Lift a finite residue
basis and express every element by successive approximation
modulo \(\ell^j\) in their \(O\)-span. The coefficients converge
in \(O\), and the errors tend to zero in the maximal ideal topology
of \(T\). This proves that \(T\), hence \(D\), is finite over
\(O\). Removing its \(O\)-torsion gives a free module with the
same generic fiber. The split fiber formula excludes the split
generic point from \(D\), and its local rings at all other
generic points coincide with those of \(T\).

At any such nonsplit character a generator \(fy\) of the invariant
augmentation is nonzero, with \(y\) a negative variable of weight
minus one and \(f\) of weight one as in Theorem~\ref{thm:negfit398}.
Thus \(y\) becomes a unit after localization and all the local
weight pieces are isomorphic to the degree-zero piece. Its
dimension is \(n_\sigma\). Conversely the split evaluation kills
all products from weights \(-2,2\). Hence the characters in the
preceding proposition are exactly these nonsplit points, with
\(m_\sigma=n_\sigma\). The fixed locus is just the split point by
Proposition~\ref{prop:fixed116}, and Corollary~\ref{cor:present116}
provides the chart, spanning and integral cotangent length there.
The proposition gives the claimed inequality. The equality follows
by taking the determinant on the local factors over \(\overline J\),
each filtered by copies of its residue field. If \(H=0\) the same
proposition with no contributing character gives
\(\operatorname{length}_O C=0\), contrary to the hypothesis.
\end{proof}

\begin{corollary}[Frobenius tests attaining the minimum valuation]\leavevmode
\label{cor:frobeniusstrict399}
In the finite case of Theorem~\ref{thm:horizontalstrict399} let
\(t\) be the trace on \(G_K\) and
\(t_0=\Psi+1+\Psi^{-1}\) its split value. Put
\[
     N(h)=\det\bigl(t(h)-t_0(h)\mid H\bigr)\in O
                               \qquad(h\in G_K).
\]
The ideal generated by all the \(N(h)\) is nonzero. Its
valuation is the minimum of \(v_\ell N(\operatorname{Frob}_w)\)
as \(w\) runs over good places of \(K\) of degree one outside
any fixed finite set, with \(v_\ell(0)=+\infty\). A set of
places of positive Dirichlet density attains this minimum. Twice
the minimum bounds \(\operatorname{length}_O C\) from above.
\end{corollary}
\begin{proof}
Over a finite extension of \(J\) realize each of the finitely many
irreducible representations \(\rho_\sigma\) at the generic points.
Indeed the invariant evaluation is integral, since \(T\) is finite
over \(O\), and is local into the integer ring by uniqueness of
the maximal ideal of \(T\). It is continuous since a power of
that maximal ideal lies in \(\ell T\). Lemma~\ref{lem:scale118}
gives the continuous matrix realization after a finite extension.
Take the joint Zariski closure of \((\Psi,(\rho_\sigma)_\sigma)\)
over \(\overline J\), and denote its identity component by \(G^0\).
The inverse image \(U\) of \(G^0\) in \(G_K\) is open and normal.
Indeed the component group is finite and its inverse images are
closed by continuity. The image of \(U\) is Zariski dense in
\(G^0\), since this component is open in the Zariski closure.

For each \(\sigma\) the regular function
\(\operatorname{tr}\rho_\sigma-1-\Psi-\Psi^{-1}\) on this joint
group is nonzero on \(G^0\). Otherwise its vanishing on \(U\)
would identify the semisimplification of \(\rho_\sigma|_U\)
with \(1\oplus\Psi|_U\oplus\Psi^{-1}|_U\) by the character
criterion in characteristic zero. The restriction itself is
semisimple. To see this, take an irreducible \(U\)-submodule and
sum its \(G_K\)-translates. These are simple \(U\)-submodules by
normality, and they span the whole space by irreducibility on
\(G_K\). The sum of simple submodules is semisimple. The three
characters are distinct on \(U\) since the cyclotomic character
has infinite image on every open subgroup and \(r>0\). Each
extends to \(G_K\), so each of the three one-dimensional
isotypic subspaces would be \(G_K\)-stable. This contradicts
absolute irreducibility. Since the connected algebraic group
\(G^0\) is irreducible, the product of these finitely many
nonzero functions is nonzero somewhere on its dense subgroup.
The determinant formula in the theorem now gives some \(h\)
with \(N(h)\ne0\).

Its nonzero values have nonnegative integral valuations, so the
least valuation \(b\) is attained. The reduction of the continuous
class function \(N\) modulo \(\ell^{b+1}\) factors through a
finite quotient of \(G_K\). Indeed its fibers are clopen and
compact. Finitely many cosets of open subgroups cover them, and
one may take an open normal subgroup refining these cosets. The
conjugacy classes of this quotient giving valuation \(b\) form a
nonempty set. The Chebotarev density theorem gives a positive
Dirichlet density of unramified places with their Frobenius in
this set. Places of degree greater than one in the quadratic
field have norm at least the square of their residual rational
prime, so their sum of inverse norms converges. They have
Dirichlet density zero. Removing them and the finitely many
excluded places preserves the positive density. Evaluating the
theorem at one of the remaining places proves the final assertion.
\end{proof}

This determinant is defined on the same scalar component as
the original cohomology group at \(x\). To sum its valuations
in Section~\ref{sec:missing} one still needs arithmetic control of
the nonsplit traces and their multiplicities, as well as control
of the residual curves. The existence of minimizing primes gives
no upper estimate for their determinant valuations as the fields
and twists vary.


\subsection{A uniform contribution to the large-prime strict length}
\label{sec:gaussianbudget401}

Use \(K=\QQ(i)\) and the original fields \(E_x\) on the intervals
of Section~\ref{sec:missing}. We now allow every \(\ell>P(x)\),
including divisors of the critical determinant, and every
\(1\le r<t(x)\). Keep \(\alpha(c)=(-1)^r\) and
\(\Psi=\alpha\kappa_\ell^r\). Let \(B=\bigoplus_j B_j\) be the
graded marked solution of Lemma~\ref{lem:algebra116}, with trace
ring \(T=B_0\), and let \(\mathfrak a\) be its trace augmentation
ideal. Write \(C_j\) for the integral cotangent in weight \(j\).
We bound the part of the total length accounted for by classical
systems with a fixed prime and a single integer throughout the
interval. This requires retaining the two-dimensional mixed
points at primes where the outer negative cotangent is nonzero.

\begin{lemma}[The even negative module at the split and mixed traces]\leavevmode
\label{lem:evenfiber401}
Put \(M_2=B_{-2}\) and \(D_2=T/\operatorname{Fitt}_0^T(M_2)\).
There is an exact sequence
\[
     \operatorname{Sym}^2_O C_{-1}\longrightarrow
           M_2/\mathfrak aM_2\longrightarrow C_{-2}\longrightarrow0.
\]
In particular \(D_2/\mathfrak a D_2\) has finite \(O\)-length.
If \(T\) is finite over \(O\), the algebra
\(H_2=D_2/D_2[\ell^\infty]\) is finite free over \(O\). Its
geometric generic local factors, with their full local rings,
are precisely the nonsplit factors of \(T\).
\end{lemma}
\begin{proof}
In an augmented homogeneous presentation a monomial of weight
minus two with no nonconstant invariant factor is either one
variable of weight minus two or two variables of weight minus
one. Indeed remove weight-zero variables and opposite pairs.
If both signs persist they have different absolute weights. A
positive two then requires at least two negative units. With
positive units and negative twos, any nonempty positive part
has even size and supplies two units to pair with a negative
two. In either case there is an invariant triple. Grouping series
by one such subproduct shows that all discarded terms belong
to \(\mathfrak aM_2\). The cotangent quotient removes the
quadratic terms that remain. A change of either factor in such
a quadratic term by an augmented product introduces terms of
degree at least three, which have just been removed. The first
map in the sequence is therefore induced by multiplication and
its image is exactly the kernel. Both negative cotangents have
finite length by Theorem~\ref{thm:cot116}. The fiber of the
zeroth Fitting ideal now has finite colength over \(O\) by
base change and diagonalization, including colength zero if
the module is zero.

A geometric generic character of \(T\) killing \(B_{-2}B_2\)
is the split one. To see this, for \(i>0\) products of opposite
weights \(-i,i\) have square in \((B_{-2}B_2)^i\). In fact in
even weight \(2i\) pair the odd factors of a monomial to make
weights zero or of absolute size two, then pair opposite even
factors. This gives \(B_{2i}=B_2^i\), with the same formula in
negative weights. It applies also to homogeneous series by the
finite generation in Lemma~\ref{lem:algebra116}. The character
therefore extends to \(B\) by sending every nonzero weight to
zero. Proposition~\ref{prop:fixed116} identifies the resulting
fixed point with the split one. Conversely the split augmentation
kills \(B_{-2}B_2\).

At any other character some product of weights \(-2,2\) is
invertible in the localized algebra. Its factor of weight minus
two is then a homogeneous unit and multiplication by it
identifies the local invariant ring with the local module
\(M_2\). Its zeroth Fitting ideal there is zero. Over the split
character the finite-length fiber found above disappears on
inverting \(\ell\). This proves the assertion about the geometric
local rings. The last quotient is finite torsion-free over the
discrete valuation ring and hence is free.
\end{proof}

\begin{theorem}[The mixed trace estimate at every large prime]\leavevmode
\label{thm:allhorizontal401}
Suppose \(S(\Psi^{-1})\ne0\) on a sufficiently distant interval.
If \(T\) is finite over \(O\), for every \(d\in\mathfrak a\)
nonzero at its nonsplit geometric generic characters one has
\[
 \operatorname{length}_O S(\Psi^{-1})+
 \operatorname{length}_O S(\eta\Psi^{-1})
 \ \le\ 2v_\ell\bigl(\det(d\mid H_2)\bigr)
 \ =\ 2\sum_{\sigma} n_\sigma v_\ell(d(\sigma)),
\]
where \(n_\sigma=\dim_{\overline J}(T\otimes_O\overline J)_\sigma\)
and the sum ranges over those characters. In particular \(H_2\ne0\).
If \(T\) is infinite over \(O\), its special fiber admits a
nonconstant trace curve strongly irreducible on \(G_K\).
\end{theorem}
\begin{proof}
The finite case of Proposition~\ref{prop:localmult399} applies by
Corollary~\ref{cor:present116} and the fixed-point calculation.
Its index set is all nonsplit characters by the lemma. We check
that its maximum over two weights is \(n_\sigma\) in every case.
At a character admitting a homogeneous unit of weight one or
minus one this follows by multiplication by the unit. For the
remaining case work over \(\overline J\), localize in the
invariant algebra at \(\sigma\), and choose a unit \(v\) of
weight two. The local invariant algebra \(T_\sigma\) is Artinian.
An element \(z\) of odd weight \(j\) has \(z^2/v^j\) in the
maximal ideal of \(T_\sigma\). Otherwise it would be a unit,
and \(zv^{(1-j)/2}\) would be a unit of weight one. Thus every
odd homogeneous element is nilpotent. The even part equals
\(T_\sigma[v,v^{-1}]\), by periodicity. Its intersection
with the nilradical has coefficients in the maximal ideal of
\(T_\sigma\), since distinct even degrees remain independent.
Equivalently, quotient by the nilpotent ideal generated by
this maximal ideal and a finite set of odd homogeneous algebra
generators. The reduced ring is \(\overline J[v,v^{-1}]\).
There are no polynomial relations on the Laurent generator over
the residue field, since its nonzero powers have distinct weights
and remain units in this nonzero graded quotient.

At any closed point of this curve all adjacent entries vanish
by their odd weights. The two outer directions both occur. In
fact elements of weights two and minus two with nonzero product
must evaluate nontrivially, and in the absence of adjacent
entries an entry of the corresponding outer weight is required
in each case. The marked projectors belong to the group span,
so the outer plane is irreducible and the middle line splits off.
These are precisely the two-dimensional mixed points treated by
Corollary~\ref{cor:corner120}. The tangent space of the full
solution at each of them has dimension one. Its local ring has
dimension one by the reduced Laurent curve and is Noetherian.
It is therefore regular, hence reduced. The localized graded
algebra is of finite type over \(\overline J\). An element
vanishing in the localizations at every maximal ideal is zero,
so it is itself reduced. Its invariant part is now \(\overline J\)
and its odd part is zero. The maximum in the proposition is
again \(n_\sigma\).

The length on the left is the integral cotangent length in
Corollary~\ref{cor:present116}. The determinant equality follows
from the lemma by filtering the finite geometric local rings.
An empty sum would contradict the positive left side. Finally,
if \(T\) is infinite, its special fiber has positive dimension
by the complete-local finiteness argument of
Theorem~\ref{thm:horizontalstrict399}. Normalizing a complete
curve quotient and applying Theorem~\ref{thm:strong118} gives
the stated curve.
\end{proof}

Choose one rational prime \(v\ge7\) split in \(K\) and prime
to the fixed conductor and to \(6\). Fix \(w\mid v\). On all
sufficiently distant intervals, \(v\nmid N_*\) and \(v<P(x)\),
since the moving prime divisors of \(N_*\) tend to infinity.
Let \(N'=4N_*^2\). At a finite trace algebra call a nonsplit
geometric point listed if one of the following descriptions
holds for its representation on \(G_K\).

\begin{enumerate}
\item It has a split middle line and the irreducible outer plane
of Theorem~\ref{thm:form117}, or it is an odd orthogonal point
of Theorem~\ref{thm:oddtrace123}.
\item It is absolutely irreducible and its twist by
\(\kappa_\ell^{-1}\) is attached to a regular algebraic
conjugate self-dual cuspidal representation over \(K\) with the
level, weights and central condition of Section~\ref{nav:III-00}
at level \(N'\). Use the normalization there for its Frobenius
traces. The association here is with one global representation.
\end{enumerate}

For a geometric character \(\alpha\) of \(E_x\) put
\(\lambda_{\ell,r,\alpha}=\operatorname{length}_O S(\Psi^{-1})\),
using the minimal unramified coefficient ring for \(\alpha\).
In the following definitions put \(U_{\ell,r,\alpha}=0\) when
\(\lambda_{\ell,r,\alpha}=0\). At an infinite trace algebra put
\(U_{\ell,r,\alpha}=\lambda_{\ell,r,\alpha}\). At a finite
one put \(d_\alpha(v)=1+v^r\alpha(v)+v^{-r}\alpha(v)^{-1}\) and
\[
 U_{\ell,r,\alpha}=\min\!\left\{\lambda_{\ell,r,\alpha},\,
    2\sum_{\sigma\ {\rm unlisted}}n_\sigma
    v_\ell\bigl(t_\sigma(\operatorname{Frob}_w)-d_\alpha(v)\bigr)
                                     \right\}.
\]
Only nonsplit points are counted and \(v_\ell(0)=+\infty\).
The lengths refer throughout to the scalar strict groups over
\(\QQ\). At these primes Theorem~\ref{thm:levelstrict345}
identifies their sum over geometric characters with the full
complementary length for the original field \(E_x\).

\begin{theorem}[One interval bound for the lengths accounted for by classical points]\leavevmode
\label{thm:gaussianbudget401}
There is a nonzero positive integer \(J_V\) with
\(\log J_V\ll_{\psi,\varepsilon,v} V^C\) for a fixed \(C\), such
that for every \(x\in[V,V^{3/2}]\) on the prescribed intervals
\[
 0\le \sum_{\ell>P(x)}\log\ell\sum_{1\le r<t(x)}
       \sum_{\alpha(c)=(-1)^r}
       (\lambda_{\ell,r,\alpha}-U_{\ell,r,\alpha})
       \le 2\log J_V .
\]
The character sum is over all geometric characters of \(E_x\)
of the indicated parity. Equivalently each residue orbit is
counted with its residue degree. In particular, for the finite
components on which every nonsplit point is listed the bound
counts their entire complementary strict length.
\end{theorem}
\begin{proof}
At every listed point the local multiplicity is one. For the
middle-line and odd orthogonal cases this follows from
Corollary~\ref{cor:corner120} and Theorem~\ref{thm:oddtrace123}.
For the cuspidal case use Corollary~\ref{cor:tangent121} with the
CM field equal to \(K\). In each case the full solution tangent
space is one-dimensional, generated by scaling. To see the
multiplicity conclusion, if there is a degree-one unit the local
graded algebra is \(T_\sigma[u,u^{-1}]\) with \(u\) that unit.
Its tangent dimension at a closed point is one plus the tangent
dimension of the invariant factor. Thus \(T_\sigma\) is a field
by Nakayama's lemma. Without the odd unit this was proved in
Theorem~\ref{thm:allhorizontal401}.

For a finite component evaluate that theorem at the Frobenius
trace minus its split value. All its geometric values are
integral since the ring is finite over \(O\). At a listed point
the trace has complex absolute value at most three, while
\(|d_\alpha(v)|\ge v^r-1-v^{-r}>3\). Hence a zero difference
can occur only in the unlisted sum, in which case
\(\lambda-U=0\). When all differences are nonzero the theorem
and the definition of \(U\) give
\[
 0\le\lambda_{\ell,r,\alpha}-U_{\ell,r,\alpha}
   \le 2\sum_{\sigma\ {\rm listed}}
           v_\ell\bigl(t_\sigma(\operatorname{Frob}_w)-d_\alpha(v)\bigr).
\]
The infinite components contribute zero to the left by definition.

To construct the one integer, take all twists \(1\le r<t_*\)
and all complex characters of \(E_*\) of the indicated parity.
For each such pair use three labelled lists at \(v\). For the
rank-three cuspidal systems use the trace multiset of
Theorem~\ref{thm:integer} at level \(N'\), with the two signs
and all the indicated embeddings. For the planes include
\(1+\alpha(v)a_v(f)/v^r\) for all primitive forms \(f\) of
weight \(2r+1\), character \(\eta\alpha^{-2}\) and level dividing
\(4N_*\). For odd orthogonal points include
\(a_v(g)^2/(v^r\alpha(v))-1\) for primitive forms \(g\) of
weight \(r+1\), character \(\eta\alpha\) and level dividing
\(4N_*\). The two formulas and level bounds follow from
Theorems~\ref{thm:form117} and~\ref{thm:oddtrace123}.
The newform roots have the common purity absolute value in their
respective weights, and \(\eta(v)=1\), giving the absolute value
bound three also for the adjoint trace. Finite characters have
unit complex absolute values. The three lists cover the listed
points. One global label with a fixed sign determines at most
one trace character for the given \(\alpha\). For rank three,
equality at the good primes over \(K\) implies equality of
traces by Chebotarev and continuity, and these traces generate
\(T\). In the other two cases the normalized newform determines
the entire plane or adjoint representation. Distinct equal
numerical values at \(w\) may still occur on different labels.

Take \(J_V\) to be the absolute value of the product of
\(v^{u_r}(b-d_\alpha(v))\) over these twists, characters and
labelled entries \(b\), with \(u_r=72(r+2)+1\). All factors
are algebraic integers by Theorem~\ref{thm:traces} and the
integrality of newform coefficients. They are nonzero by the
archimedean bound. The list in the entire product is Galois
stable, including the pairing of a newform character with
\(\alpha\). The rank-three list contains the conjugates as in
Theorem~\ref{thm:integer}. Thus the product is a nonzero rational
integer before taking absolute value. The number of labelled
entries per pair is polynomial in \(N'(r+2)\), by the same
theorem and the dimension estimate for newforms in its proof.
Each factor has logarithmic absolute value \(O_v(r+1)\). Since
\(N_*,[E_*:\QQ],t_*\) are polynomial in \(V\) by
Theorem~\ref{thm:conductor232}, this proves the asserted size.

At a given \(\ell>P(x)\), choose one embedding for all the
labels. The rational prime \(v\) is an \(\ell\)-adic unit.
Inflating characters from \(E_x\) to \(E_*\) is injective.
Every listed contribution for each \(\alpha,r\) therefore uses
a distinct allowed factor of the product, and all unused factors
have nonnegative valuation. Their sum of valuations is at most
\(v_\ell(J_V)\). Summing with \(\log\ell\) gives the claimed
bound. Conjugation over \(\QQ_\ell\) preserves valuations and
all three global lists, so summing the geometric characters is
equivalent to the residue-degree convention stated above.
\end{proof}

Thus the remaining large-prime contribution is measured by the
\(U_{\ell,r,\alpha}\) on those same fields. Bounding their sum
requires further information about the nonconstant residual
curves and the unlisted characteristic-zero traces. The
small-prime tails and the uniform comparison for the signed
periods in Section~\ref{sec:missing} are separate remaining
estimates.


\subsection{Solvable descent to the level of the common interval}
\label{sec:gaussiansolvable402}

Take an absolutely irreducible characteristic-zero trace point in
the marked problem of Section~\ref{sec:gaussianbudget401}. Denote
its representation on \(G_K\) by \(U\), where \(K=\QQ(i)\).
An integral trace point has a continuous integral matrix
realization after finite coefficient extension by
Lemma~\ref{lem:scale118}. This applies in particular to every
geometric generic point of a trace algebra finite over \(O\).
Fix a coefficient isomorphism \(\iota:\overline{\QQ}_\ell\simeq\CC\).
We use the rank-three normalization of
Theorem~\ref{thm:auxdescent361}. In it, \(U\otimes\kappa_\ell^{-1}\)
is the Galois representation attached to the automorphic system,
and the good Frobenius traces of \(U\) are its unitary traces,
with inverse Frobenius giving the opposite sign of the label.

\begin{theorem}[Descent through every solvable CM Galois extension]\leavevmode
\label{thm:gaussdesc402}
Let \(L/K\) be a finite solvable Galois extension, with \(L\) a
CM field. Suppose \((U\otimes\kappa_\ell^{-1})|_{G_L}\) is
attached to a regular algebraic conjugate self-dual cuspidal
representation on \(\mathrm{GL}_3(\mathbb A_L)\). Then
\(U\otimes\kappa_\ell^{-1}\) is attached to such a cuspidal
representation over \(K\). Its unitary central character is
trivial. The assertion includes extensions of degree divisible
by three.
\end{theorem}
\begin{proof}
The marked polarization gives \(U^c\simeq U^\vee\) by equality
of traces and absolute irreducibility. Lemma~\ref{lem:strong119}
gives \(\det U=1\) and absolute irreducibility after any finite
restriction inside \(G_K\). The unitary central character of the
cuspidal representation over \(L\) is therefore trivial. At
every good finite place it agrees with the determinant of the
unitarily normalized Galois representation. Equality there
determines a Hecke character by strong multiplicity one for
\(\mathrm{GL}_1\).

Take a subnormal series of \(\operatorname{Gal}(L/K)\) with
prime cyclic quotients. Its intermediate fields are CM. Indeed
complex conjugation on a CM field commutes with every field
automorphism. This follows upon composing any complex embedding
with that automorphism and using the defining conjugation on
complex embeddings. It therefore preserves each fixed field in
the series. The fixed subfield of conjugation is totally real,
and conjugation is nontrivial on each intermediate field since
it contains \(K\). It also commutes with the group of each
cyclic step.

Consider one step \(F/F_0\) of prime degree \(d\), with the
cuspidal system \(\Pi_F\) already attached to the restriction
of \(U\) in the indicated normalization. At good places its
parameters are invariant under \(\operatorname{Gal}(F/F_0)\)
because its Galois representation comes by restriction from
\(G_{F_0}\). Strong multiplicity one gives invariance of
\(\Pi_F\). By cyclic base change it has a cuspidal descent
\(\Theta\) over \(F_0\), and all its cuspidal descents are the
twists of \(\Theta\) by characters of the cyclic group, viewed
by class field theory as Hecke characters and composed with the
determinant. Base change here commutes with automorphisms,
contragredients and central characters. At the infinite places
the local fields are complex on both sides and the infinitesimal
parameters are inherited, so every descent is regular algebraic.
These are the cyclic base change and fiber statements of
Arthur and Clozel.\,\cite{VIII-VI-I-ac123}

We choose a conjugate self-dual descent. If \(d=2\), the central
character of \(\Theta\) base changes to one. The cyclic norm
theorem identifies it with a character of the cyclic quotient.
Twisting by a character \(\nu\) multiplies the central character
by \(\nu^3\). Cubing is bijective on the group of two characters,
so exactly one descent has trivial central character. Its
conjugate contragredient is again such a descent, hence is
isomorphic to it.

If \(d\) is odd, write
\((\Theta^c)^\vee\simeq\Theta\otimes\delta\) for a character
of the cyclic quotient, since both have base change \(\Pi_F\).
Conjugation fixes these quotient characters. Squaring is
bijective on them, so take \(\nu^2=\delta\). The conjugate
contragredient of \(\Theta\otimes\nu\) is
\(\Theta\otimes\delta\nu^{-1}\simeq\Theta\otimes\nu\).
This chooses a conjugate self-dual descent also when \(d=3\).

Let \(V\) be the normalized Galois representation of the chosen
descent, with the Tate factor restored. It exists by the
rank-three regular algebraic conjugate self-dual theorem and
satisfies \(V^c\simeq V^\vee\) in this normalization, by
compatibility at good places and the character criterion.
Its restriction to \(G_F\) agrees with \(U|_{G_F}\): the good
parameters agree by cyclic base change, and Chebotarev and
semisimplicity identify their semisimplifications. The latter
is irreducible of dimension three, so the restrictions are
indeed isomorphic.\,\cite{VIII-VI-V-shin,VIII-VI-I-ac123}

After identifying their restrictions, \(V\) and \(U\) on
\(G_{F_0}\) differ by a character \(\xi\) of the cyclic quotient.
For each group element the two matrices induce the same
conjugation on the irreducible restriction, so their ratio
is scalar by Schur's lemma. The scalars multiply and are one
on \(G_F\). There is no nontrivial self-twist of \(U|_{G_{F_0}}\)
by a character of this quotient, since its intertwiner would
commute with \(U|_{G_F}\) and hence be scalar. For \(d=2\) both
normalized determinants are one by the choice of central
character and good-place compatibility. Thus \(\xi^3=1\), which
forces \(\xi=1\). For odd \(d\) apply the conjugate-dual operation
to \(V\simeq U\otimes\xi\). Since \(\xi^c=\xi\), this gives a
self-twist of \(U\) by \(\xi^2\). It follows that \(\xi^2=1\),
again forcing \(\xi=1\). The descent therefore gives exactly
\(U\) in the normalized convention. Its unitary central character
is trivial by its determinant at the good places. Induction down
the cyclic tower proves the theorem.
\end{proof}

\begin{proposition}[The level and weights after descent]\leavevmode
\label{prop:gausslevel402}
On the polynomial-conductor intervals of
Theorem~\ref{thm:conductor232}, the cuspidal representation over
\(K\) in Theorem~\ref{thm:gaussdesc402} has fixed vectors at the
principal level \(U_K(N_*^2)\). At both embeddings its coefficient
weights satisfy the bound \(4(r+2)\) in Section~\ref{nav:III-00}.
Consequently it is on the list in Theorem~\ref{thm:gaussianbudget401}
for the entire interval, independently of the level, field degree
or ramification used for its potential automorphy.
\end{proposition}
\begin{proof}
At \(z\mid\ell\) the original marked representation is
crystalline with cyclotomic exponents \(r,0,-r\), including
\(r=1\) by Theorem~\ref{thm:crystalline121}. Its Weil--Deligne
parameter has trivial inertia and zero monodromy. Local--global
compatibility at \(\ell\) applies with equality up to Frobenius
semisimplification by the theorem of Barnet-Lamb, Gee, Geraghty
and Taylor, since the weight is Shin-regular in odd rank. The
local automorphic component is spherical. At the complex
embeddings regular algebraic compatibility gives, in the
unitary convention, infinitesimal exponents \(r,0,-r\).
Subtracting the half-sum of positive roots \((1,0,-1)\) in
dominant order gives coefficient weights of absolute value at
most \(r+1\), also after taking the dual. The trivial unitary
central character supplies the required central condition.
\,\cite{VIII-VI-lgp361,VIII-VI-V-shin}

At a finite place \(z\nmid\ell\) of \(K\) the marked inertia
kernel makes the action factor through the cyclic image of
\(\alpha\) on that inertia group. This image has order prime to
\(\ell\), since the conductor modulus \(N_*\) is less than
\(\ell\) throughout a sufficiently distant interval. Its
characters in the integral representation are therefore the
unique lifts of the three residual inertia characters, with
multiplicity. They are \(\alpha,1,\alpha^{-1}\) restricted to
this inertia, since the cyclotomic character is unramified
here. Finite inertia implies zero monodromy. The Artin conductor
exponent is the sum of the three character exponents, as is
seen by taking the dimensions of invariants in the ramification
filtration.

Write \(q\) for the rational prime under \(z\), and
\(a_z=v_z(N_*)\) in the normalized integral valuation. In local
reciprocity the finite character on \(K_z^\times\) is pulled back
by the norm to \(\QQ_q^\times\). If \(q\nmid N_*\) it is
unramified. Otherwise an element of \(1+N_*\mathcal O_{K_z}\)
acts by multiplication on the free \(\ZZ_q\)-module
\(\mathcal O_{K_z}\) by a matrix congruent to the identity modulo
\(q^{v_q(N_*)}\). Its determinant, the field norm, has the same
congruence. This proves triviality on \(1+\mathfrak p_z^{a_z}\),
including at the ramified Gaussian prime over two. The Artin
conductor exponent is thus at most \(2a_z\).

The local component of a cuspidal representation of
\(\mathrm{GL}_3\) is generic. Local--global compatibility away
from \(\ell\) and preservation of conductors under the local
Langlands correspondence identify its conductor exponent with
that of the parameter just calculated. The local newform
theorem gives a vector fixed by the subgroup whose last row
is congruent to \((0,0,1)\) at the conductor exponent. This
subgroup contains the principal subgroup of exponent \(2a_z\).
For exponent zero it is the full maximal compact group. Taking
the finite tensor product of these vectors with the spherical
vectors elsewhere gives the claimed level. Principal level
\(4N_*^2\) is contained in this level, so the list used in
the interval estimate includes the resulting trace system.
\,\cite{VIII-VI-V-shin}\ \href{https://lesesvre.perso.math.cnrs.fr/newforms-references/jpss.pdf}{Jacquet, Piatetski-Shapiro and Shalika, \emph{Conducteur des repr\'esentations du groupe lin\'eaire}.}
\end{proof}

\begin{corollary}[Full lengths on the components with solvable potential automorphy]\leavevmode
\label{cor:gausssum402}
For a fixed \(x\in[V,V^{3/2}]\), take all contributing triples
\((\ell,r,[\alpha])\) of Section~\ref{sec:gaussianbudget401}
for which the trace algebra is finite and every absolutely
irreducible nonsplit point either has odd orthogonal type or
satisfies the hypothesis of Theorem~\ref{thm:gaussdesc402}.
The CM extension may be chosen independently at each point.
For this set \(\mathcal A(x)\) the same interval integer gives
\[
 \sum_{(\ell,r,[\alpha])\in\mathcal A(x)}
  [O_\alpha:\ZZ_\ell]\log\ell\,
             \operatorname{length}_{O_\alpha}S(\Psi^{-1})
             \ \le\ 2\log J_V\ll_{\psi,\varepsilon,v}V^C .
\]
\end{corollary}
\begin{proof}
Every reducible nonsplit point in the finite problem is a
middle line and an irreducible outer plane. Indeed the matrix
projectors force invariant subspaces to be coordinate subspaces,
and reflection of the adjacent entries restricts a reducible
representation to an upper triangular, lower triangular or
outer-block shape, as in the proof of
Theorem~\ref{thm:strong118}. Its triangular trace is the split
one by the fixed-point calculation. At the nonsplit outer block
Theorem~\ref{thm:form117} supplies the newform. Odd orthogonal
points are already on the list, and the theorem and proposition
put every other point of these triples on the rank-three list.
Thus \(U_{\ell,r,\alpha}=0\) there. Restrict the nonnegative sum
of Theorem~\ref{thm:gaussianbudget401} to these triples and
group the geometric characters into residue orbits. This gives
the asserted bound for their full lengths.
\end{proof}

No quantitative bound on the auxiliary solvable extensions
enters the integer \(J_V\). To include the remaining finite
components by this argument one still has to produce potential
automorphy for their remaining points. The construction of such
an extension from an arbitrary ordinary Galois point, and the
exclusion of nonconstant residual curves, remain necessary for
a uniform bound by this route.

There is a local restriction on the use of an ordinary lifting
theorem here. At every finite place away from \(\ell\) the original
point has finite inertia, also after restriction to any finite
extension. Its Weil--Deligne monodromy operator is zero. A twist
of the rank-three Steinberg representation has monodromy with
one Jordan block of size three. Thus the Steinberg place in
Thorne's Theorem~7.1 cannot be supplied by extending the base
field of this point. The passage through auxiliary deformations
in the source occurs in
\srcloc{Development/AutomorphicArithmetic/AutomorphicVarieties/ModularityArgument/ThresholdArguments/DenseThresholdCases/DnEnds}{35}.
\href{https://www.dpmms.cam.ac.uk/~jat58/reducible_lifting.pdf}{J. Thorne, \emph{Automorphy lifting for residually reducible \(l\)-adic Galois representations}.}


\subsection{Cohomology bounds over varying auxiliary CM fields}
\label{sec:cmcores403}

Fix an integer \(d\ge2\). Let \(F\) range over CM subfields of a fixed
algebraic closure of \(\QQ\) that contain \(K=\QQ(i)\), with
\([F:\QQ]\le d\) and absolute discriminant \(D_F\le X\), where
\(X\ge2\). No normality over \(K\) is imposed. We will allow a
different \(F\) at each characteristic-zero trace point. To test
all these points, fix the prime \(v\) and the place \(w\mid v\) of
\(K\) in Section~\ref{sec:gaussianbudget401}, and put
\(m=\operatorname{lcm}(1,\ldots,d)\). A Frobenius at \(w\) raised
to the power \(m\) can be evaluated by spherical Hecke operators
over every \(F\), also when \(v\) ramifies in \(F\).

\begin{lemma}[A uniform cell count in bounded degree]\leavevmode
\label{lem:cmcells403}
There are at most \((2X)^{C_d}\) such embedded fields. For each
of them and each positive integer \(N\), the adelic locally
symmetric space for \(\mathrm{GL}_3/F\), modulo the positive
real split center and at principal level \(3N\), has the homotopy
type of a finite cell complex with at most \((2D_FN)^{C_d}\) cells.
In this section \(C_d\) denotes a positive constant depending
only on \(d\), whose value can increase from one bound to the next.
\end{lemma}
\begin{proof}
Write \(b=[F:\QQ]\) and use the usual Euclidean norm on the
Minkowski embedding of \(F\). The nonzero vectors of
\(\mathcal O_F\) have norm bounded below by a constant depending
only on \(b\), by the product formula for the field norm of an
algebraic integer. Its covolume is a constant depending on the
signature times \(D_F^{1/2}\). Minkowski's second theorem gives
\(b\) independent lattice vectors of lengths at most
\(C_d D_F^{1/2}\).

We record how to pass to a basis. In a Euclidean lattice of rank
\(s\) having \(s\) independent vectors of length at most \(B\),
a shortest nonzero vector is primitive and has length at most
\(B\). Projection perpendicular to it is a lattice of rank
\(s-1\); the given vectors project to a spanning set of lengths
at most \(B\). Inductively take a short basis of the projection
and lift it, subtracting multiples of the primitive vector until
each parallel component has length at most \(B/2\). These lifts
and the primitive vector are a basis of lengths at most
\(C_s B\). Apply this to get a basis \(a_1,\ldots,a_b\) of
\(\mathcal O_F\) with all conjugates bounded by \(C_dD_F^{1/2}\).

For two different embeddings of \(F\), equality of their values
on \(\sum_j T^{j-1}a_j\) holds at at most \(b-1\) rational
integers \(T\). The difference polynomial is nonzero since the
\(a_j\) span \(F\). Some integer \(0\le T\le d^3\) separates
all pairs. The resulting primitive integer of \(F\) has a monic
minimal polynomial of degree at most \(d\) with coefficient
magnitudes bounded by \((2X)^{C_d}\), by the elementary symmetric
functions. There are polynomially many such polynomials in \(X\)
for fixed \(d\), each with at most \(d\) roots. This counts the
embedded fields as asserted.

At full integral level the finite adelic classes correspond to
rank-three \(\mathcal O_F\)-lattices \(M\) in \(F^3\) up to
\(\mathrm{GL}_3(F)\). Indeed a finite adele sends the standard
local lattices to local lattices whose intersection with \(F^3\)
is the required global lattice. The group on the corresponding
symmetric space is \(\operatorname{Aut}_{\mathcal O_F}(M)\).
Use Ash's equivariant well-rounded retract for this lattice,
with trivial weights on the vectors. Normalize the minimum
length to one. Its cells are determined by the finite sets of
shortest vectors, these sets span \(F^3\) over \(F\), and passing
to a face enlarges the set. Barycentric simplices correspond
to flags of these cells. This is the number-field retract for
marked lattices in the construction of Ash and McConnell.
\href{https://arxiv.org/abs/math/9611220}{A. Ash and M. McConnell,
\emph{Cohomology at infinity and the well-rounded retract for
general linear groups}.}

Here is a uniform count of the flags over all the lattice
classes. Take the shortest-vector set of the smallest face of
a flag, with its normalized metric. It contains three
\(F\)-independent vectors \(u_1,u_2,u_3\). The \(3b\) vectors
\(a_j u_i\) are integral and independent over \(\QQ\). Their
lengths are at most \(C_dD_F^{1/2}\), since multiplication by
\(a_j\) on an archimedean Hermitian factor has operator norm
\(\lvert\sigma(a_j)\rvert\). The short-basis argument gives a
\(\ZZ\)-basis of \(M\) with that bound up to \(C_d\). In an
orthonormal coordinate system let \(A\) be its real basis
matrix. The balls of radius \(1/2\) centered on the lattice
have disjoint interiors, so the covolume \(\lvert\det A\rvert\)
is bounded below by the volume of such a ball. Cofactor
expansion now bounds \(\|A^{-1}\|\) by \((2D_F)^{C_d}\).

Record in this basis the integer matrices for multiplication by
all the \(a_j\) and the integer coordinates of the shortest
vectors in the chosen face. Every entry has magnitude at most
\((2D_F)^{C_d}\) by the two basis bounds. There are at most
\(3^{3b}\) shortest vectors, since the disjoint balls around
them lie in the ball of radius \(3/2\). Record also the nested
subsets belonging to the other faces in the flag. The number
of such subsets and the length of a strict chain are bounded
solely in terms of \(d\). There are thus at most \((2D_F)^{C_d}\)
records. Two coinciding records give a basis isomorphism
commuting with every \(a_j\). It is \(\mathcal O_F\)-linear and
extends to an \(F\)-linear isomorphism of the rational spaces.
It identifies the flags since it identifies every specified
shortest-vector set. Consequently each record accounts for at
most one orbit, even when all projective lattice classes are
counted together.

Passing to principal level \(3N\) introduces at most
\(\lvert\mathrm{GL}_3(\mathcal O_F/3N)\rvert\le(3N)^{9b}\)
lifts per orbit: the possible additional level coordinates
form a quotient of the integral finite adelic group by its
principal subgroup. Stabilizers of points on the retract are
finite, since the minimum is normalized and a lattice has
finitely many isometries. At this level they act trivially,
by the torsion-freeness of the principal subgroup modulo three
on the underlying integral lattice. The quotients of the
retracts therefore give cell complexes of the required total
size and homotopy type.
\end{proof}

\begin{proposition}[A uniform list of Frobenius power traces]\leavevmode
\label{prop:cmpower403}
Let \(v\nmid 3N\) and \(h\ge1\). Take all the above fields and
all regular algebraic conjugate self-dual cuspidal representations
of \(\mathrm{GL}_3(\mathbb A_F)\) with trivial unitary central
character, principal level \(N\) fixed vectors, and cohomological
coefficient weights bounded in absolute value by \(h\). For every
\(z\mid w\) put \(f=[k(z):\mathbb F_v]\). Record each
representation, place and sign, with value
\[
             b_{\pm}=\sum_{j=1}^3 a_{z,j}^{\,\pm m/f},
\]
where the \(a_{z,j}\) are the unitary Satake parameters. Include
each conjugate of each value, keeping the labels distinct. The
size of this list is at most \((2XN(h+1))^{C_d}\); every value
has absolute value at most three. Multiplication by
\(v^{m(7dh+d)}\) makes each value an algebraic integer.
\end{proposition}
\begin{proof}
For one field, one weight has coefficient dimension at most
\((2h+2)^{3b}\) by the \(\mathrm{GL}_3\) Weyl dimension
formula at the \(b\) embeddings. There are at most
\((2h+1)^{3b}\) weights. At level \(3N\), cellular cochains on
the complex in the lemma bound the total cohomology dimension
by its cell count times the coefficient dimension. Every one
of our cuspidal representations contributes to that cohomology
by the cuspidal direct summand of Franke's theorem, as in
Lemma~\ref{lem:coh}. This bounds their number.
\,\cite{VIII-VI-V-franke}

To bound the degrees and denominators, define the coefficient
system over the normal closure of \(F\), of degree at most
\(d!\). At a finite adelic representative take the tensor
product of the integral coefficient lattices at the embeddings
of its underlying lattice \(M\), using Schur modules and
determinant powers as in Lemma~\ref{lem:int}. For the double
coset of \(\operatorname{diag}(\pi_z,1,1)\) or its inverse the
neighboring lattices satisfy \(vM\subseteq M'\subseteq v^{-1}M\).
This holds at \(z\) even in the ramified case since \(v\) belongs
to its maximal ideal. The same inclusions hold for the duals.
At each embedding the polynomial part of the weight uses at
most \(4h\) tensor factors; its determinant power uses at most
another \(3h\), including negative powers by duality. Hence
\(v^{7dh}\) times either Hecke operator preserves the integral
cohomology lattice modulo torsion, by integral pullback and
transfer. The unscaled spherical eigenvalues are
\(v^f s_{\pm}\), with \(s_{\pm}=\sum_j a_{z,j}^{\pm1}\).
Thus \(v^{7dh+d}s_{\pm}\) are algebraic integers, and their
degrees over \(\QQ\) are at most \(d!\) times the cohomology
dimension. The bound includes nonfree coefficient lattices by
localizing their characteristic polynomials at every finite
prime.

The product of the three Satake parameters is one. For either
sign the three roots have polynomial
\(Y^3-s_\pm Y^2+s_\mp Y-1\). Newton's identities express their
\(j\)th power sum as an integral polynomial in \(s_+,s_-\) of
total degree at most \(j\). This follows successively from the
first three identities and the cubic recurrence. Taking
\(j=m/f\) proves the denominator bound. The degree of the
compositum of the two trace fields is at most the product of
their degrees. This costs another polynomial factor when listing
all conjugates. There are at most \(d\) choices of \(z\) per
field. Shin's purity theorem in odd rank gives absolute value
one for every conjugate of each unitary root. Summing over fields
using the lemma proves all the claims.\,\cite{VIII-VI-V-shin}
\end{proof}

\begin{theorem}[The strict order for points over bounded-degree CM fields]\leavevmode
\label{thm:cmcores403}
Use the common fields and intervals of
Section~\ref{sec:gaussianbudget401}. For each \(x\in[V,V^{3/2}]\),
take those large-prime triples \((\ell,r,[\alpha])\) for which
\(T\) is finite over \(O_\alpha\) and the following holds at every
absolutely irreducible nonsplit point for every geometric character
in the indicated residue orbit. It is odd orthogonal,
or there are CM fields \(K\subseteq F\subseteq L\) with
\([F:\QQ]\le d\), \(D_F\le X\) and \(L/F\) solvable Galois, such
that \((U\otimes\kappa_\ell^{-1})|_{G_L}\) is attached to a
regular algebraic conjugate self-dual cuspidal representation.
Denote the set of triples by \(\mathcal A_{d,X}(x)\).
There is a positive integer \(J_{V,d,X}\), independent of \(x\),
such that
\[
 \sum_{(\ell,r,[\alpha])\in\mathcal A_{d,X}(x)}
 [O_\alpha:\ZZ_\ell]\log\ell\,
 \operatorname{length}_{O_\alpha}S(\Psi^{-1})
       \ \le\ 2\log J_{V,d,X}
       \ \ll_{d,\psi,\varepsilon,v} (2VX)^{C}
\]
for a fixed exponent \(C\). Neither the degree nor the
ramification of \(L/F\) enters the estimate.
\end{theorem}
\begin{proof}
For a point of the second kind, the cyclic descent proof of
Theorem~\ref{thm:gaussdesc402} applies over \(F\). Indeed
\(U|_{G_F}\) has determinant one, remains absolutely irreducible
on every open subgroup and is carried to its dual by CM
conjugation. These are exactly the properties of the base
representation used at the cyclic steps. The intermediate
fields in \(L/F\) are CM because they contain \(F\) and CM
conjugation commutes with its Galois group. We obtain a
cuspidal system for the actual restriction to \(F\), with
trivial unitary central character.

It has level \(N=N_*^2\) fixed vectors. At \(z\nmid\ell\)
the finite inertia characters are the restrictions of
\(\alpha,1,\alpha^{-1}\). For \(z\mid q\mid N_*\), the norm
to \(\QQ_q\) sends \(1+N_*\mathcal O_{F_z}\) into
\(1+N_*\ZZ_q\) by the determinant of multiplication on the
local integral basis. The conductor is at most \(2v_z(N_*)\);
at the other places away from \(\ell\) it is zero. At \(\ell\)
restriction preserves crystallinity and pulls each labeled Hodge
weight back along the embedding. Thus the spherical and
highest-weight conclusions of Proposition~\ref{prop:gausslevel402}
also apply over \(F\), with \(h=4(r+2)\). The tangent result
of Corollary~\ref{cor:tangent121} applies with this CM field,
giving multiplicity one in the invariant algebra by the proof
of Theorem~\ref{thm:gaussianbudget401}. For reducible nonsplit
points and odd orthogonal points that theorem gives the same
multiplicity and the two classical modular descriptions.

Choose a Frobenius lift \(g\) at \(w\) and evaluate a single
element of the trace augmentation ideal at every point,
\[
           t(g^m)-d_{\alpha,m},\qquad
 d_{\alpha,m}=1+(v^r\alpha(v))^m+(v^r\alpha(v))^{-m}.
\]
For \(z\mid w\), a Frobenius at \(z\) maps into the conjugacy
class of a lift of the \(f\)th power of Frobenius at \(w\),
modulo inertia there. The original representation at \(w\)
is unramified. The trace at \(g^m\) is therefore one of the
power traces in the proposition. Moreover one global cuspidal
representation over a fixed embedded \(F\), with the embedding
of coefficients fixed, accounts for at most one trace point
of a fixed triple. To see this, two representations giving
that label agree on \(G_F\) by Chebotarev. Identify them there.
On the group of the normal closure of \(F/K\) their common
restriction is absolutely irreducible. Schur's lemma then makes
their ratio on \(G_K\) a character \(\xi\), trivial on \(G_F\).
The determinants force \(\xi^3=1\). If \(\xi\ne1\), its reduction
still has order three. Multiplication by this residual
character preserves \(\{1,\bar\Psi,\bar\Psi^{-1}\}\), the common
residual semisimplification on \(G_K\); the orbit of \(1\)
would exhaust this set. But on inertia at \(\ell\) the order
of \(\bar\Psi\) is at least \((\ell-1)/r>3\), since \(K\) is
unramified there, \(\alpha\) is unramified and
\(r<t(x)\ll P(x)\). This excludes the twist. The group traces
generate \(T\), so equality of the representations identifies
their trace characters. If several fields or labels work for
a given point, select one.

We describe the integer explicitly to keep the prime sum uniform.
For each twist \(r<t_*\) starting at one and each geometric
character \(\alpha\) of the common field \(E_*\) of the required
parity, take all labels of the proposition at \(N_*^2\) and
\(h=4(r+2)\). Add the power traces of the two modular families
in Theorem~\ref{thm:gaussianbudget401}, recording their
newform labels and both Frobenius conventions. For the plane
with roots \(\beta_1,\beta_2\) of the weight-\(2r+1\) form,
these are the \(m\)th power sums of
\(1,\alpha(v)\beta_1/v^r,\alpha(v)\beta_2/v^r\) and their
inverses. For the adjoint with roots \(\gamma_1,\gamma_2\) of
the weight-\(r+1\) form they are those of
\(1,\gamma_1/\gamma_2,\gamma_2/\gamma_1\). The products of
the two modular roots are a root of unity times \(v^{2r}\)
and \(v^r\), respectively, so \(v^{mr}\) clears the denominators
in both power traces. At the split prime \(v\) these are exactly
the two families with their unitary normalizations. Deligne
purity bounds every conjugate by three. The modular dimension
and degree bounds of Theorem~\ref{thm:traces} remain polynomial
after taking these fixed powers and adjoining the character
values of \(E_*\). Distinct points using the same modular label
coincide because it specifies the global plane or adjoint.

Set \(u_r=m(7dh+d+r)\). Multiply \(v^{u_r}(b-d_{\alpha,m})\)
over these labeled entries \(b\), twists and characters. Take
all conjugates of the rank-three values as in the proposition,
and take the modular lists with their matching Nebentypus,
unitary scalars and Galois conjugates as in
Theorem~\ref{thm:gaussianbudget401}. The product is Galois
invariant; conjugation permutes fields of character values,
newforms and the corresponding character labels together.
Each factor is an algebraic integer. It is nonzero since
\(\lvert d_{\alpha,m}\rvert\ge v^{rm}-1-v^{-rm}>3\)
at every embedding. Take \(J_{V,d,X}\) to be the absolute value
of the product, with empty product one. Each factor has
logarithmic absolute value \(O_{d,v}(r+1)\). The number of
entries is polynomial in \(X,V\) by the proposition and the
bounds for \(N_*,t_*\) and \([E_*:\QQ]\) on the polynomial
conductor intervals. This proves the claimed size.

For any included triple of positive strict length apply
Theorem~\ref{thm:allhorizontal401} to the displayed trace
difference. All nonsplit points occur with multiplicity one,
and the absolute values just proved ensure that its determinant
is nonzero. Distinct nonsplit characters can be assigned
distinct labels by the preceding uniqueness arguments. A fixed
embedding into \(\overline{\QQ}_\ell\) identifies the geometric
characters at \(x\) with a subset of those of \(E_*\);
\(v\) is an \(\ell\)-adic unit. Thus at a given \(\ell\) the sum
of the lengths for all included geometric characters and twists
is at most \(2v_\ell(J_{V,d,X})\). Each included orbit has as many geometric characters as the
residue degree by unramifiedness of \(O_\alpha/\ZZ_\ell\). Summing with \(\log\ell\) gives the stated
residue-degree formula.
\end{proof}

\begin{corollary}[Ramification on the common conductor]\leavevmode
\label{cor:cmram403}
For fixed \(d\), the full large-prime strict contribution in
the theorem is polynomial in \(V\) if at each point an
intermediate field \(F\) of degree at most \(d\) can be chosen
unramified over \(\QQ\) outside the primes dividing \(2N_*\).
More generally intermediate fields of discriminant at most
\(\exp(V^\varepsilon/\log(2V))\) give a total
\(o(\exp(x^\varepsilon/2))\) uniformly for
\(V\le x\le V^{3/2}\).
\end{corollary}
\begin{proof}
At a rational prime \(q\), pass in \(F_z/\QQ_q\) to its
maximal unramified subfield. The remaining extension of
ramification index \(e\) is generated by a uniformizer with
Eisenstein minimal polynomial. Its different is generated by
the derivative at that uniformizer. The nonzero derivative
terms have distinct valuations modulo \(e\), so its valuation
is at most that of the leading term, \(e-1+e v_q(e)\).
The unramified subextension has trivial different. Summing
over \(z\mid q\) with residue degrees bounds the discriminant
exponent by \(d(1+\log_2 d)\). Under the given ramification
hypothesis this yields \(D_F\le(2N_*)^{d(1+\log_2 d)}\),
polynomial on the common interval. For the second claim the
logarithm of the bound \((2VX)^C\) is
\(O(\log(2V)+V^\varepsilon/\log(2V))=o(V^\varepsilon)\).
Both conclusions follow from the theorem.
\end{proof}

In the source the auxiliary field, the seed representations
and the passage to other primes of the ordinary deformation
ring enter at
\srcloc{Development/AutomorphicArithmetic/AutomorphicVarieties/ModularityArgument/ThresholdArguments/DenseThresholdCases/LinkedEndE}{60}
and
\srcloc{Development/AutomorphicArithmetic/AutomorphicVarieties/WeilGroupCharacters/UniversalOrdinaryDeterminantRing}{284}.
The estimate above allows substantial variation of the field
where cuspidality is proved. Showing that every required
fixed-weight point has such a field, and treating the residual
curves of an infinite trace algebra, remain open in the present
argument.


\subsection{One totally real extension removing the conductor on inertia}
\label{sec:realram404}

On a polynomial-conductor interval of Theorem~\ref{thm:conductor232}
the period fields lie in their largest member \(E_*\), of conductor
dividing \(N_*\). Write \(G=\operatorname{Gal}(E_*/\QQ)\). It matters
where one removes the conductor of \(\alpha\). The auxiliary
ordinary problem uses a nontrivial cyclotomic character on the
kernel of \(\bar\Psi=\bar\alpha\bar\kappa_\ell^r\). We construct
a real field on which every \(\alpha\) is unramified at the finite
places and on which its global values are preserved.

\begin{theorem}[Removing finite ramification while preserving the character group]\leavevmode
\label{thm:realtwist404}
There are at most four distinct odd primes \(b_i>N_*\) and a
totally real abelian field \(D/\QQ\) of degree \(|G|\) and conductor
dividing \(N_*\prod_i b_i\). Both the conductor and the primes
\(b_i\) are bounded by a fixed power of the interval parameter
\(V\). If a finite Galois extension \(M/\QQ\) is unramified at all
the \(b_i\), then \(D\) and \(M\) are linearly disjoint. In
particular \(E_*D/D\) has Galois group \(G\) and is unramified
at all finite places. The field \(D\) is fixed throughout the
interval, including when the coefficient prime \(\ell>P(x)\)
varies.
\end{theorem}
\begin{proof}
By Proposition~\ref{prop:field12}, \(E_*\) is the compositum of
the cyclic field of \(\psi\), a cyclic real cyclotomic layer, and
two quadratic fields. Restriction embeds \(G\) in the product
of the four cyclic groups. Every subgroup of such a product
is generated by at most four elements, by projection onto one
factor and induction on the kernel. The structure theorem for
finite abelian groups gives independent characters
\(\chi_i\) generating the dual of \(G\), of orders \(n_i>1\),
with \(\prod n_i=|G|\) and at most four indices. The trivial
group uses an empty list.

For each \(i\) we need a character \(\theta_i\) of prime conductor
\(b_i\) with order \(n_i\) and the same parity as \(\chi_i\).
When \(\chi_i\) is even impose \(b_i\equiv1\pmod{2n_i}\).
When it is odd, \(n_i\) is even and we impose
\(b_i\equiv1+n_i\pmod{2n_i}\). Both residue classes are units.
In the cyclic group modulo such a prime an order-\(n_i\)
character takes \(-1\) to the desired sign, by evaluating at
the \((b_i-1)/2\)th power of a generator.

For a common size bound choose an auxiliary prime \(t\) between
\(4(N_*+s+10)\) and twice that number by Bertrand's postulate,
where \(s\) is the fixed prime in Proposition~\ref{prop:quadratic345}.
Use different residues strictly between \(N_*+s\) and \(t\) as
additional congruences for the \(b_i\) modulo \(t\). The Chinese
remainder theorem gives a unit class modulo \(2n_i t\), since
\(n_i\le |G|\le N_*\). Linnik's least-prime theorem bounds a
prime in each class by \(C(2n_i t)^A\) for absolute constants
\(A,C\). These primes are distinct and exceed \(N_*\) and \(s\).
Only four occur and \(N_*\) is polynomial in \(V\).
\href{https://academic.oup.com/plms/article-abstract/s3-64/2/265/1528322}{D. R. Heath-Brown,
\emph{Zero-free regions for Dirichlet \(L\)-functions, and the
least prime in an arithmetic progression}.}

Define \(D\) by the characters \(\chi_i\theta_i\). They are even.
At \(b_i\) the character \(\theta_i\) maps inertia onto all its
\(n_i\)th roots of unity. All the other \(\theta_j\) and all the
\(\chi_j\) are unramified there. The images of these inertia
groups in the product of the character groups of values therefore
contain the individual coordinate factors. The full image is
their product. This proves the degree assertion and shows that
the inertia groups at the \(b_i\) generate \(\operatorname{Gal}(D/\QQ)\).
Their images in \(D\cap M\) are trivial, hence \(D\cap M=\QQ\);
Galois theory gives linear disjointness. The conductor bound
follows by multiplication of the Dirichlet conductors.

At a prime over \(p\mid N_*\), inertia in \(G_D\) kills
\(\chi_i\theta_i\), and \(\theta_i\) is unramified at \(p\).
It thus kills \(\chi_i\) for each \(i\). At every other prime of
\(D\) the \(\chi_i\) were already unramified over \(\QQ\).
This proves the finite unramifiedness of \(E_*D/D\). Disjointness
applies to \(M=E_*\). For sufficiently distant intervals all
the displayed conductors and primes are below \(P(V)\).
\end{proof}

\begin{proposition}[Exact restriction on the original strict groups]\leavevmode
\label{prop:strictdescent404}
For \(B=\QQ\) or \(E_x\), write \(B'=BD\) and
\(\Delta=\operatorname{Gal}(B'/B)\). Take a discrete
\(\ell\)-primary Galois module \(W\) and let \(S_B(W)\) be the
kernel of restriction in \(H^1(G_B,W)\) to the full local
groups above \(\ell\) and to inertia at the other finite places.
If \(\ell\nmid |G|\), restriction gives
\[
                   S_B(W)\simeq S_{B'}(W)^\Delta .
\]
It commutes with the actions of the characters of the abelian
period fields whenever \(W\) is obtained by restriction from
\(\QQ\). In particular the original \(\mathcal S_{\ell,r}(E_x)\)
and all its large-prime character summands are recovered
with their full lengths, at every twist and every point of the
same interval. The assertion also applies over \(\QQ\) to the
finite and divisible character modules in the marked problem.
\end{proposition}
\begin{proof}
Disjointness gives \(|\Delta|=|G|\). Multiplication by this
integer is invertible on any \(\ell\)-primary module under
consideration. Averaging annihilates all positive-degree
cohomology of \(\Delta\) on such modules. Inflation--restriction
therefore identifies global \(H^1\) with the invariant part
upstairs. For full local groups the same argument gives
injectivity on restriction since the local degree divides
\(|\Delta|\). The intersection of a chosen inertia group with
\(G_{B'}\) is the inertia group upstairs, of index dividing
\(|\Delta|\). Inflation--restriction again gives injectivity.
A class corresponding under the global isomorphism meets
all the local zero conditions downstairs exactly when its
restriction does upstairs. The operations are functorial for
conjugation. Finite-level character idempotents where defined
commute with restriction and averaging. For the Tate module
\((\QQ_\ell/\ZZ_\ell)(-r)\) used here, inertia away from
\(\ell\) acts trivially. A cocycle restricting to zero in its
cohomology vanishes there, so it descends to the global
quotient unramified at these places, as in
Section~\ref{sec:coho12}. This identifies the stated groups.
\end{proof}

\begin{theorem}[The ordinary problem at hyperspecial tame level]\leavevmode
\label{thm:hyperspecialframe404}
In the notation of Theorem~\ref{thm:commonframe365}, put
\(F'_a=F_a^+D\), \(L'_a=L_aD\), and retain
\(K_s=\QQ(\sqrt{-s})\). For each \(\ell>P(V)\), a single sufficiently
deep layer \(a\) works for all the twists and conductor characters
on the interval. On this layer the local, cohomological and
numerical assertions of that theorem hold with \(F'_a,L'_a\)
and a possibly different integer \(C_0\). The marked global
ramification quotient over \(F'_a\) is now the group of the
maximal extension unramified at the finite places away from
\(\ell\). The cuspidal ordinary adjoints of
Theorem~\ref{thm:adjointseed361} over \(L'_a\) have hyperspecial
fixed vectors at every finite place away from \(\ell\) on the
definite unitary group over \(F'_a\).
\end{theorem}
\begin{proof}
Choose the auxiliary primes of Proposition~\ref{prop:auxfield361}
outside the \(b_i\). Its field \(L_aE_*(\mu_\ell)\) is unramified
at these primes at every depth. The theorem on \(D\) shows that
adjoining it preserves the full product of Galois groups of
\(L_a\) and \(E_*(\mu_\ell)\). Thus \(L'_a\) remains disjoint
from \(E_*(\mu_\ell)\). In particular some element of \(G_{L'_a}\)
acts as complex conjugation on \(E_*(\mu_\ell)\). Its values are
\(\bar\kappa_\ell=-1\) and \(\bar\Psi=1\) for every pair since
\(\alpha(-1)=(-1)^r\). This same element rules out every
identity \(\bar\kappa_\ell=\bar\Psi^j\) on the sign kernel.

The real field \(F'_a\) has CM quadratic extension
\(L'_a=K_sF'_a\). It is the base change of \(L_a/F_a^+\),
hence unramified at finite places and split above \(\ell\).
The ramification index of \(L'_a\) over \(\QQ\) at \(\ell\)
is still \(\ell^a\), as \(D\) is unramified there. Full rational
residual cyclotomic inertia survives. The separating inertia
element and the local principal-unit bases in the proof of
Theorem~\ref{thm:commonframe365} therefore exist over these
completions as well. In detail, the local norm on the principal
units has rank one, the local principal units have no
\(\ell\)-power torsion since they contain no \(\ell\)th root
of unity other than one, and all but one basis element can be
chosen in the norm kernel. The full local dimension is at least
that of \(F_{a,u}^+\). At infinity both the sign and the parity
are preserved by real base change. Away from \(\ell\) the
quadratic sign and \(\bar\Psi\) are now both unramified on
\(G_{F'_a}\). Killing the intersection of their inertia kernels
kills every such inertia group. The resulting Euler characteristic
is the ordinary global one for the group unramified outside
\(\ell\), with the same local correction terms as in
Lemma~\ref{lem:frameeuler365}.

The two nonzero strict classes of Proposition~\ref{prop:strictclasses365}
remain nonzero on \(G_{F'_a}\), since the degree over \(F_a^+\)
is prime to \(\ell\). Their zero restrictions above \(\ell\)
and their inertia conditions away from it persist. In particular
a cocycle on the old sign and character quotient vanishes on
the inertia groups killed on the new field. For the uniform upper
bounds set \(T'_a=L'_aE_*(\mu_\ell)\). These are cyclotomic
layers of the abelian field \(T'_0\), totally ramified up the
tower at every place above \(\ell\). Indeed the extra factors
at depth zero are unramified locally over \(\QQ_\ell(\mu_\ell)\).
As in Proposition~\ref{prop:classrank365}, Ferrero--Washington
and Iwasawa's structure theorem bound the \(\ell\)-ranks of all
their class groups by one integer \(h'\). The unramified class
extensions of \(T'_a\) are disjoint from the totally ramified
tower, so their groups are quotients of its unramified abelian
pro-\(\ell\) module. Let \(g'\) be the constant number of places
over \(\ell\). Restriction to \(T'_a\) is injective on the
strict residual groups by prime-to-\(\ell\) degree, and their
additive characters factor through the class group. This bounds
the relevant dimensions by \(h'\). For trivial coefficients
over \(F'_a\) the same restriction followed by Kummer theory
bounds \(H^1\) by \(1+g'+h'\): in the minus eigenspace the
ordinary units modulo powers contribute at most one dimension,
their quotients by their conjugates being roots of unity, and
the allowed valuations contribute at most \(g'\).
Taking \(C_0=h'+g'\), the global degree eventually exceeds
\(3C_0+7\) and the local degrees exceed five. Finiteness of
\(H^1\) with finite coefficients on open subgroups follows by
the same finite-prime Kummer argument as in
Theorem~\ref{thm:commonframe365}. The bound annihilating
everywhere unramified abelian characters on \(L'_a\) is its
finite class number. These verify the stated frame assertions.
\,\cite{VIII-VI-II-fw13,VIII-VI-II-iwasawa15}

For the adjoints, use cyclic base change along the abelian
extension \(L'_a/L_a\). Cuspidality persists at every prime
cyclic step: a self-twist by a nontrivial quotient character
would, by good-place compatibility and Chebotarev, give such a
self-twist of the normalized Galois representation. On the
smaller open subgroup its intertwiner would be scalar by
absolute irreducibility there. The twisting character would
then be one. Cyclic base change and its cuspidality criterion
apply successively. The associated representation at each step
is the restriction, by the same good-place comparison.
The local adjoint at an old conductor prime had finite
inertia \(\alpha,1,\alpha^{-1}\) and zero monodromy.
All three restrictions to inertia on \(L'_a\) are trivial.
It was unramified at the other finite primes away from \(\ell\)
and is still unramified there. The crystalline ordinary flag at
\(\ell\) restricts with its specified parallel weights. Its
spherical refinement is ordinary by the comparison of the
ordered Frobenius slopes with the ordered Hodge slopes in
Section~\ref{sec:ordinaryfamily363}. The unitary descent used
there applies in odd rank since the CM extension is unramified
at finite places and split above \(\ell\). It has spherical
components at every finite place away from \(\ell\).
\,\cite{VIII-VI-I-ac123,VIII-VI-geraghty363,VIII-VI-V-shin}
\end{proof}

\begin{corollary}[A common tame level for the original Gaussian points and the adjoints]\leavevmode
\label{cor:onelevel404}
Set \(M_a=L'_a\QQ(i)\), with maximal totally real subfield \(M_a^+\).
The normalized adjoints above and the restriction of every
large-prime marked Gaussian representation have the same
residual semisimplification on \(G_{M_a}\) and are unramified
away from \(\ell\). At the original weights the cuspidal adjoint
gives a nonzero ordinary module for this residual system at
hyperspecial tame level. One may also impose any of the split
Taylor--Wiles levels of Section~\ref{sec:levelmaps370} at additional
places. The dual ordinary module at their diagonal cover is
finite free over \(\Lambda[\Delta_N]\) and its diamond
coinvariants are the Iwahori module. Here \(\Lambda\) is the
three-character principal weight algebra over \(M_a^+\).
The analogous module statement holds over \(F'_a\).
\end{corollary}
\begin{proof}
We have \(M_a=K_s M_a^+=\QQ(i)M_a^+\). Its quadratic extension
over \(M_a^+\) is the base change of \(L'_a/F'_a\), hence
unramified at all finite places and split above \(\ell\).
The preceding argument gives cuspidal base change of the adjoint
over the extension of degree at most two from \(L'_a\), followed
by spherical unitary descent. An original marked Gaussian
representation on \(G_{\QQ(i)}\) kills the kernel of \(\alpha\)
on inertia away from \(\ell\). On \(G_{M_a}\) this is the
entire inertia by construction. Both residual semisimplifications
are \(\bar\Psi\oplus1\oplus\bar\Psi^{-1}\), and the adjoint
exists at weight \(m=r\).

For integral ordinary control choose in advance an odd rational
prime \(p_0\) completely split in \(M_a\) whose class modulo
\(\ell\) has order \(\ell-1\). This is possible by Chebotarev:
\(M_a\) has \(\ell\)-power ramification index at \(\ell\), while
each nontrivial subfield of \(\QQ(\mu_\ell)\) has prime-to-\(\ell\)
ramification index greater than one there. Choose all additional
places away from \(p_0\). At one split place over \(p_0\)
first impose the principal subgroup modulo \(p_0\).
It is torsion-free by the matrix logarithm, giving a sufficiently
small level for ordinary control. Its quotient in the maximal
compact subgroup has order
\(p_0^3(p_0-1)(p_0^2-1)(p_0^3-1)\), a unit at \(\ell>3\)
on these distant intervals. Averaging over that quotient
integrally restores hyperspecial level. It commutes with weights,
ordinary projection, every added diamond action and change
of level at the added places. The integral freeness and
coinvariant assertions of Proposition~\ref{prop:diamond370}
therefore pass to this direct summand. The completed weight
ring and its finite abelian \(\ell\)-group ring are local, so
the finite projective summands are free. The same reasoning
applies to the common residual Hecke summand, cut out by the
local factor in the finite Hecke algebra over the complete
weight ring. Nonvanishing follows from the spherical adjoint
and its ordinary refinement. A spherical vector is still
fixed at all additional smaller levels. Over \(F'_a\) one uses
\(L'_a\) in the identical prime choice.
\end{proof}

For a strongly irreducible characteristic-zero Gaussian point,
cuspidal automorphy of its normalized restriction to \(M_a\)
would descend to \(\QQ(i)\) by Theorem~\ref{thm:gaussdesc402},
since \(M_a/\QQ(i)\) is abelian. The available adjoint supplies
a point in the ordinary module. Controlling all the components
through the prescribed residual system still requires the
propagation argument.


\subsection{The ordinary determinant on the three integral levels}
\label{sec:typeaction405}

Fix a sufficiently distant interval as in Section~\ref{sec:missing}
and one of its large coefficient primes \(\ell\). We use one of three
CM extensions \(L/F\). They are \(L_a/F_a^+\) of
Proposition~\ref{prop:auxfield361}, \(L'_a/F'_a\) of
Theorem~\ref{thm:hyperspecialframe404}, and \(M_a/M_a^+\) of
Corollary~\ref{cor:onelevel404}. Fix a pair \(r,\alpha\) from
the interval. In the first extension the allowed inertia away
from \(\ell\) is the finite quotient given by \(\alpha\). In the
other two the quotient is trivial. At the places above \(\ell\)
choose one of the two places of \(L\) over each place of \(F\)
and the three ordered residual characters of the marked problem.
The inertia characters are distinct there. Write \(\Lambda\) for
the completed weight algebra of the principal diagonal units,
with the finite prime-to-\(\ell\) diamonds fixed by these data.
We use the cyclotomic twist of Section~\ref{sec:basesupport371}
when comparing the characteristic polynomials of a Galois
representation with the spherical Hecke operators.

Let \(Q\) be a finite set of split places of \(F\), disjoint from
the earlier conductors and \(\ell\), whose residue cardinalities
are \(1\bmod\ell\). At \(Q\) use, in turn, the hyperspecial,
Iwahori, and diagonal-cover levels from
Section~\ref{sec:levelmaps370}. Denote their dual ordinary
modules, with the old local conditions specified below, by
\(M_K,M_I,M_1\). We choose the local factor of the good Hecke
algebra for the residual system
\(\bar\Psi\oplus1\oplus\bar\Psi^{-1}\). All three localizations
use good Hecke operators outside the same set of places.

\begin{lemma}[Integral selection of the old inertia]\leavevmode
\label{lem:oldidempotent405}
Over \(L_a/F_a^+\) one can retain precisely the local
automorphic types having inertia
\(\tau_v=\alpha\oplus1\oplus\alpha^{-1}\) and zero monodromy
at the old conductor places, by applying integral idempotents
at those places. These idempotents are independent of \(Q\).
Over \(L'_a/F'_a\) and \(M_a/M_a^+\) one may use hyperspecial
level throughout the old tame places. In each case \(M_1\)
is finite free over \(\Lambda[\Delta_Q]\), and its diamond
coinvariants are \(M_I\). All three modules are finite free
over \(\Lambda\) and contain the ordinary eigensystem of the
adjoint at the original weight.
\end{lemma}
\begin{proof}
First work over \(F_a^+\). Every \(v\) above a rational prime
\(p\mid N_*\) splits in \(L_a\), so the local unitary group
identifies with \(\mathrm{GL}_3(F_v)\). Its residue cardinality
is \(q_v=p^f\). The residue degree divides \(4\ell^{a+1}\)
by Proposition~\ref{prop:auxfield361}, hence \(f=c\ell^j\)
with \(c\mid4\). For \(1\le h\le3\), Fermat's theorem gives
\(q_v^h\equiv p^{ch}\bmod\ell\). Moreover
\(1<p^{ch}\le N_*^{12}<\ell\) on all sufficiently distant
intervals because \(N_*\) grows polynomially in \(V\).
Thus \(\ell\) divides none of \(q_v^h-1\). The finite groups
\[
       H_{v,t}=\mathrm{GL}_3(\mathfrak o_v/\mathfrak p_v^t)
                      \qquad(t\ge1)
\]
have orders prime to \(\ell\). Reduction modulo the maximal
ideal has a \(p\)-group kernel, and the order over the residue
field is \(q_v^3\prod_{h=1}^3(q_v^h-1)\).

Use the inertial local Langlands type \(\sigma_v\) attached
to \(\tau_v\) with zero monodromy. We recall the precise
criterion being used. If an irreducible admissible local
representation contains this compact type, the Jordan
partitions of its inertial parameter are dominated by those
of the specified parameter in each finite inertial block,
with equal total sizes. Conversely equality of the full
inertial parameter gives occurrence. The partitions of the
zero operator consist entirely of ones, so each is minimal
in dominance order. Occurrence is therefore equivalent to
the stated inertia, including when two characters coincide.
This applies also to wild finite characters. With geometric or
arithmetic reciprocity chosen consistently, simultaneous
inversion does not change this self-dual triple.
\href{https://maths.dur.ac.uk/users/jack.g.shotton/papers/gln.pdf}{J. Shotton,
\emph{The Breuil--M\'ezard conjecture when \(\ell\ne p\)},
Definition 3.4 and Theorem 3.7.}

By smoothness and finite dimension \(\sigma_v\) factors through
\(H_{v,t}\) for some \(t\). Start at principal level \(t\).
Enlarge the coefficient DVR \(O\) to contain the character
values. For \(H=H_{v,t}\) the projector on the \(\sigma_v\)
isotypic space is
\[
       \frac{\dim\sigma_v}{|H|}
             \sum_{h\in H}\operatorname{tr}(\sigma_v(h^{-1}))[h].
\]
All coefficients are integral, since the traces are sums of
roots of unity and \(|H|\) is a unit. Character orthogonality
proves idempotence in the group algebra. Use the transpose
projector on a dual. The product over the old places commutes
with every operation at \(Q\) and \(\ell\). For unramified
types one can take hyperspecial level directly. The other
two CM extensions have only these latter types by
Theorem~\ref{thm:realtwist404}.

At an extra odd split prime of primitive residue class modulo
\(\ell\), first impose principal level and then average back to
hyperspecial level as in Corollary~\ref{cor:onelevel404}.
This can be done away from the indicated finite sets. Ordinary
control and Proposition~\ref{prop:diamond370} apply at the
sufficiently small intermediate level. Each operation used on
the result is a commuting idempotent, including the common
residual localization. This gives the coinvariant assertion
and finite projectivity. The respective coefficient rings are
local, so the modules are free. The adjoint has the specified
finite inertia with zero monodromy at the old places and is
spherical at \(Q\). Its ordinary refinement and stable base
changes are given by Theorems~\ref{thm:adjointseed361}
and~\ref{thm:hyperspecialframe404}. Its vector survives all
the indicated projections.
\end{proof}

\begin{proposition}[Classical systems distinguishing endomorphisms]\leavevmode
\label{prop:ordinaryclass405}
At each of the three levels one may test identities of good
Hecke operators, acting as endomorphisms over \(\Lambda\), at
all the ordinary classical systems of sufficiently separated
dominant algebraic weights with the specified tame diamonds.
Each of these systems has a continuous semisimple polarized
Galois representation of \(G_L\). It is unramified outside
\(\ell,Q\) and the old conductor. At each old conductor place
its inertia is the prescribed finite representation. At
the chosen places over \(\ell\) it has an ordinary flag with
the ordered inertial characters given by the weight.
At hyperspecial level the representations are also unramified
at \(Q\).
\end{proposition}
\begin{proof}
Choose algebraic weights with every consecutive gap in the
archimedean infinitesimal character greater than two. They
are dense in the weight algebra. Indeed on integer exponent
tuples with any prescribed residues, including residues at
an arbitrarily high power of \(\ell\), one can increase each
successive exponent by a multiple of the common modulus so
as to enforce all the gaps. On a small principal-unit
neighborhood the logarithmic embedding matrix at each place
has nonzero determinant by separability of the local field.
The exponent tuples are dense in an open subset of
\(\mathbf Z_\ell^{\,3[F:\QQ]}\). A convergent power series
in logarithmic coordinates vanishing on that subset vanishes
by the one-variable identity theorem successively in the
coordinates. The invertible linear substitution by the
embedding matrix preserves this identity. This proves the
required detection of power series, as in
Proposition~\ref{prop:upfamily372}. At these weights ordinary
control identifies the fibers with classical spaces. On each
classical fiber the good spherical algebra is semisimple,
by the decomposition of automorphic forms and the spherical
character on each summand. Since the modules are free over
\(\Lambda\), testing at all these fibers detects every matrix
entry of a good Hecke operator.

We spell out the compatibility with inertia which is needed
in applying this test. Use the discrete automorphic transfer
for definite unitary groups in three variables and the
Galois attachment of Geraghty. The transfer at a split place
has the parameter of its \(\mathrm{GL}_3\) component. At a
hyperspecial place its unramified parameter restricts to an
unramified parameter over \(L\). A discrete global parameter
is a sum of cuspidal blocks \(\pi_b\boxtimes[d_b]\), where
\([d_b]\) contributes consecutive archimedean exponents
with spacing one if \(d_b>1\). This is excluded by the gaps.
Thus the local transfer is the sum of the local parameters
of the cuspidal constituents, and each is regular at infinity
in its cohomological normalization. This also explains why
endoscopic systems enter the argument. The normalized global
Galois representation is the semisimple sum of the attached
constituents, as follows from its good Frobenius polynomials
and Chebotarev. In rank \(n_b\) the twist
\(|\det|^{(n_b-3)/2}\) of the unitary normalization aligns
the local Langlands normalization \(|\det|^{(1-n_b)/2}\)
with the common \(|\det|^{-1}\) in rank three. The additional
common cyclotomic twist in our determinant is unramified
away from \(\ell\).
\cite{VIII-VI-geraghty363}
\href{https://press.princeton.edu/books/ebook/9781400882441/automorphic-representation-of-unitary-groups-in-three-variables-pdf}{J. Rogawski,
\emph{Automorphic Representation of Unitary Groups in Three Variables}.}

At a place of residue characteristic different from \(\ell\),
local--global compatibility for a regular algebraic cuspidal
constituent gives the semisimple Weil parameter and bounds
its Galois monodromy by that of its automorphic parameter
in dominance order. For a rank-one constituent this is local
class field theory. Zero monodromy on the transferred sum
means zero on every constituent, hence zero on the Galois
side. The finite inertia is now determined by its
semisimplification in characteristic zero. At each of the
old places the preceding lemma gives just such a parameter.
The same argument gives unramifiedness at every hyperspecial
place away from \(\ell\), including the inert places and,
at level \(K\), the places of \(Q\).
\href{https://www.math.toronto.edu/~ila/lgellneqp7}{I. Varma,
\emph{Local--global compatibility for regular algebraic
cuspidal automorphic representations when \(\ell\ne p\)},
Theorem 1.}

At the places over \(\ell\) the ordinary flag statement for
regular classical ordinary forms is Geraghty's Galois
attachment and ordinary comparison, Proposition~2.7.2 and
Corollary~2.7.8 in his original numbering. Its hypotheses
here are the splitting above \(\ell\), the regular dominant
weight, and the unit ordinary eigenvalues for the normalized
operators. These hold for each ordinary eigensystem being
tested. In particular the ordered inertial characters are
the algebraic characters of the diagonal units, with the
fixed cyclotomic normalization. Together with conjugate
duality and the good-place formula from the same attachment
this proves the proposition. \cite{VIII-VI-geraghty363}
\end{proof}

For the fixed host field write \(R_Q\) for the universal
ordinary polarized determinant ring of
Section~\ref{sec:basesupport371}, allowing the extra
ramification at \(Q\) and its conjugates. Write \(R_0\)
for the original ring. Both have the full principal weight
algebra, the same ordered residual characters and the same
fixed finite diamonds. Killing the added inertia gives the
surjection \(R_Q\to R_0\). Let \(B_Y\) be the \(\Lambda\)-subalgebra
of the endomorphisms of \(M_Y\), for \(Y=K,I,1\), generated
by the normalized spherical characteristic coefficients and
inverse determinants at good split places outside all the
levels. These algebras use only operators common to the
three levels.

\begin{theorem}[Integral actions of the ordinary determinant]\leavevmode
\label{thm:alllevels405}
There are continuous surjections \(R_Q\to B_Y\) for all
three levels. They give the good Frobenius polynomials and
the specified ordered inertial weights on the entire modules.
The action on \(M_K\) factors through \(R_0\). Passage
from \(M_1\) to \(M_I\) by diamond coinvariants is compatible
with these actions.
\end{theorem}
\begin{proof}
Each \(B_Y\) is a submodule of a finite matrix algebra over
the Noetherian ring \(\Lambda\), hence is finite over
\(\Lambda\). Its evaluations at the classical systems in
the proposition embed it in the product of their coefficient
valuation rings, by the detection assertion. In particular
it is reduced. In the topology of its finite \(\Lambda\)-module
it is compact and these evaluations are continuous, so the
embedding is a homeomorphism onto a closed subspace.
Good split Frobenius classes are dense up to conjugacy in
\(G_L\). Indeed Chebotarev gives positive density for any
class in a finite quotient, and the primes of \(L\) above
inert places of \(F\) have density zero. At these classes
the tuple of classical traces belongs to \(B_Y\) by the good
Hecke formula. At an arbitrary \(h\in G_L\) the trace tuple
belongs to the same closed image by continuity. This defines
a continuous trace \(t_Y\) with values in \(B_Y\).

The characteristic polynomial identities hold in the product
and therefore in \(B_Y\). Since \(3!\) is invertible, the
degree-three determinant on a group algebra element is recovered
from the traces of its first three powers. Applying this to
a generic linear combination of group elements checks the
identities of polynomial laws over arbitrary coefficient
extensions. The trace identities, continuity and polarization
thus give the polarized determinant integrally.
The reduction at the closed point has the prescribed
semisimplification by the good-prime localization and
Chebotarev. For \(g\) in any inertia kernel imposed on the
global quotient and \(h\in G_L\), one has \(t_Y(hg)=t_Y(h)\)
by the proposition. Conjugation, multiplication and closure
extend this to the normal subgroup generated by these kernels.
Consequently the determinant factors through the required
quotient. At hyperspecial level the same proof kills the
additional inertia at \(Q\). \cite{VIII-VI-chenevier364}

For clarity the full ordered local characters also take
values in this good Hecke algebra. At each chosen place
above \(\ell\) select \(s\) in inertia separating the three
residual characters, and let \(a_i\) be the universal
inertial weight values there. Their pairwise differences
are units. For \(h\) in the local decomposition group put
\[
      \xi_i(h)=t_Y\bigl(P_i(s)h\bigr),\qquad
      P_i(X)=\prod_{j\ne i}(X-a_j)/(a_i-a_j).
\]
The trace here is extended linearly. At every classical
test point \(s\) is diagonalizable with these distinct
eigenvalues. The ordered flag is the successive span of its
eigenlines. Thus \(\xi_i(h)\) is the \(i\)th diagonal character.
Multiplicativity, the value at the identity and the prescribed
restriction on inertia follow by detection. For arbitrary
local elements \(h_1,h_2,h_3\), the ordered product of the
three differences \(h_i-\xi_i(h_i)\), with \(i=3\) acting
first on the right, lowers these three flag steps successively
to zero. The trace after multiplication by every global
group element vanishes identically by the same test. This
is the ordinary relation in the polarized determinant
problem. The universal property now gives the asserted
maps. Their images contain the weights and all good
characteristic coefficients. The inverses of the unit
determinants come from inverse group elements. The images
therefore contain the generators of \(B_Y\).

The diamond action commutes with the good operators and
weights. The quotient to \(M_I\) consequently commutes
with the continuous traces at every global element by the
same good Frobenius density. The displayed projector formula
then gives compatibility of the ordered local characters.
These data determine the maps from \(R_Q\).
\end{proof}

\begin{corollary}[Annihilation of additional inertia on coinvariants]\leavevmode
\label{cor:annram405}
Write \(J_Q=\ker(R_Q\to R_0)\). For the local Bernstein
center at \(v\in Q\) use the variables \(X_{v,i}\) and
residue cardinality \(q_v\) of Section~\ref{sec:levelmaps370}.
On the integral ordinary Iwahori module one has
\[
       C_Q\,J_Q M_I=0,\qquad
       C_Q=\prod_{v\in Q}\;\prod_{i\ne j}
                         (X_{v,i}-q_vX_{v,j}).
\]
After inverting \(C_Q\), the Iwahori module is the
hyperspecial module with the roots of the local spherical
polynomials adjoined in all orders. The isomorphism respects
the full action of \(R_Q\). These assertions apply to the
actual diamond coinvariants \(M_1/(\Delta_Q-1)M_1\).
\end{corollary}
\begin{proof}
Write \(Z_Q\) for the tensor product of the local centers
and \(S_Q\) for the corresponding algebras with ordered
roots, tensored over \(O\). Apply the two maps of
Theorem~\ref{thm:iwahori370} to \(S_Q\otimes_{Z_Q}M_K\)
and \(M_I\). The averaging idempotents, ordered variables
and local Hecke operators used in these maps are supported
at \(Q\). They therefore commute with the good operators
outside the levels and with the ordinary weights. They
are continuous on these finite weight modules. For every
element of \(G_L\), approximate its conjugacy class by
the same good Frobenius classes on both sides. It follows
that the maps intertwine the two universal trace actions.
The formula for \(\xi_i\) proves the same for every local
character. A determinant with \(3!\) invertible and its
ordered local characters determine its universal action,
so both maps are \(R_Q\)-linear. Equivalently, the traces
of group algebra elements and the values of the local
characters topologically generate the representing ring.
This follows by forming its closed subalgebra generated by
these coefficients and applying the universal property to
the resulting pair over that subalgebra, or directly from
the generators of the universal trace algebra.

Denote the map from \(M_I\) to the ordered hyperspecial
module by \(\Psi\) and the reverse map by \(\Phi\).
If \(a\in J_Q\), then \(\Psi(ay)=a\Psi(y)=0\) by the theorem.
Applying \(\Phi\) and using \(\Phi\Psi=C_Q\) gives
\(C_Qay=0\) for every \(y\). The second composite is also
\(C_Q\), so \(C_Q^{-1}\Psi\) is an inverse to \(\Phi\)
after localization. The identification with coinvariants
was proved in Lemma~\ref{lem:oldidempotent405}.
\end{proof}

For the auxiliary sets and the characteristic-\(\ell\) curves
of Corollary~\ref{cor:levelcurve370} on the original host,
\(C_Q\) specializes to a unit times \(z^{D_0}\), with \(D_0\)
independent of the level. Thus the extra inertia on these
specialized coinvariants has an integral annihilator of
bounded order in the curve parameter. The identities occur
on the same modules and use the same global determinants
as the two level maps. Propagation from their support to every
required mixed-characteristic component, and the projected
period estimate on the fields \(E_x\), remain open here.


\subsection{Frobenius and inertia on the ordered diagonal cover}
\label{sec:orderedinertia406}

Keep the fields, coefficient DVR \(O\), tame types, ordinary
weight algebra \(\Lambda\), and modules of
Section~\ref{sec:typeaction405}. We give the relations before
taking diamond coinvariants. In particular the inertia
characters in them can have order growing with the auxiliary
level. Write \(t\) for the trace from
Theorem~\ref{thm:alllevels405}, acting on \(M_1\). Adjoin
to its commuting operators the normalized diagonal
Hecke operators at \(Q\). Trace on the group algebra with
coefficients in this commutative algebra means linear
extension of \(t\). The variables at \(Q\) act on the level
module, while the group elements inside the trace belong to
the global Galois group.

Choose one place \(w\) of \(L\) above each \(v\in Q\), and
let \(q_v\) be its residue cardinality. Use a geometric
Frobenius lift \(\phi_v\) and an inertia lift \(\tau_v\) whose
image in the maximal \(\ell\)-quotient of the residue units
generates it. Local reciprocity sends a uniformizer to
geometric Frobenius. Use the uniformizer, including its
unit factor, representing the actual abelian image of
\(\phi_v\), and the unit representing \(\tau_v\).

\begin{lemma}[Hecke polynomials at elements of the local Weil group]\leavevmode
\label{lem:weiltorus406}
There are commuting integral operators \(T_{v,i}(h)\),
\(1\le i\le3\), for \(h\) in the local Weil group at \(w\).
For each \(i\) they give a character of its abelianization
into the units of the commutative operator algebra. Their
values on \(\tau_v\) are the three coordinate diamonds
\(D_{v,i}\). Put \(U_{v,i}=T_{v,i}(\phi_v)\). At every
classical test system of Proposition~\ref{prop:ordinaryclass405},
with normalized Galois representation \(\rho\), one has
the following identity of polynomials acting on its whole
local fixed-vector space:
\[
                 \prod_{i=1}^3(X-T_{v,i}(h))
                     =\det(X-\rho(h))\,\mathrm{id}.
\]
The coefficient identities also hold on \(M_1\), with the
global determinant coefficients in place of the classical
ones.
\end{lemma}
\begin{proof}
Use the Bernstein homomorphism for the diagonal cover of an
Iwahori. It gives a homomorphism from the group algebra of
the diagonal torus modulo the kernel of the residue diamonds
to the cover Hecke algebra. On diagonal units it is the
diamond action. The normalized coordinate homomorphisms of
Newton and Thorne use the unramified factors on uniformizers
for the local Langlands normalization. Their products and
inverses are integral, since the residue cardinality is an
\(\ell\)-adic unit. The chosen square root of \(q_v\equiv1\bmod
\ell\), where needed, also belongs to \(O^\times\).
For a dual module take the Hecke action transposed from the
one on forms, with the same eigenvalue convention. The
torus algebra is commutative. Normalize each coordinate by
the same additional unramified character to account for the
fixed global cyclotomic twist. This does not affect the
coordinate diamonds.

The local theorem used here is the characteristic-polynomial
formula (4.1) of Newton and Thorne. It says that each elementary
symmetric function of the three torus operators at an element
of the local multiplicative group acts on the entire space
of fixed vectors of an irreducible representation by the
corresponding coefficient of its semisimple Weil parameter.
The formula holds for every Weil element with that image under
reciprocity. Its scalar assertion follows from the description
of the center of the cover Hecke algebra. The local operators
and the original proof of this assertion are in Allen et al.,
Section 2.2, equation (2.2.7) and Proposition 2.2.8. On the
Iwahori subspace the torus operators for a uniformizer are
the Bernstein variables of Section~\ref{sec:levelmaps370},
with the common normalization just specified. For positive
translations the unipotent representatives of the Hecke
cosets are identical at both levels, since the two subgroups
differ only in the diagonal units. Multiplicativity gives
the assertion for an arbitrary torus element.
\cite{VIII-VI-VI-nt119}
\href{https://arxiv.org/pdf/1812.09999}{P. B. Allen et al.,
\emph{Potential automorphy over CM fields}, Section 2.2.}

The split-place transfer and local--global compatibility used
in Proposition~\ref{prop:ordinaryclass405} identify the
semisimplified Weil parameters at every \(v\in Q\).
Monodromy does not alter the characteristic polynomial of
a Weil element. Indeed the Weil action preserves the
filtration by the kernels of the powers of its monodromy
operator, and the exponential of a multiple of this operator
acts as the identity on the successive quotients. Thus the
local formula gives the polynomial of the actual Galois
matrix. Every coefficient on \(M_1\) is an endomorphism of
a finite free \(\Lambda\)-module. At the dense weights its
matrix specializes to the identity asserted on each classical
summand. Detection of the matrix entries gives the identity
on the integral module. This argument only needs the good
Hecke algebra to act semisimply at the test weights.
\end{proof}

For a fixed \(v\) temporarily write \(\phi,\tau,U_i,D_i,q\)
for this place's data. Put
\[
 A_v=\prod_{i\ne j}(U_i-U_j),\qquad
 C_v=\prod_{i\ne j}(U_i-qU_j),\qquad f_v=A_v C_v .
\]
The products are over ordered pairs. They are symmetric in
the \(U_i\). For any global group element \(g\) put
\[
 \begin{split}
 L_v(g)&=t\!\left(g\prod_{i=1}^3(\phi-U_i)\right),\\
 H_{v,i}(g)&=t\!\left(g\prod_{j\ne i}(\phi-U_j)
                                  (\tau-D_i)\right).
 \end{split}
\]

\begin{theorem}[Integral inertia relations with exponent one]\leavevmode
\label{thm:inertiatrace406}
On the full diagonal-cover module \(M_1\), for every \(g,v,i\),
\[
                   L_v(g)=0,\qquad f_v H_{v,i}(g)=0 .
\]
There is no dependence of the exponent on the residue
cardinality, the depth, the weight, or the number of
additional places. The assertions persist after arbitrary
specialization of the coefficient algebra.
\end{theorem}
\begin{proof}
Work on one classical system and its space of vectors at the
given level. The first identity follows from the lemma and
the Cayley--Hamilton theorem for the matrix \(F=\rho(\phi)\).
One may use the scalar characteristic polynomial of \(F\)
in the product before multiplying by \(\rho(g)\) and taking
the trace. For the second identity the lemma implies that
\(f_v\) acts by a scalar, since it is a symmetric function.
If this scalar is zero the assertion holds on the entire
space. Suppose it is nonzero. The three roots \(u_i\) of the
polynomial of \(F\) are then distinct, and no two have ratio
\(q\) or \(q^{-1}\).

Write \(N\) for Galois monodromy at this place. The relation
\(FNF^{-1}=q^{\pm1}N\) forces \(N=0\), since conjugation by
\(F\) has only the pairwise root ratios as eigenvalues. Thus
the inertia image is finite. We verify that it commutes with
\(F\). Over an algebraically closed coefficient field the
commuting torus operators on the nonzero local fixed-vector
space have a simultaneous eigencharacter. Evaluating the
lemma there gives the character of the semisimple Weil
representation as a sum of three characters of the local
abelianization. On the finite inertia image the actual
representation is semisimple. Its inertia characters therefore
extend to characters of the Weil group and are fixed under
conjugation by \(F\). It follows that \(F\) preserves their
isotypic spaces and commutes with inertia. In the three
eigenlines of \(F\) write \(d_i'\) for the eigenvalue of
\(\rho(\tau)\).

Each of the operators \(U_i\) satisfies the same square-free
cubic, by substituting it for \(X\) in the identity of the
lemma. They are simultaneously diagonalizable. The \(D_i\)
commute with them and are diagonalizable in characteristic
zero because they have finite order. Choose any joint
eigenspace. Order the \(u_i\) by the eigenvalues of the three
operators \(U_i\) there, and write \(d_i\) for the eigenvalue
of \(D_i\). The cubic proves that all three roots occur in
this order. Apply the trace coefficient of the lemma at
\(\phi^k\tau\) for \(k=0,1,2\). Its two sides read
\[
                  \sum_i u_i^k d_i=\sum_i u_i^k d'_i .
\]
The Vandermonde determinant is nonzero, so \(d_i=d_i'\) for
all \(i\). On the Galois space the matrix
\(\prod_{j\ne i}(F-u_j)(\rho(\tau)-d_i)\) vanishes. Indeed
the first factor vanishes on the other eigenlines and the
second vanishes on the \(i\)th. Taking the trace after any
\(\rho(g)\) proves the desired identity on this joint
eigenspace. These eigenspaces exhaust the fixed vectors in
the case under consideration. The calculation is unchanged
in the tensor product with the vectors at all other places.

The dense classical fibers detect entries of matrices on
\(M_1\), so both identities hold over \(\Lambda\), including
above residual coincidences among the roots. Base change
preserves zero endomorphisms. This proves the two integral
relations. In the notation of
\srcloc{Development/AutomorphicArithmetic/AutomorphicVarieties/PatchingModules/ThetaCoefficientMapEq}{224},
multiplying by the other \(f_w\) consequently gives both
trace relations there with exponent one.
\end{proof}

\begin{proposition}[Selecting any ordering of the residual roots]\leavevmode
\label{prop:orderedpiece406}
At each \(v\in Q\) prescribe any ordering of the three
residual roots of \(\det(X-\bar\rho(\phi_v))\). Repetitions
are allowed. The associated summand \(M_1^{\rm ord}\) of
\(M_1\) is finite free over \(\Lambda[\Delta_Q]\) and its
coinvariants are the corresponding Iwahori summand
\(M_I^{\rm ord}\). If \(B=R_0\) and \(P^{\rm ord}\) is the
residual local factor of the ordered splitting algebra of
the universal Frobenius cubics over \(B\), then
\(P^{\rm ord}\otimes_B M_K\) is the corresponding summand
on the ordered hyperspecial side of Corollary~\ref{cor:annram405}.
\end{proposition}
\begin{proof}
Consider the commuting algebra generated by the \(U_{v,i}\),
the diamonds and \(B_1\) in the finite matrix algebra over
\(\Lambda\). It is finite over \(\Lambda\), which is complete
local. It is the product of its complete local factors, by
lifting the idempotents of its finite Artinian special fiber.
On that fiber the three \(U_{v,i}\) have joint residual values
given by orderings of the roots, by Lemma~\ref{lem:weiltorus406}.
All roots belong to the residue field after the fixed
coefficient extension since the residual characters are
defined there. Take the sum of the factor idempotents giving
the prescribed tuple. It commutes with \(\Lambda[\Delta_Q]\),
so its image is projective and hence free over that local
ring. Reducing the decomposition through the diamond
augmentation gives exactly the selected residual factors
downstairs. This proves the coinvariant assertion, since the
torus operators are compatible under change to Iwahori level.

Over \(B\), adjoin one ordered root of each monic cubic and
then one root of its remaining monic quadratic. The third
root is prescribed by the sum. All roots are units by the
unit constant term. This gives a finite free algebra of rank
\(6^{|Q|}\) over \(B\), including cubics with repeated roots.
Its local factors are obtained by the same argument with
the special fiber. In particular each prescribed ordering
gives a nonzero factor, finite free over the local ring
\(B\). On \(M_K\) the characteristic coefficients of the
universal Frobenius at \(v\) are the coefficients of the
hyperspecial Hecke polynomial, by the classical test and
the spherical case of the lemma. Substituting these values
into the presentation of the splitting algebra gives the
ordered module in the corollary. The residual decompositions
agree because they are determined by identical coordinate
operators. The two maps in that corollary commute with the
ordered coordinates and hence preserve these summands.

They are nonzero on the Iwahori side as well. Evaluate at
the adjoint. At each \(v\in Q\) this unramified adjoint has
Frobenius eigenvalues of a single complex absolute value,
by the purity of the two-dimensional modular form and
its adjoint. Consequently \(C_v\) is nonzero on its
hyperspecial vectors. The ordered splitting module there
has a point in each indicated residual ordering, by choosing
the three integral unit roots in a finite coefficient
extension and putting them in the required order. The
two composites with value \(C_Q\) then give a nonzero
Iwahori summand, and the surjection from \(M_1^{\rm ord}\)
gives its nonvanishing too.
\end{proof}

\begin{theorem}[A finite-level diagram over the universal base]\leavevmode
\label{thm:finitelevel406}
Let \(q_v\equiv1\bmod\ell^N\) for every \(v\in Q\), put
\(a_v=v_\ell(q_v-1)\), and prescribe the residual orderings
of the preceding proposition. There is a complete Noetherian
local \(O\)-algebra \(\mathcal A_Q\) carrying the universal
ordinary pair pulled back from \(R_Q\), three lifted root
coordinates and three diamond coordinates at every place.
It acts on \(M_1^{\rm ord}\) and maps to \(P^{\rm ord}\).
Under this map the root coordinates go to the ordered
universal roots, the diamond coordinates go to one, and
the ordinary pair is the pullback of the pair over \(B=R_0\).
There are \(\mathcal A_Q\)-linear maps in both directions
between
\[
   P^{\rm ord}\otimes_B M_K \quad\text{and}\quad
                  M_1^{\rm ord}/(\Delta_Q-1)M_1^{\rm ord}
\]
whose two composites equal \(\prod_{v\in Q}f_v\) on their
respective modules. On the diagonal cover the two trace
relations of the preceding theorem hold. On the curves
and auxiliary sets of Corollary~\ref{cor:levelcurve370},
the common composite has valuation \(2D_0\) before
truncation to the prescribed precision.
\end{theorem}
\begin{proof}
Lift the chosen residual values to units \(b_{v,i}\in O\).
Take
\[
            \mathcal A_Q=R_Q[[Z_{v,i},Y_{v,i}:
                                        v\in Q,\ 1\le i\le3]] .
\]
Use the root coordinates \(b_{v,i}+Z_{v,i}\) and the
diamond coordinates \(1+Y_{v,i}\). On \(M_1^{\rm ord}\)
send \(Z_{v,i}\) to \(U_{v,i}-b_{v,i}\) and \(Y_{v,i}\)
to \(D_{v,i}-1\), together with the action of \(R_Q\)
from Theorem~\ref{thm:alllevels405}. The indicated
differences are nilpotent modulo the maximal ideal of
\(\Lambda\) on the selected summand. They therefore tend
to zero under repeated powers on that finite module. All
operators commute, so this defines a continuous formal
power series action. The map to \(P^{\rm ord}\) uses
\(R_Q\to B\to P^{\rm ord}\), sends \(Z_{v,i}\) to the
\(i\)th universal root minus \(b_{v,i}\), and sends every
\(Y_{v,i}\) to zero. The residual ordering ensures
convergence there. The maps on the ordinary determinant,
its local characters and its principal weights follow
from their definitions.

Over the diamond and weight variables the diagonal-cover
module is finite free over
\[
 \Lambda[[Y_{v,i}]]\big/
       \bigl((1+Y_{v,i})^{\ell^{a_v}}-1:\ v\in Q,\ 1\le i\le3\bigr).
\]
Indeed this is the completed group ring of the product of
the cyclic groups generated by the chosen diamonds, each
of order \(\ell^{a_v}\). The assertion is exactly the
freeness of Proposition~\ref{prop:orderedpiece406}.
Putting \(Y=0\) gives the indicated coinvariants. The
support over \(B\) of the base module \(M_K\) is cut out by
\(\ker(B\to B_K)\), since \(B_K\) is its faithful algebra
of good Hecke operators and the module is finite over \(B\).

Use the two maps \(\Phi,\Psi\) in
Corollary~\ref{cor:annram405} on the selected summands.
They intertwine \(R_Q\) and the root coordinates. They
also intertwine \(Y_{v,i}\), which are zero on both sides.
Continuity gives the same for their formal power series.
Multiply \(\Psi\) by \(\prod_{v\in Q} A_v\). The two resulting
composites are exactly \(\prod_v f_v\). If the normalization
of the three Bernstein variables differs by a common unit
from that of the Hecke comparison, rescaling all three
changes \(C_v\) by the sixth power of that unit. Multiplying
the reverse map by this unit power restores the stated
scalar. The full-cover trace identities follow from
Theorem~\ref{thm:inertiatrace406}. Finally, on the specified
curves the residue characteristic is \(\ell\), so \(q_v=1\)
as a scalar and \(A_v=C_v\). The product of the latter is
a unit times \(z^{D_0}\) by
Corollary~\ref{cor:levelcurve370}. Squaring proves the
claim, with the same identities after truncation, even at
precisions at which this square vanishes.
\end{proof}

The construction gives the finite-level maps and their
trace relations for each auxiliary set, including on the
original \(L_a/F_a^+\) used to test the cohomology curves.
Passing from these maps to support of every required
mixed-characteristic component requires a further
commutative-algebra argument at the curve primes.


\subsection{The limiting Galois group and its ordinary cohomology}
\label{sec:liminertia407}

Return to the field \(F=F_a^+\), \(L=L_a\) on which the
simultaneous Frobenius tests of Theorem~\ref{thm:primes367}
were constructed. Fix one pair \(r,\alpha\). Denote its
Galois quotient from Section~\ref{sec:commonframe365} by
\(\widetilde\Gamma_0\), and allow additional inertia at
the \(q\) places of \(\mathcal Q_N\) to obtain
\(\widetilde\Gamma_N\). At all other places use exactly the
previous local conditions. The sign \(\eta\) cuts out \(L\).
The residual characters \(\bar\Psi,\bar\kappa_\ell,\eta\)
give a common finite quotient \(\Delta\) with fixed field
\(T\) in the sense that \(\Delta=\operatorname{Gal}(T/F)\).
In particular \(T\) contains \(L\) and \(\mu_\ell\).
The integer \(N\) indexes the auxiliary level and varies
with this field \(T\) fixed.

\begin{lemma}[A uniformly generated quotient for the residual problem]\leavevmode
\label{lem:relativepro407}
In \(\widetilde\Gamma_N\) replace the kernel of the map
to \(\Delta\) by its maximal pro-\(\ell\) quotient, including
for \(N=0\). Denote the resulting groups by \(\Gamma_N\)
and \(\Gamma_0\). This preserves the ordinary determinant
problems of Section~\ref{sec:typeaction405} and the curve
cocycles of Sections~\ref{sec:selmerlength368}
and~\ref{sec:selmerlift369}. These groups have generating
sets of bounded cardinality as \(N\) varies. The natural
map \(h_N:\Gamma_N\to\Gamma_0\) is onto, and its kernel is
the closed normal subgroup generated by one tame
\(\ell\)-inertia generator at each of the \(q\) new places.
\end{lemma}
\begin{proof}
The kernel of the maximal pro-\(\ell\) quotient is characteristic
in the subgroup being reduced, and so is normal in the full
group. To verify the assertion for determinants, first take
an Artinian local coefficient ring with finite residue field.
It is finite. A continuous trace with these values factors
through a finite group quotient, since it is locally constant
on a compact group. Explicitly there is an open normal subgroup
whose right multiplication preserves the trace at every point.
As \(3!\) is invertible, Newton's formula on arbitrary finite
group algebra sums shows that the determinant factors as well.
Use the Cayley--Hamilton algebra of this determinant over the
finite group quotient. This is a quotient of a finite group
algebra over the coefficient ring and is finite as a set.
The determinant descends to it, by the Cayley--Hamilton
ideal theorem for determinants. On the CM subgroup above
the identity of \(\Delta\), the residual characteristic
polynomial of each group element is \((X-1)^3\). Its image
in the algebra modulo the coefficient maximal ideal is thus
unipotent, by Cayley--Hamilton. Each such unit has
\(\ell\)-power order, by taking \(\ell\)-powers in
characteristic \(\ell\). The units congruent to one modulo
the coefficient maximal ideal also form an \(\ell\)-group,
as their successive congruence quotients are additive groups
in characteristic \(\ell\). The finite image of the whole
subgroup is therefore an \(\ell\)-group by Cauchy's theorem.
The determinant consequently factors through the asserted
relative quotient. The same conclusion holds under conjugation
by the other coset of the CM subgroup, since the kernel being
removed is characteristic. Inverse limits apply to complete
coefficient rings. The ordered characters are prescribed on
the same local absolute groups, with the same residual and
torsion restrictions. All their determinant relations are
preserved under this factorization. \cite{VIII-VI-chenevier364}

The finite curve representations and their coefficient modules
have the same property. On the residual kernel their semisimple
constituents in residual characteristic are trivial. This
follows first on the marked representation from its three
residual characters and then on its endomorphisms, direct
summands and the successive extensions used to form the
paired coefficients. A stable residual flag and its congruence
filtration put the image of this kernel in an \(\ell\)-group.
Adjoining a cocycle to such an action amounts to mapping into
a semidirect product by its finite additive module. The
image of the residual kernel is still an \(\ell\)-group.
Hence the finite-precision representations and cocycles
descend, also with the prescribed local conditions.
Inverse limits give the claim for integral curve coefficients.

Let \(S_T\) contain the places of \(T\) above the old primes
and \(\ell\). At level \(N\) add at most \(|\Delta|q\) more
places. Since \(T\) contains \(\mu_\ell\), Kummer theory
bounds the dimension of the continuous homomorphisms of the
kernel of \(\Gamma_N\to\Delta\) to \(\mathbf F_\ell\) by
\[
   r_1(T)+r_2(T)+|S_T|+|\Delta|q+
                         \dim_{\mathbf F_\ell}\Cl(T)[\ell].
\]
Indeed the Kummer classes have valuations divisible by
\(\ell\) outside this set. The quotient by the classes of
units with these allowed denominators injects in the
\(\ell\)-torsion of the corresponding class group, by taking
the \(\ell\)th root of the ideal away from the set. The
\(\ell\)-rank of that class group is at most that of
\(\Cl(T)\), since it is a quotient. Dirichlet's unit theorem,
including the roots of unity modulo \(\ell\)th powers,
bounds the remaining dimension by the first four terms.
Any extra conditions at old inertia only decrease the
dimension. In a pro-\(\ell\) group, lifts of a basis of its
quotient by powers and commutators generate topologically.
Indeed failure to generate is detected in a finite
\(\ell\)-group, where every maximal subgroup has index
\(\ell\). The displayed bound and lifts of generators of
\(\Delta\) give a uniform number of generators for
\(\Gamma_N\).

At a new place inertia has trivial image in \(\Delta\).
Its pro-\(\ell\) image is thus a quotient of tame
\(\mathbf Z_\ell(1)\), because the residue characteristic
differs from \(\ell\). Imposing triviality on one generator
there kills the whole inertia. Imposing this at all new
places before taking the relative quotient gives exactly
the original ramification problem. Conversely every quotient
of that original group which is pro-\(\ell\) over \(\Delta\)
is already a quotient of the auxiliary group. The two
quotient operations therefore give the claimed kernel
and surjectivity.
\end{proof}

\begin{theorem}[The kernel in a limit with marked inertia]\leavevmode
\label{thm:limitkernel407}
One can choose the auxiliary sets and an increasing sequence
of levels tending to infinity so that they have a limiting
profinite group \(\Gamma_\infty\), a surjection
\(h:\Gamma_\infty\to\Gamma_0\), and elements
\(\phi_j,\tau_j\) for \(1\le j\le q\). The group is finitely
generated and pro-\(\ell\) over \(\Delta\). The fixed local
absolute groups map to it compatibly over \(\Gamma_0\), as
does the chosen complex conjugation. For the elements
\(s_j=n_j'\) used in Theorem~\ref{thm:primes367}, viewed in
\(\Gamma_0\), one has
\[
 h(\phi_j)=s_j,\qquad h(\tau_j)=1,\qquad
 \phi_j\tau_j=\tau_j\phi_j,\qquad
 \ker h=\overline{\langle\!\langle\tau_1,\dots,\tau_q
                                           \rangle\!\rangle}.
\]
Here the normal closure is taken in the full group over \(F\).
The last equality identifies the entire unramified quotient
of the limit.
\end{theorem}
\begin{proof}
A finitely generated profinite group has only finitely many
open subgroups of each index, by its actions on finite sets.
Consequently the intersections of its open normal subgroups
of index at most \(m\) are a descending cofinal sequence.
In the Chebotarev construction at level \(N\) also prescribe
the images of the \(s_j\) in the \(N\)th such quotient of
\(\Gamma_0\). These prescriptions use the very same elements
of the cyclotomic kernel as the test on the curve. Thus they
are compatible with the required congruence
\(q_{N,j}\equiv1\bmod\ell^N\) and with all cocycle and
discriminant tests there. Chebotarev applied to the combined
finite quotient over \(L\) supplies the primes, as in
Theorem~\ref{thm:primes367}. The base images of the chosen
Frobenius elements now tend to \(s_j\).

We construct one presentation mapping to every level.
Choose \(d\) fixed generators of \(\Gamma_0\) and lifts of
them in each \(\Gamma_N\). The closed subgroup generated by
the lifts maps onto \(\Gamma_0\) by compactness. By the lemma
take a tuple of at most \(e\) generators for each \(\Gamma_N\),
for one \(e\). Multiply each member of that tuple by an
element of the subgroup of lifts to bring its image in
\(\Gamma_0\) to one. The corrected tuple together with the
lifts still generates. Let \(\Phi\) be the free profinite
group on \(d+e+q\) generators with one additional generator
\(c\) subject to \(c^2=1\). Map the first \(d+e\) generators
as above, the next \(q\) to the designated inertia generators
and \(c\) to the fixed complex conjugation. The maps
\(\pi_N:\Phi\to\Gamma_N\) are onto and their composites to
\(\Gamma_0\) are one fixed map \(\pi\). Write \(K_N=\ker\pi_N\),
\(K=\ker\pi\), and \(t_j\in\Phi\) for the designated inertia
generators. For every \(N\) the preceding lemma gives
\[
         K=\overline{\langle\!\langle K_N,t_1,\ldots,t_q
                                           \rangle\!\rangle}.
\]

Use the cofinal characteristic open subgroups \(V_m\) of
\(\Phi\) described above. In \(\Phi/V_m\) there are finitely
many possibilities for \(K_NV_m/V_m\). Once this subgroup
is fixed, the images of the \(q\) Frobenius elements in the
resulting quotient also take finitely many values. The
same is true of the maps from each fixed local absolute
group at \(\ell\) and the old places to that finite quotient.
A finite extension of a \(p\)-adic field has only finitely
many extensions of any bounded degree up to isomorphism,
by the local finiteness theorem (Krasner's lemma applied to
Eisenstein polynomials). Thus there are finitely many such
continuous homomorphisms, as each factors through one of
finitely many finite Galois groups. Fix compatible embeddings
of the local closures into the global one at the old places
throughout. By successively passing to infinite subsequences
and then choosing a diagonal subsequence, all the indicated
data are eventually constant for each \(m\).

The stabilized kernel images are compatible and surject onto
one another. Their inverse limit is a closed normal subgroup
\(K_\infty\subset\Phi\) with exactly those images. Put
\(\Gamma_\infty=\Phi/K_\infty\). Since all \(K_N\subset K\),
one has \(K_\infty\subset K\). Each finite congruence
quotient of the limit is obtained from a sufficiently late
\(\Gamma_N\). Taking the quotients also to remember the fixed
map to \(\Delta\) shows that its kernel in the limit is
pro-\(\ell\). The stabilized markings give \(\phi_j,\tau_j\)
and continuous maps on the fixed local groups by inverse
limits. Their composites to \(\Gamma_0\) are the original
markings, with the asserted Frobenius limits. In particular
\(\tau_j\) is the image of \(t_j\).

At finite level the tame relation for geometric Frobenius
is \(\phi_{N,j}^{-1}\tau_{N,j}\phi_{N,j}
=\tau_{N,j}^{q_{N,j}}\). In each fixed finite quotient the
orders of these inertia images are bounded powers of
\(\ell\), so the right side equals \(\tau_{N,j}\) for all
large \(N\). This proves commutation in the limit. Finally
reduce the displayed identity for \(K\) modulo \(V_m\).
The image of \(K_N\) eventually equals that of \(K_\infty\),
and the \(t_j\) are fixed. Thus \(K\) and the closed normal
subgroup generated by \(K_\infty,t_1,\ldots,t_q\) agree
in every \(\Phi/V_m\). Equality of the closed subgroups
proves the last assertion. The added finite Chebotarev
conditions and subsequences can be chosen together for
any finite list of pairs or curves already under consideration.
\end{proof}

\begin{proposition}[The unramified specialization of the limiting determinant]\leavevmode
\label{prop:unramifiedlimit407}
Let \(R_\infty^{\rm ord}\) denote the universal ordinary
polarized determinant ring on the CM subgroup of
\(\Gamma_\infty\), with its fixed local maps. Write
\(t_\infty\) for its trace. Take the closed ideal generated by
\[
    t_\infty(gu)-t_\infty(g),\qquad
    g\in\Gamma_\infty^0,\quad
    u\in\{\tau_j,c\tau_jc^{-1}:1\le j\le q\}.
\]
The quotient by this ideal is canonically the original
ordinary determinant ring \(R_0\), including its ordered
local characters.
\end{proposition}
\begin{proof}
The limiting group and its open subgroups are topologically
finitely generated, by the finite-index generator theorem.
They satisfy the finiteness condition for profinite
determinants. Together with the usual local character rings
and the closed ordinary, tame and polarization relations
this gives the displayed universal ring, as in
Section~\ref{sec:basesupport371}. Alternatively the following
argument identifies the representing functors on Artinian
coefficient rings directly. \cite{VIII-VI-chenevier364}

A trace satisfying the indicated identities also satisfies
them for the inverses, by substitution for \(g\). Cyclic
invariance of the trace gives the same identities for
conjugates by the CM subgroup. Iteration gives them for
products, and continuity gives them on the closed normal
subgroup they generate there. Since the CM subgroup has
index two with representative \(c\), the theorem identifies
this subgroup with the entire kernel of \(h\). The trace
therefore factors uniquely through the original CM group.
Newton's identities give the same factorization for the
determinant, tested on generic finite sums in its group
algebra. The ordinary local characters are still maps on
the fixed local absolute groups. The ordinary word and
polarization identities pass through the surjection, and
the maps at the old primes agree by construction. Conversely
the pullback of every pair on \(\Gamma_0\) obeys all these
identities. These operations are inverse and commute with
coefficient homomorphisms. This proves the ring assertion.
\end{proof}

\begin{theorem}[Ordinary cocycles on the limiting group]\leavevmode
\label{thm:limitselmer407}
On any selected curve let \(A=k'[[z]]\), let \(W,E,M,U\)
be the paired, diagonal, off-diagonal and extension modules
of Sections~\ref{sec:selmerlength368}
and~\ref{sec:selmerlift369}, and write \(b\) for the
rank of the ordinary off-diagonal cohomology over the base.
All coefficients on \(\Gamma_\infty\) are pulled back by \(h\).
Use the marked ordinary local conditions at the old places
and the coefficient prime. Then the following hold.
\begin{enumerate}
\item The \(A\)-rank of the integral polarized ordinary cocycles
with values in \(W\), modulo polarized coboundaries, is at most
\(3([F:\QQ]+q)\).
\item For \(0\le i\le b,\ 1\le j\le b\), let \(y_{ij}\)
be any ordinary polarized \(E\)-valued cocycles on the limit.
With the exponent \(C\) of Theorem~\ref{thm:selmerlift369}
there are \(c_0,\ldots,c_b\in A\), at least one a unit, and
one ordinary polarized \(U\)-valued cocycle \(Y\) on the limit
satisfying
\[
                     (\pi Y)_j=z^C\sum_{i=0}^{b}c_i y_{ij}
                                      \qquad(1\le j\le b).
\]
\item Pullback identifies the integral polarized ordinary
\(M\)-valued cocycles and their coboundaries with those on
the original group. In particular the chosen basis
\(\nu_1,\ldots,\nu_b\) on the base still spans cohomology
after inverting \(z\).
\end{enumerate}
\end{theorem}
\begin{proof}
Fix \(m\) and use the finite modules over \(A/(z^m)\). Pull
cocycles back to the fixed presentation \(\Phi\), or to its
index-two CM subgroup, which also has a finite generating
set. Only finitely many cocycles exist on this presentation
with each prescribed finite action. Indeed adjoining a
cocycle to the action gives a homomorphism to a finite
semidirect product, determined by the values on the
generators. All the candidate homomorphisms can be tested on
one finite quotient of \(\Phi\), taking the core of their
kernels when starting on the CM subgroup. Descent of a
candidate through \(K_N\) or \(K_\infty\) is detected by the
image of this kernel in the same finite quotient. The local
maps to that quotient stabilize as well. Projection to upper
triangular entries up to local coboundaries is a finite test,
including the possible local witnesses, as is the symmetry
under \(c\). Old inertia conditions are tested by the local
maps restricted to the fixed inertia groups. It follows
that at this precision the polarized ordinary cocycles for
the limit and for every sufficiently late level give
identical subsets on the presentation. The same holds for
their coboundaries and the coefficient projection \(\pi\)
between the modules. The finite quotient argument can be
applied to all the modules appearing here at once.

Put \(D=3([F:\QQ]+q)\). Corollary~\ref{cor:bfirst368} now
bounds the length of the limiting finite-precision polarized
ordinary cohomology by \(mD+C_1\). Suppose \(e\) integral
classes are independent over the fraction field. Evaluate
their representatives on a fixed finite generating tuple
of the CM presentation, in a finite free \(A\)-module
\(A^t\). Let \(B'\) be the image of the matrix taking
polarized witnesses to coboundaries on that tuple.
The allowable witnesses here are cut out by the involution
at \(c\) and form a direct summand since \(2\) is invertible.
Their reductions give all finite-precision coboundaries.
Modulo the saturation of \(B'\) the \(e\) vectors are
linearly independent over the fraction field, by the
cocycle identity and the fact that the tuple generates.
The quotient by that saturation is free. Smith normal
form shows that the reductions of the \(e\) vectors generate
a module there of length at least \(em-C_2\) for a constant
independent of \(m\). Dividing just by actual coboundaries
can only increase that length. Comparison as \(m\) tends
to infinity proves \(e\le D\).

For the second assertion reduce the given cocycles modulo
\(z^m\). The finite test just proved and
Theorem~\ref{thm:selmerlift369} provide the coefficients
and the desired lift at every \(m\), over the limiting
group. Keep all possible solutions at each precision,
including local witnesses. They form finite nonempty sets.
Reduction maps send a solution at higher precision to a
solution at lower precision, preserving primitivity of the
coefficient vector. A compatible sequence exists by compactness
of a product of finite sets. Indeed every finite set of
compatibility conditions is witnessed by reduction of a
solution at its largest precision. The sequence defines
the integral \(c_i\), the continuous cocycle \(Y\), and
its integral local witnesses. Its relation holds modulo
every power of \(z\), hence in \(A\).

Finally a cocycle with coefficients in \(M\) is a continuous
additive homomorphism on \(\ker h\), since the coefficient
action is trivial there. Commutation of \(\tau_j\) with
\(\phi_j\) gives
\((\operatorname{Ad}r(s_j)-1)y(\tau_j)=0\).
Each \(s_j\) has three distinct eigenvalues over a splitting
field of the curve fraction field. Thus the spectra in
different constituent blocks are disjoint. There is no
nonzero intertwiner between two different blocks, so the
operator in this equation has zero kernel on \(M\) over the
fraction field. The module \(M\) is torsion-free, and
\(y(\tau_j)=0\) integrally. Polarization gives the same at
\(c\tau_jc^{-1}\). The cocycle identity gives the corresponding
zero at all CM conjugates, and continuity gives zero on
the full kernel by Theorem~\ref{thm:limitkernel407}. Hence
\(y\) descends uniquely to the original group. Local
conditions and coboundaries agree on both sides by the
fixed local markings. If there is a single constituent,
the off-diagonal module is zero and the assertion holds
as well.
\end{proof}

The exponent and the rank bounds in this theorem are measured
along a curve in characteristic \(\ell\). To obtain the
mixed-characteristic length of Section~\ref{sec:missing}
one still has to identify which components of the original
ordinary determinant ring occur in the support of the
patched module and specialize at the prescribed weights.

\section{Hecke support and ordinary extension classes}\label{nav:chapter-ordinary-hecke}
\subsection{A patched module at the discriminant of the cohomology curves}
\label{sec:patchcomparison408}

The ranks of the integral ordinary modules at different
auxiliary levels need not have a common bound in the
construction so far. We use the maps in
Theorem~\ref{thm:finitelevel406} to obtain a finite module
after inverting the element which is nonzero at the chosen
curve. We keep the same prime sets, weight algebra and
residual ordering in this passage.

Let \(\Phi,\Gamma_\infty,\Gamma_0,h\) be as in
Theorem~\ref{thm:limitkernel407}. Enlarge the residue field
first if necessary to contain all residual roots. There
is a complete Noetherian local \(O\)-algebra \(R\) for the
following data: a degree-three continuous determinant on
\(\Phi^0\) with the given reduction, the three local characters
on each of the fixed ordinary local groups lifting the
specified residual characters and finite diamonds, and
variables \(U_{v,i}-b_{v,i},Y_{v,i}\) in the maximal ideal.
Here \(v\) in a subscript now denotes a position in the
ordered list of \(q\) auxiliary places. The units \(b_{v,i}\)
lift the designated residual roots of \(s_v\). The residual
roots at the finite levels agree with these after omitting
finitely many levels. To construct \(R\), take the
determinant ring of the finitely generated profinite group
\(\Phi^0\), the completed group rings for the local characters,
and their completed tensor product over \(O\), followed by
the indicated formal variables. The local unit groups have
finite \(\ell\)-rank. These rings are quotients of power
series rings with finitely many variables by the finiteness
condition for continuous determinants and local reciprocity.
Thus the same holds for their completed tensor product.
The ring admits the weight map \(\Lambda\to R\) by its
ordered inertial characters. \cite{VIII-VI-chenevier364}

Let \(\mathfrak n,\mathfrak m\) be the maximal ideals of
\(R,\Lambda\), and put \(\mathfrak a=(Y_{v,i})\subset R\).
The maps from \(\Phi\) to the finite-level groups give
actions of \(R\) on the residual summands
\(M_N=M_{1,N}^{\rm ord}\) of Proposition~\ref{prop:orderedpiece406}.
On these modules \(1+Y_{v,i}\) acts by its diamond and
\(U_{v,i}\) by its ordered Frobenius operator. Define in \(R\)
\[
        f=\prod_{v}\left(\prod_{i\ne j}(U_{v,i}-U_{v,j})\right)^2 .
\]
The finite-level composite \(f_N=\prod_v f_v\) of
Theorem~\ref{thm:finitelevel406}, written in the same variables,
satisfies \(f_N-f\in\ell^N R\). Indeed \(q_{N,v}\equiv1\bmod
\ell^N\) in every factor.

\begin{lemma}[The fixed comparison fiber]\leavevmode
\label{lem:fixedfiber408}
Let \(P\) be the selected local factor of the ordered
splitting algebra over \(R_0\) for the characteristic
polynomials at \(s_v\), and put \(X=P\otimes_{R_0}M_K\).
Give \(X\) the \(R\)-action from \(\Phi\to\Gamma_0\), these
ordered roots and \(Y=0\). It is finite over \(\Lambda\).
For fixed \(j\) and all sufficiently large \(N\), the
corresponding ordered hyperspecial module \(X_N\) at
the finite level is canonically \(X\) modulo \(\mathfrak m^j\)
as an \(R\)-module. The two comparison maps between this
module and \(M_N/(\mathfrak a,\mathfrak m^j)M_N\) have
composites \(f\). There is an integer \(d\), independent
of \(N\), such that \(\mathfrak n^d X_N\subset\mathfrak m X_N\).
\end{lemma}
\begin{proof}
Omitting finitely many good places leaves the same good
Hecke algebra on the base module. Indeed the remaining good
Frobenius conjugacy classes are dense, so the continuous
trace and its characteristic coefficients recover the
omitted spherical polynomials. The same residual factor
of this algebra and the same \(M_K\) thus serve all \(N\).
On \(M_K\) the polynomials at the finite Frobenius elements
tend to the polynomials at \(s_v\) by the convergence in the
fixed base group. Their coefficients therefore agree as
matrices modulo \(\mathfrak m^j\) for all large \(N\).
The presentations adjoining successive roots of monic
cubics and quadratics now agree modulo that power. The
isomorphism keeps the ordered root variables and commutes
with residual factor projection by uniqueness of idempotent
lifting. All traces of elements of \(\Phi^0\) agree on
the two base modules since the maps to \(\Gamma_0\) are
one fixed map. The local characters likewise agree. These
facts prove the asserted \(R\)-linearity, including on
power series by continuity. The resulting identifications
are compatible with reduction in \(j\).
The finite-level composites reduce to \(f\) for \(N\ge j\).

The modules \(X_N,X\) are summands of at most \(6^q\)
copies of \(M_K\) as weight modules. Their reductions
modulo \(\mathfrak m\) therefore have dimension at most
\(d=6^q\operatorname{rank}_\Lambda M_K\) over the residue
field. On the selected factor, \(R\) acts locally, with
maximal ideal acting nilpotently on this finite-dimensional
space. For example the action factors through the local
factor of the finite ordered-root algebra over the base
Hecke algebra. Each nonzero step in its maximal ideal
filtration has positive codimension by Nakayama's lemma.
This proves the uniform bound with the specified \(d\).
\end{proof}

\begin{proposition}[The successive diamond quotients]\leavevmode
\label{prop:partialpatch408}
Choose a nonprincipal ultrafilter on the subsequence of
Theorem~\ref{thm:limitkernel407}, and put, for \(k\ge1\),
\[
 V_{k,j}=\prod_{\mathcal U}
          M_N/(\mathfrak a^k M_N+\mathfrak m^j M_N),\qquad
 V_k=\varprojlim_j V_{k,j}.
\]
These are \(R\)-modules. One has \(V_1[1/f]\simeq X[1/f]\).
If \(h_i\) is the number of degree-\(i\) monomials in the
\(3q\) diamond variables, there are exact sequences
\[
             0\longrightarrow V_1^{\,h_i}
                \longrightarrow V_{i+1}\longrightarrow V_i
                \longrightarrow0\qquad (i\ge1).
\]
The first map multiplies lifts by their indicated monomials.
In addition, if \(j,k\) are fixed, then for all sufficiently
large \(N\)
\[
      f^k\mathfrak n^{djk}M_N
                \ \subset\ \mathfrak m^j M_N+\mathfrak a^k M_N .
\]
\end{proposition}
\begin{proof}
An ultraproduct identifies two sequences when they agree
on a set in the ultrafilter. It preserves exact sequences
of modules by choosing representatives and lifting at each
index where a preimage exists. A constant finite module
has ultraproduct equal to itself, since exactly one part
of every finite partition belongs to the ultrafilter.
At \(k=1\) the lemma gives maps in both directions between
\(X/\mathfrak m^j X\) and \(V_{1,j}\). The maps are the
ultraproducts of the actual level maps, using the fixed
finite presentation of the ordered roots modulo
\(\mathfrak m^j\). They commute with reduction, since
for any two precisions the finite-level identifications
and maps commute on a cofinite set of indices. Taking inverse
limits gives the two maps with composites \(f\). Here
\(X=\varprojlim_jX/\mathfrak m^jX\) by finiteness over the
complete weight algebra. Inverting \(f\) proves the first
assertion.

For fixed \(k,j\), and all large \(N\), the relations
\((1+Y_{v,i})^{\ell^{a_{N,v}}}-1\) in the diamond ring
belong to \((Y)^k+\mathfrak m^j\). Indeed the coefficient of
\(Y^t\) for \(1\le t<k\) is divisible by
\(\ell^{a_{N,v}-v_\ell(t)}\), by multiplying the binomial
coefficient by \(t\), and \(a_{N,v}\ge N\).
Hence \(M_N/(\mathfrak a^k,\mathfrak m^j)M_N\) is free over
\((\Lambda/\mathfrak m^j)[Y]/(Y)^k\) for all such \(N\),
by Theorem~\ref{thm:finitelevel406}. The top degree-\(i\)
part in the truncation at \(i+1\) is canonically a copy of
the degree-zero fiber for each monomial of degree \(i\).
This is checked on a free basis, and the map itself just
multiplies a lift of a degree-zero element by that monomial.
Taking ultraproducts proves the indicated exact sequence
at precision \(j\). Its left-hand modules map onto each
other as \(j\) decreases, by the corresponding surjections
at finite levels. Inverse limit therefore preserves the
surjection on the right. Explicitly, lift successively in
\(j\) and correct each mismatch in the left-hand module.
This gives the exact sequence for \(V_i\). The same proof
at any two truncations identifies every intermediate
graded term by multiplication of the corresponding monomials.

Finally \(\mathfrak n^{dj}X_N\subset\mathfrak m^j X_N\)
by iterating the lemma. Apply both comparison maps on
\(M_N/\mathfrak aM_N\), on which their composite is \(f_N\).
It follows that \(f\mathfrak n^{dj}M_N\subset
(\mathfrak m^j+\mathfrak a)M_N\) for \(N\ge j\).
Iterate this inclusion \(k\) times and use
\((\mathfrak m^j+\mathfrak a)^k\subset\mathfrak m^j+\mathfrak a^k\).
This proves the last estimate. Its power of \(f\) is
independent of the weight precision \(j\).
\end{proof}

\begin{theorem}[A finite patched module and its fiber]\leavevmode
\label{thm:modulepatch408}
Write \(\widehat R\) for the \(\mathfrak a\)-adic completion
of \(R[1/f]\). The module
\[
                  \mathcal P=\varprojlim_k V_k[1/f]
\]
is finite over \(\widehat R\) and has
\[
 \mathcal P/\mathfrak a^k\mathcal P\simeq V_k[1/f],
 \qquad
 \operatorname{gr}_{\mathfrak a}\mathcal P
        \simeq (X[1/f])[\overline Y_{v,i}].
\]
In particular its fiber at zero diamonds is \(X[1/f]\).
The diamond variables are a regular sequence on \(\mathcal P\)
when this fiber is nonzero. For each of the cohomology curves
of Theorem~\ref{thm:markedcurve366}, the fiber is nonzero
at the corresponding ordered generic point.
\end{theorem}
\begin{proof}
Each localized transition map in the proposition is onto.
The actions on these quotients, killed by \(\mathfrak a^k\),
give the action of the completed ring. The fiber \(X[1/f]\)
is finite over \(R[1/f]\), since \(X\) is finite over
\(\Lambda\) and the weights act through \(R\).
Lift finitely many of its generators compatibly up the
surjections. By the proposition, all subsequent graded
layers are generated by the corresponding monomials
times these same lifts. Every element of \(\mathcal P\)
can now be written as a linear combination of them with
coefficients in \(\widehat R\). Indeed at each stage correct
the error by coefficients homogeneous in the diamond
variables of the current degree, with coefficients in
\(R[1/f]\). The resulting series converge in the completion.
Starting at degree \(k\) for an element with zero image in
\(V_k[1/f]\) gives coefficients in \(\mathfrak a^k\widehat R\).
This proves finiteness and identifies the kernel at every
stage. The canonical monomial identifications then give the
asserted graded module. The notation on the right means a
polynomial module on \(X[1/f]\); multiplication in the
graded module is the indicated multiplication by monomials.

In this graded module each variable is injective after
quotient by the preceding ones. This yields regularity on
the complete module as follows. A nonzero element has a
first nonzero graded term. Multiplication by the first
variable cannot annihilate that term. More precisely the
order of its product increases by exactly one. It follows
that the graded module after quotient by that variable is
the quotient of the graded module by its initial variable:
an element of the principal image in degree at least \(k\)
has a preimage in degree at least \(k-1\). The quotient is
again complete and separated, as it is finite over the
Noetherian complete ring \(\widehat R\) with its defining
ideal. Induction proves the assertion for all variables.

For nonvanishing return to the dense cuspidal base changes
of Lemma~\ref{lem:bcpoints371}. At the old split places
they have the finite inertia of \(\alpha\oplus1\oplus\alpha^{-1}\)
and zero monodromy. They therefore contain the types of
Lemma~\ref{lem:oldidempotent405}, by the local parameters
of these cuspidal base changes and the compact-type criterion
there. At the other tame places they admit spherical vectors
by the same base change and descent. Thus the full ordinary
module \(M_K\) used here contains these systems. Its
annihilator in \(R_0\) vanishes at their evaluations and
therefore on the domain of their family by density. It
lies in the prime of every indicated curve by the
specialization in Theorem~\ref{thm:stableprime371}.
The finite module consequently has nonzero fiber at that
prime by Nakayama. Extending its residue field to split
the cubics and taking an order of the roots compatible
with the prescribed residues gives the nonzero ordered
fiber. Such orderings exist by integral reduction of the
roots over the curve DVR. Each cubic has distinct roots
generically by Theorem~\ref{thm:primes367}. Hence \(f\)
is nonzero there. Localizing by it preserves the fiber.
\end{proof}

\begin{theorem}[Galois relations on the patched module]\leavevmode
\label{thm:patchrelations408}
On \(\mathcal P\) the universal trace on \(\Phi^0\) descends
to \(\Gamma_\infty^0\) and satisfies the polarized ordinary
relations with its limiting local maps and the old tame
conditions. For \(g\in\Gamma_\infty^0\) it satisfies, at
each marked pair \(\phi_v,\tau_v\),
\[
 \begin{split}
 t\!\left(g\prod_i(\phi_v-U_{v,i})\right)&=0,\\
 t\!\left(g\prod_{j\ne i}(\phi_v-U_{v,j})
                         (\tau_v-1-Y_{v,i})\right)&=0 .
 \end{split}
\]
These are identities of scalar endomorphisms on the finite
\(\widehat R\)-module. The determinant and local characters
consequently give a pair over the quotient of \(\widehat R\)
by its annihilator on \(\mathcal P\), with the determinant
identities understood coefficientwise.
\end{theorem}
\begin{proof}
We first give a uniform way to pass an approximate relation
to this module. Suppose \(u\in R\) and there is \(e=0\) or
\(1\) such that for every \(t\), at all sufficiently large
levels there is \(u_N\in R\) annihilating \(M_N\) with
\(u_N\equiv f^e u\bmod\mathfrak n^t\). Fix \(k,j\) and take
\(t=djk\). The last inclusion in
Proposition~\ref{prop:partialpatch408} then shows that
\(f^{k+e}u\) kills \(V_{k,j}\). The exponent does not vary
with \(j\), so the same element kills \(V_k\). After inverting
\(f\), \(u\) kills \(V_k[1/f]\) for every \(k\) and therefore
kills \(\mathcal P\).

At a prescribed \(\mathfrak n^t\)-precision the universal
trace on the presentation factors through a finite quotient
of \(\Phi^0\). This follows from continuity, compactness
and finiteness of \(R/\mathfrak n^t\). Take its normal core
in \(\Phi\) and also remember the sign. If \(a\in K_\infty\),
choose \(a_N\in K_N\) with the same image in this finite
quotient, using Theorem~\ref{thm:limitkernel407}.
Then \(t(ga)-t(g)\) is congruent to \(t(ga_N)-t(g)\),
which acts as zero on \(M_N\). The preceding argument
descends the trace. Polarization is tested in the same way
or directly, using the common conjugation on \(\Phi\) and
the fixed cyclotomic character from \(\Gamma_0\).
For an ordered local word lift its three group elements to
\(\Phi^0\) using their specified images in the limit. At
every sufficiently large level the local maps have the same
images in the fixed finite quotient. There are thus lifts
of the actual finite-level local elements congruent to the
three chosen lifts there. The local character values in
the coefficients are those of the same elements of the
fixed local absolute groups. Expanding the word and taking
the trace after any global element makes its congruence with
the corresponding finite-level trace explicit. That trace
annihilates \(M_N\) by Theorem~\ref{thm:alllevels405}.
For an old inertia relation use the same argument with a
single local element in the specified kernel. This proves
the ordinary and old tame assertions.

For the two displayed words lift their marked Frobenius
and inertia to \(\Phi^0\). The finite-level markings can
again be lifted to agree modulo the chosen finite quotient.
The first trace then agrees to this precision with an
annihilator by Theorem~\ref{thm:inertiatrace406}. For the
second trace, multiply its finite-level version by the total
\(f_N\). This is an annihilator by the same theorem, and
\(f_N\equiv f\bmod\mathfrak n^t\) when \(N\ge t\). The
approximation argument with \(e=1\) now proves the second
identity. Conjugating the local words by the polarization
gives the conjugate conditions as well. Finally the
determinant coefficients are recovered by Newton's identities
on generic finite group algebra sums, so trace descent
gives descent of all the determinant identities. This
argument asserts algebraic identities on the localized,
diamond-adically completed ring. The original continuous
identities at the finite coefficient levels were used
before localization.
\end{proof}

The module constructed here has a nonzero fiber at each
tested curve, with all the added diamond directions regular.
Determining its support along every component passing through
that fiber is a further issue. The finite-level and limiting
identities above do not yet yield the mixed-characteristic
length, or the comparison with the normalized theta periods,
required in Section~\ref{sec:missing}.


\subsection{Full Hecke support at an absolutely irreducible curve}
\label{sec:irrcurvesupport409}

Take a curve from Theorem~\ref{thm:markedcurve366} whose
generic CM representation remains absolutely irreducible
on the auxiliary field \(L=L_a\). Put \(F=F_a^+\) and
\(d=[F:\QQ]\), and let \(q\) be the number of new places
in Theorem~\ref{thm:primes367}. Let \(\mathfrak p\) be the
prime it defines in the ordinary polarized determinant ring
\(R_0\) over this field. Our aim is to determine the ordinary
Hecke algebra near \(\mathfrak p\). Set \(D=3(d+q)\).

Take a finite extension of the curve fraction field which
contains the ordered roots at the \(s_j\) and over which
the generic representation is realized with a stable
integral lattice. Its valuation ring is a complete
equal-characteristic DVR \(A=k'[[z]]\), where \(k'\) is finite.
Use the embedding of this coefficient field in \(A\).
Extend the coefficient DVR \(O\) unramified to \(O'\) with
residue field \(k'\), enlarging \(k'\) first if necessary.
At every level adjoin one formal parameter \(s\) and extend
coefficients from \(O\) to \(O'\). At the curve send
\(s\) to \(z\). Write \(K=k'((z))\). The parameter supplies a
specified copy of this residue field when taking tangents
in residue characteristic \(\ell\).

Here and below this extension of the patching construction
means the following. At the base take \(B=(R_0\otimes_O O')[[s]]\).
Use \(\Lambda'=(\Lambda\otimes_O O')[[s]]\) as weight algebra.
Replace the modules and finite-level rings by the same
coefficient extension and formal power series extension.
They have the same finite freeness and comparison identities
by base change. On the Galois groups the residual data are
unchanged after extension of constants. On curve cohomology
the fraction-field rank estimate of
Theorem~\ref{thm:limitselmer407} is unchanged under the finite
extension of \(k'((z))\) just used. Indeed a continuous
cocycle with values in a finite extension splits in any field
basis into continuous cocycles in the original field.
The same applies to coboundaries and to the linear equations
for local witnesses. Their values have bounded denominators
on the compact groups, and multiplying by a power of a
uniformizer makes them integral. Thus extension of scalars
commutes with the cohomology relevant to that rank estimate.

Let \(P\) be the ordered splitting factor over \(B\), and
let \(\mathfrak x\) be its prime at the extended curve,
with an ordering compatible with the residual one.
Repeat Section~\ref{sec:patchcomparison408} using \(\Lambda'\).
The finite truncations of \(\Lambda'\) are still finite, the
modules remain free over the indicated diamond group rings,
and the convergence of the base Frobenius coefficients
continues to hold at each finite weight precision. These
are the inputs to Propositions~\ref{prop:partialpatch408}
and Theorem~\ref{thm:modulepatch408}; their proofs therefore
give the module \(\mathcal P\) with augmented base fiber
\((P\otimes_{R_0}M_K)[1/f]\). The presentation ring there
is now augmented by \(s\); denote its localized
diamond-adic completion by \(\widehat R'\).
Let \(T\) be its quotient by the trace descent, polarization,
old inertia, ordered local character-word, and two auxiliary
trace relations in Theorem~\ref{thm:patchrelations408}.
All are imposed with arbitrary global multipliers as there.
Use the full coefficient identities for polarization. The
module \(\mathcal P\) is a finite \(T\)-module.

There is a surjection \(T\to P[1/f]\) which sends all
diamond variables \(Y_{v,i}\) to zero. Indeed the original
ordinary determinant pulled back to the presentation, its
local characters and the ordered universal roots give this
map. At the base, the image of \(\tau_v\) is one and the
first Frobenius trace relation is Cayley--Hamilton. Thus
the relations defining \(T\) hold on the entire ordered
base. The universal traces and local character values
generate \(R_0\) topologically, since all group elements
have lifts to the presentation. This is also immediate
from the quotient construction of the ordinary determinant
ring. The root coordinates generate the ordered factor
over the augmented base. Before localization the image of
the complete presentation ring is closed by compactness,
hence contains all these generated rings. Inverting \(f\)
and setting \(Y=0\) during completion preserves surjectivity.
Let \(\mathfrak P\subset T\) be the preimage of \(\mathfrak x\)
in this map, using the same symbol for its extension after
inverting \(f\). Its residue field is \(K\). All evaluations
before localization lie in \(A\), and the image contains
\(k'[[z]]\) through \(O'[[s]]\).

\begin{lemma}[Depth and relative tangent directions at the curve]\leavevmode
\label{lem:parameterdepth409}
The nonzero finite module \(\mathcal P_{\mathfrak P}\) has
depth at least \(D+1\) over \(T_{\mathfrak P}\). Let \(e\)
be the embedding dimension of this local ring, let
\(\varpi\) be a uniformizer of \(O'\), and let \(\mathfrak p_T\)
be its maximal ideal. The \(K\)-space of lifts of its residue
map to \(K[\epsilon]/(\epsilon^2)\) fixing \(O'[[s]]\)
in the given residue map has dimension
\[
              \dim_K\mathfrak p_T/
                        (\mathfrak p_T^2+\varpi T_{\mathfrak P})
                        \ \ge e-1 .
\]
On the original presentation ring, the first-order parts
of all these lifts are continuous on group traces and have
uniformly bounded denominators in \(A\).
\end{lemma}
\begin{proof}
The \(3q\) diamond variables already form a regular sequence
on the patched module by Theorem~\ref{thm:modulepatch408}.
Its fiber localized at \(\mathfrak P\) is nonzero by that same
theorem and the choice of \(\mathfrak x\). Write the remaining
\(3d\) principal weight variables as \(w_i\), and let
\(\widetilde a_i(s)\in O'[[s]]\) lift their power series values
in \(A\), with zero constant term modulo \(\varpi\). Translation
by these series in the weight algebra shows that
\(\varpi,w_i-\widetilde a_i(s)\) form a regular sequence
of length \(3d+1\) there. The ordered base module is a finite
free \(\Lambda'\)-module by the monic splitting presentation
and its residual direct summand, as in
Lemma~\ref{lem:fixedfiber408}. The same sequence is regular
on it. Inverting \(f\) and localizing the base module at
\(\mathfrak P\) preserve injectivity at each step. The
final quotient is nonzero by the nonzero fiber over \(K\).
Together the two sequences give the stated depth.

Modulo \(\varpi\), the map from \(k'[[s]]\) into the local
ring is injective, since evaluation sends \(s\) to \(z\).
Every nonzero element of this subring is inverted at the
curve prime. Thus \(K\) is a coefficient field of
\(T_{\mathfrak P}/\varpi\) splitting its residue map. Over
this field lifts to the dual numbers are exactly the linear
functionals on the displayed cotangent space, by taking
linear terms. Killing one element can decrease the embedding
dimension by at most one.

For the continuity assertion present the ring before the two
operations of localization and diamond completion as a quotient
of \(O'[[s,x_1,\ldots,x_n]]\). Its evaluation sends \(x_i\)
to series in \(zA\), which have lifts \(b_i(s)\in O'[[s]]\)
in the maximal ideal. Any given dual-number lift has values
\(x_i\mapsto b_i(z)+\epsilon c_i\), with \(c_i\in K\).
In this expression \(b_i(z)\) denotes the reduction evaluated
in \(A\). Translation by \(b_i(s)\) is an automorphism of
the ambient complete local power series ring. Modulo the
square of the ideal generated by the translated variables,
the value of every series is consequently its constant
term plus its linear terms in those variables. This proves
the first-order evaluation formula even for an algebraic
map from the localization. Its denominators on the entire
original ring are bounded by those of the finitely many
\(c_i\). On the \(n'\)th power of the original maximal ideal,
the linear terms have valuation at least
\(n'-1+\min(0,\min_i v_z(c_i))\). This proves continuity.
The diamond variables map into \(\epsilon K\), so such lifts
on the completion are determined modulo the square of the
diamond ideal. Completion induces an isomorphism on that
quotient of the localized ring. Thus no further choice of a
lift is introduced by completion.
\end{proof}

\begin{proposition}[Tangents detected by ordinary adjoint cohomology]\leavevmode
\label{prop:tangentcurve409}
The relative tangent space of the preceding lemma injects
linearly into the polarized ordinary cohomology of
\(\Gamma_\infty\) with coefficients in
\(\operatorname{End}_K(\rho)\). Here \(\rho\) is the pullback
of the generic representation on the curve. Consequently
this tangent space has dimension at most \(D\).
\end{proposition}
\begin{proof}
A lift gives a degree-three determinant over the dual numbers.
Its reduction is the determinant of the split absolutely
irreducible \(\rho\) over \(K\). Chenevier's Theorem 2.22(i)
over a Henselian local ring, as used in
Lemma~\ref{lem:detangent372}, supplies a matrix representation
lifting \(\rho\), unique up to conjugation by a matrix
reducing to one. The trace descent relations ensure that it
factors through \(\Gamma_\infty^0\). More explicitly, lift
nine group elements whose matrices span the full matrix
algebra over \(K\). Their trace Gram matrix is invertible
also over the dual numbers. Hence traces against these nine
detect zero matrices, including the difference from the
identity at any element of the descent kernel. The
coordinates of an arbitrary group matrix in this basis
are given by the inverse Gram matrix applied to nine
translated traces. The lemma makes their first-order
coefficients continuous and uniformly bounded in
denominator. A fixed integral lattice for \(\rho\) then
gives the same properties for the cocycle \(z(g)\) defined
by \(\rho_\epsilon(g)=(1+\epsilon z(g))\rho(g)\).
Strict conjugation changes \(z\) by a coboundary. The trace formulas apply over the curve field since \(3!\) is a unit here.
\cite{VIII-VI-chenevier364}

Polarization of the determinant makes the class invariant
under conjugation by \(c\) followed by negative adjunction
for the polarized base representation, with its fixed
multiplier. Indeed the conjugate-dual lift has the same
determinant and differs by strict conjugation after
identifying the base representations by their polarization.
Averaging a representative with its image under this
involution, using \(2\in K^\times\), gives a polarized
cocycle. Polarized representatives differing by a
coboundary differ by one with a polarized witness by the
same averaging.

At each chosen ordinary place take an inertia element
\(\gamma\) whose three curve characters have distinct
values. The ordinary word for three copies of \(\gamma\)
has zero trace against every group element, so is the zero
matrix by the Gram argument. Its three lifted roots are
unit-separated in \(K[\epsilon]\). Their polynomial
projectors give three free eigenlines reducing to the
original ones. In this basis put any local \(g\) in position
\(k\) of the word, leaving the other two letters equal to
\(\gamma\). On a row of index at least \(k\) the factors
to the left have invertible product, and on a column of
index at most \(k\) the factors on the right have invertible
product. It follows that the corresponding entries of
\(\rho_\epsilon(g)-\psi_k(g)\mathrm{id}\) vanish. Varying
\(k\) proves the ordinary flag with its specified diagonal.
After choosing a lifting basis this gives the ordinary
cocycle condition up to a local coboundary. The old tame
kernel is killed by the same trace pairing. Changing the
global representative by averaging preserves the local
conditions on its class. There are finitely many local
witnesses, so after multiplication by a power of \(z\)
they and the cocycle all take values in the integral
lattices. With one constituent the coefficient \(W\) of
Theorem~\ref{thm:limitselmer407} is the full adjoint
endomorphism space. Part 1 of that theorem therefore
bounds the dimension of the target by \(D\).

It remains to see that zero class forces the original
tangent to vanish. After conjugation the lift of \(\rho\)
is then constant, so its determinant is constant. On the
ordinary flag the eigenvalues of \(\gamma\) are uniquely
determined lifts of simple roots of its characteristic
polynomial. Each local character is recovered by taking
the trace of the corresponding eigenprojector times the
group matrix, so these too are constant. At every new
place the Frobenius has distinct eigenvalues. The first
auxiliary trace word, by the Gram pairing, gives
\(\prod_i(\rho(\phi_v)-U_{v,i})=0\). Evaluating on each
eigenline shows that its matching root coordinate is
constant, since each other factor is invertible. The
second word now gives
\(\prod_{j\ne i}(\rho(\phi_v)-U_{v,j})
(\rho(\tau_v)-1-Y_{v,i})=0\). The curve representation
comes from the unramified base, so \(\rho(\tau_v)=1\).
Evaluation on the \(i\)th eigenline gives \(Y_{v,i}=0\).
All generating data of the presentation ring and the
parameter \(s\) are thus fixed. Equality on its topological
generators gives equality on the ring by the continuity
in the lemma, then on the localization and on the
first-order image of the completion. This proves injectivity.
Adding first-order matrix variations adds their cocycles,
and changing basis contributes a coboundary. Thus the map
to cohomology is linear. Equivalently the construction commutes with
linear maps from a square-zero coefficient module, including
addition on two copies of \(K\).
\end{proof}

\begin{theorem}[All components through an absolutely irreducible curve]\leavevmode
\label{thm:irrcurvesupport409}
At the prime \(\mathfrak p\) of every selected curve with
absolutely irreducible generic representation, the localized
ordinary module \((M_K)_{\mathfrak p}\) is finite free of
positive rank over \((R_0)_{\mathfrak p}\). In particular
the natural surjection onto the localized good Hecke algebra
is an isomorphism
\[
                   (R_0)_{\mathfrak p}
                        \ \simeq\ (B_K)_{\bar{\mathfrak p}},
\]
where \(\bar{\mathfrak p}\) is the corresponding Hecke prime.
Every prime of \(R_0\) contained in \(\mathfrak p\) belongs
to the support of this ordinary module.
\end{theorem}
\begin{proof}
The two preceding results bound the embedding dimension of
\(T_{\mathfrak P}\) above by \(D+1\) and the depth of its
nonzero finite module below by \(D+1\). For a Noetherian
local ring the depth of such a module is at most the
dimension of the ring, which is at most its embedding
dimension. Hence both dimensions equal \(D+1\) and
\(T_{\mathfrak P}\) is regular local. Every finite module
over a regular local ring has finite projective dimension.
The Auslander--Buchsbaum formula, using the depth equality,
gives projective dimension zero for
\(\mathcal P_{\mathfrak P}\). It is therefore free of
positive rank. \href{https://stacks.math.columbia.edu/tag/065U}
{The Stacks Project, \emph{Regular rings and global dimension}.}
\href{https://stacks.math.columbia.edu/tag/090U}
{The Stacks Project, \emph{Auslander--Buchsbaum}.}

Tensor with \(P_{\mathfrak x}\) through the displayed
surjection from \(T\). The resulting module is
\(P_{\mathfrak x}\otimes_{R_0}M_K\). Indeed the diamond
variables are zero in \(P\), and the fiber identification
of Theorem~\ref{thm:modulepatch408} already gives
\((P\otimes_{R_0}M_K)[1/f]\) with the action through the
surjective map to \(P[1/f]\). This ordered module localized
at \(\mathfrak x\) is consequently free of positive rank.
The map \((R_0)_{\mathfrak p}\to P_{\mathfrak x}\) is flat
and local, hence faithfully flat. The coefficient extension
is finite free, adjoining \(s\) is flat over the Noetherian
base, the ordered splitting factor is finite projective
by the monic presentation, and localization preserves
flatness. Contraction of \(\mathfrak x\) gives precisely
the curve prime on \(R_0\). Faithfully flat descent of
flatness for the finitely presented module \((M_K)_{\mathfrak p}\)
shows that it is finite flat and therefore free. Its rank
is positive by the ordered calculation upstairs.

The good Hecke algebra \(B_K\) is the image of \(R_0\) in
the endomorphisms of \(M_K\). Its kernel is the annihilator.
The localized positive-rank free module is faithful, so
this kernel localizes to zero, proving the isomorphism.
For every smaller prime a further localization keeps this
module free of the same positive rank. This proves the
assertion about whole components. It gives the inclusion
of ideals in the first, irreducible case of
\srcloc{Development/AutomorphicArithmetic/AutomorphicVarieties/WeilGroupCharacters/UniversalOrdinaryDeterminantRing}{284}
for the curves and fields specified here.
\end{proof}

The other generic possibility in Theorem~\ref{thm:rescurve364}
is the sum of a character and an irreducible two-dimensional
representation. At such a point trace against the group
matrices kills the off-diagonal blocks, so the nine-matrix
pairing in Proposition~\ref{prop:tangentcurve409} is
degenerate. Theorem~\ref{thm:limitselmer407} supplies
information about the corresponding extension classes.
Their effect on component support, and the specialization
to the original integral cohomology lengths, remain to be
established.


\subsection{The off-diagonal algebra at a curve with two constituents}
\label{sec:oddblock410}

Let the representation on a selected curve over \(L=L_a\) be
\(r=r_1\oplus r_2\), where \(r_1\) is a character and \(r_2\)
is an absolutely irreducible representation of dimension two.
Both are individually polarized, and they remain irreducible
and inequivalent on the cyclotomic kernel by
Theorem~\ref{thm:markedcurve366}. Let \(K=k'((z))\) be the
curve field. We may first extend it by a finite extension
to put the symmetric polarization in each block into
orthonormal coordinates. Indeed a nondegenerate symmetric
matrix in characteristic different from two can be diagonalized,
and adjoining square roots of its nonzero diagonal coefficients
suffices. Use a uniformizer of the extended field for \(z\),
its residue field for \(k'\), and repeat the construction
with \(s\mapsto z\) in Section~\ref{sec:irrcurvesupport409}.
The stable integral lattice used for continuity can be transported
through this change of basis. Put
\[
            H=\widehat{T_{\mathfrak P}},\qquad
            \mathfrak m=\mathfrak m_H,\qquad H/\mathfrak m=K .
\]
Here \(T\) is defined by determinant relations before passage
to the annihilator of the patched module. These same relations
and the surjection onto \(P[1/f]\) apply to the present curve.
We study the Galois algebra over this complete local ring.

\begin{lemma}[The finite ordinary Cayley--Hamilton algebra]\leavevmode
\label{lem:finitech410}
There is a finite \(H\)-algebra \(C\) generated as an
\(H\)-module by group elements of \(\Gamma_\infty^0\).
Its determinant is the base change of the given determinant
on the presentation. The old inertia identities, the ordinary
matrix words and the two auxiliary matrix words of
Theorem~\ref{thm:patchrelations408} hold in \(C\), as do
their images under the adjoint involution. The semisimple
quotient of \(C/\mathfrak m C\) is
\(K\times M_2(K)\), with its transposition involution.
\end{lemma}
\begin{proof}
Start over the complete presentation ring, before localization,
with the universal continuous Cayley--Hamilton algebra on
\(\Phi^0\) pulled back along its determinant. This algebra is
finite over the presentation ring by the finiteness theorem
of Wang-Erickson. Its hypothesis on residual first cohomology
holds since values on the finite set of profinite generators
determine a cocycle into the finite residual adjoint module.
The images of the group span the base changed algebra.
Indeed they have dense linear span before base change in
the quotient of the completed group algebra. A submodule
of a finite module over a complete Noetherian local ring
is closed in its adic topology, by Artin--Rees. Ordinary
tensor product then preserves spanning. The topology at the
initial complete ring makes the group map continuous.
\href{https://people.brandeis.edu/~cwe/pdfs/algfam.pdf}
{C. Wang-Erickson, \emph{Algebraic families of Galois representations},
preprint, Propositions 2.2.2 and 3.2.2.}

Extend scalars to \(H\). Quotient by the two-sided ideal
generated by the differences \(g-1\) for the kernel of
\(\Phi^0\to\Gamma_\infty^0\), the old inertia differences,
the ordinary words and the two auxiliary words. Include also
the adjoints of the generators. On the group algebra adjunction
sends \(g\) to \(\mu(g)\sigma(g)^{-1}\), where \(\sigma\)
is conjugation by complex conjugation and \(\mu\) is the
fixed multiplier in the polarization. It fixes all scalars.
Polarization gives \(t(x^*)=t(x)\) on group algebra sums.
Together with cyclicity this gives the same characteristic
polynomial for \(x^*\) as for \(x\), by Newton's formulas.
Consequently the Cayley--Hamilton ideal is invariant under
adjunction. To see this also with the completed group algebra,
one can first impose the polarization equations on the
universal determinant over the complete coefficient ring.
On the completed group algebra adjunction is continuous
and the polarized identities extend from the dense group
algebra. Pullback to \(H\) gives the stated involution.

The trace against every group element of each imposed word
is zero by definition of \(T\). Cyclicity, polarization and
linear spanning give the same vanishing on the two-sided
ideal, even after arbitrary scalar extension. Applying
Newton's identities to the first three powers shows that
the determinant descends to the quotient \(C\), since
\(3!\) is invertible. It is Cayley--Hamilton there. The
semisimple curve representation realizes the specialization
of this determinant and satisfies the matrix words. At
the auxiliary places it comes from the unramified base
and its Frobenius has the ordered roots. Its image spans
\(K\times M_2(K)\), by the irreducibility and inequivalence
on the cyclotomic kernel. The kernel of this surjection in
\(C/\mathfrak m C\) is the determinant kernel. Indeed the
trace pairing on the displayed product is nondegenerate,
and the determinant is computed from its traces. This
kernel is nilpotent by the Cayley--Hamilton radical theorem
in residue characteristic greater than three. It follows
that it is the Jacobson radical. On the quotient the
involution is transposition by our choice of basis.
\cite{VIII-VI-chenevier364}
\end{proof}

\begin{proposition}[A single module for both off-diagonal blocks]\leavevmode
\label{prop:orthogonalblocks410}
There is a finite \(H\)-module \(N\) and a symmetric
\(H\)-bilinear form \(p:N\times N\to\mathfrak m\) satisfying
\[
                p(x,y)w=x\,p(y,w) \qquad (x,y,w\in N).
\]
Give \(H\oplus N\) the commutative algebra structure in which
\(x y=p(x,y)\) for \(x,y\in N\). With suitable matrix units,
\(C\) identifies, as an algebra with involution, with
\[
      \begin{pmatrix}
         H & M_{1,2}(N)\\
         M_{2,1}(N) & M_2(H)
      \end{pmatrix}\ \subset M_3(H\oplus N)
\]
with ordinary transposition. Its determinant agrees with
the matrix determinant after extension to \(H\oplus N\).
\end{proposition}
\begin{proof}
Let \(J_0\) be the kernel of the map from \(C\) to its
semisimple residue algebra. A power of \(J_0\) is contained
in \(\mathfrak m C\) by the lemma. The \(J_0\)-adic and
\(\mathfrak m\)-adic topologies on \(C\) coincide, and \(C\)
is complete in both. Lift the three primitive diagonal
idempotents to orthogonal self-adjoint idempotents summing
to one. For the first lift, replace any preimage by its
average with its adjoint. Newton iteration for \(X^2-X\)
converges because its error starts in \(J_0\) and its
derivative \(2X-1\) is a unit modulo \(J_0\). All iterates
commute with their rational expressions and remain
self-adjoint. Repeat inside the complementary corner and
take the residual complement for the last idempotent.

For each lifted idempotent \(e\), the determinant on the
corner \(eCe\) is Cayley--Hamilton of degree one. In fact its
degree can be read off from \(D(1-e+Xe)\) after reduction,
where the residue idempotent has rank one. Chenevier's
corner lemma gives Cayley--Hamilton on that corner.
The degree-one identity now says \(y=t(y)e\) for every
\(y\in eCe\), and \(t(e)=1\). Thus \(eCe=He\) with the
indicated copy of \(H\) embedded. \cite{VIII-VI-chenevier364}

Write \(e_1,e_2\) for distinguished primitive idempotents
in the character and the plane, respectively. For the other
idempotent \(e_j\) in the plane lift a standard matrix unit
to \(w\in e_j C e_2\). Then \(w^* w=d e_2\) for an element
\(d\equiv1\pmod{\mathfrak m}\). Divide \(w\) by the square
root of \(d\) congruent to one, which exists by Hensel's
lemma. Now \(w^* w=e_2\) and \(ww^*\) is an idempotent
of \(H e_j\) reducing to \(e_j\), so \(ww^*=e_j\).
Multiplication by these matrix units identifies the corners
within the plane and exchanges the off-diagonal corners
by transposition. Take \(N=e_1 C e_2\). For \(x,y\in N\)
define \(xy^*=p(x,y)e_1\). Its reduction on the semisimple
quotient is zero. Adjunction gives symmetry since it fixes
\(H e_1\). Cyclicity of trace and the rank-one corner
calculation give \(y^*x=p(x,y)e_2\). Computing \(xy^*w\)
by the two associations gives the asserted identity. It
also proves associativity of the multiplication on \(H\oplus N\),
by checking products of three elements of \(N\).
Decomposing \(C\) by the three orthogonal idempotents now
gives exactly the displayed algebra. Terms in distinct
corners have trace zero off the diagonal, by cyclicity.
Thus the matrix trace is \(t\) on every element and after
base change. Newton's identities give equality of the
determinants. Compare the polarized block structure in
\srcloc{Development/AutomorphicArithmetic/AutomorphicVarieties/PatchingModules/PatchedSupportCompletions}{78}.
\end{proof}

\begin{theorem}[Irreducible and nonsemisimple fibers at the curve]\leavevmode
\label{thm:oddsupport410}
Put \(J=(p(N,N))\subset H\) and \(I=\operatorname{Ann}_H N\).
For a prime \(Q\subset H\), use an algebraic closure of
\(\kappa(Q)\) for its geometric fiber. The determinant there
is absolutely irreducible precisely when \(J\not\subset Q\).
When \(J\subset Q\), the determinant has constituents of
dimensions one and two. The ordinary Cayley--Hamilton
fiber is then nonsemisimple precisely when \(I\subset Q\).
Consequently \(V(I)\) consists exactly of the irreducible
fibers together with the reducible nonsemisimple fibers.
\end{theorem}
\begin{proof}
Extend the block description to the geometric residue
field \(k\). If the pairing has nonzero image, there is an
\(x\in N\otimes_H k\) with \(p(x,x)\ne0\), since two is
invertible. The identity in the proposition gives
\(p(x,x)y=p(x,y)x\) for every \(y\) in this space. Thus it
is a line. Rescale \(x\) over \(k\) so that \(p(x,x)=1\).
The map \(k\oplus kx\to k\) sending \(x\) to one identifies
the displayed block algebra with all of \(M_3(k)\) on its
matrix entries. Its group elements span this algebra,
and the resulting representation is irreducible and realizes
the determinant.

If the pairing vanishes, the two off-diagonal blocks form
an ideal of square zero and the quotient is \(k\times M_2(k)\).
It is the full radical. The product determinant on this
quotient is the given determinant, since its trace agrees
on all group algebra sums. Each simple factor occurs with
multiplicity one. The radical is nonzero precisely when
\(N\otimes_H k\ne0\). As \(N\) is finite this is equivalent
to \(N_Q\ne0\), by Nakayama, and to \(I\subset Q\). In the
first case as well \(N_Q\ne0\) by the nonzero pairing.
These alternatives give the asserted locus.
\end{proof}

\begin{corollary}[The number of off-diagonal generators]\leavevmode
\label{cor:selmergenerators410}
Let \(b\) be the rank of ordinary polarized off-diagonal
cohomology on the original group of this curve as in
Theorem~\ref{thm:limitselmer407}. Then
\[
               \dim_K N/\mathfrak m N=b,\qquad b\ge1 .
\]
In particular \(J\) is generated by at most \(b(b+1)/2\)
elements. If \(b=1\), there is a \(d\in\mathfrak m\)
such that \(J=(d)\) and
\[
             H\oplus N\ \simeq\ H[X]/(IX,\ X^2-d),
             \qquad I d=0 .
\]
\end{corollary}
\begin{proof}
Put \(U=N/\mathfrak m N\). At this fiber the pairing is
zero. For \(\lambda\in U^\vee\), send every off-diagonal
entry \(u\) in the block description to \(\epsilon\lambda(u)\)
and each scalar entry to its residue, where \(\epsilon^2=0\).
This yields a polarized lift of \(r\). In the chosen coordinates
its first-order part is a cocycle with only off-diagonal
entries. It is continuous with bounded denominators over the
curve DVR. Indeed the group map comes from the finite module
over the original complete presentation ring in the lemma.
At a fixed finite set of module generators its specialized
coordinates have bounded denominators. Powers of the original
maximal ideal evaluate into the corresponding powers of
\(z\), which gives continuity for the finite module map.

The matrix words in \(C\) impose the ordinary condition on
this cocycle. At an ordinary place diagonalize the matrix of
the element with three distinct inertia eigenvalues over the
dual numbers using its simple-root projectors. The ordinary
word with two occurrences of that element and a local element
in the remaining position then gives upper triangularity with
the prescribed diagonal, by the row and column calculation
of Proposition~\ref{prop:tangentcurve409}. The change of
the flag from its original one is an infinitesimal local
conjugation. Projecting its coboundary witness onto the
off-diagonal blocks preserves the condition, since the
original ordinary basis respects the two constituents and
block projection preserves upper triangularity. The old
inertia differences vanish in \(C\). Multiplication by a
power of \(z\) puts the cocycle and all its local witnesses
in the chosen lattices. Part 3 of
Theorem~\ref{thm:limitselmer407} descends the resulting
ordinary class to the original group after inverting \(z\).
We have thus obtained a linear map \(U^\vee\) to a space
of dimension \(b\).

This map is injective. If the cocycle is a coboundary, its
lift on the group is conjugate to the constant representation.
The group spans \(C/\mathfrak m C\) linearly over \(K\),
so the two algebra homomorphisms also agree up to that
conjugation. The constant one kills the radical, hence the
lift kills every off-diagonal entry \(\epsilon\lambda(u)\).
This gives \(\lambda=0\).

For the opposite bound take the \(b\) ordinary polarized
classes on the base with integral representatives
\(\nu_1,\ldots,\nu_b\) chosen in
Section~\ref{sec:pairedtests367}. Write \(E\) for the
block diagonal algebra over \(K\) and \(M\) for both
off-diagonal matrix blocks. The representation
\[
                 g\longmapsto
                 \left(1+\sum_{j=1}^b e_j\nu_j(g)\right)r(g),
                 \qquad e_i e_j=0 ,
\]
inflated to the limit, gives an algebra map from
\(C/\mathfrak m C\) to \(E\oplus M^{\oplus b}\). We verify
factorization. All mixed terms of the determinant having
two off-diagonal entries vanish by the square-zero relation.
A single such entry has coefficient zero at every block
diagonal matrix by expansion of its determinant. Thus the
determinant is the original one on all algebra sums and
after scalar extension. The matrix Cayley--Hamilton identities
follow. Continuity and a stable integral lattice for these
matrices give a representation of the original completed
group algebra. The chosen local witnesses give upper
triangular matrices with the same diagonal at each ordinary
place, and polarization gives the adjoint words. Old inertia
conditions hold by construction. At the new places inertia
acts trivially and Frobenius has the original characteristic
polynomial, giving both auxiliary relations. This proves the
factorization, including through the residue map \(H\to K\).
Lemma~\ref{lem:spandual367} says the group already spans
the entire target algebra on the cyclotomic kernel. The
surjective map just obtained gives
\(5+4\dim_K U\ge 5+4b\) by the block dimensions. Together
with the injection this proves equality.

To prove positivity use the stable prime of
Theorem~\ref{thm:stableprime371} below the curve in \(R_0\).
The flat local map \((R_0)_{\mathfrak p}\to P_{\mathfrak x}\)
from the proof of Theorem~\ref{thm:irrcurvesupport409}
is independent of the number of constituents. Choose a
prime over the stable prime there by faithful flatness.
Pull it back by the surjection \(T_{\mathfrak P}\to P_{\mathfrak x}\),
then lift it to \(H\) by faithful flatness of completion.
The determinant at this prime of \(H\) is absolutely
irreducible, since it extends the determinant at the stable
prime. The theorem implies \(J\ne0\), hence \(N\ne0\) and
\(b\ge1\) by Nakayama. Finally choose \(b\) generators of
\(N\). Their pairwise products, with equal pairs allowed,
generate \(J\) by symmetry and bilinearity. For one generator
\(n\), put \(d=p(n,n)\). The map \(X\mapsto n\) gives
the asserted presentation. A polynomial reduces to \(a+cX\)
using \(X^2=d\), and this expression maps to zero exactly
when \(a=0\) and \(c\in I\), by the direct sum decomposition.
Also \(Id=0\) by bilinearity.
\end{proof}

The ideal \(I\) describes the support of the off-diagonal
Galois module. To propagate Hecke support across these
reducible curves one must still place the corresponding
primes in the support of the patched module \(\mathcal P\).
The preceding calculation identifies the primes and the
extension classes that enter that comparison.


\subsection{Relations among the ordinary extension classes}
\label{sec:oddrelations411}

Keep the curve and the rings \(H\) and \(H^\sharp=H\oplus N\)
of Section~\ref{sec:oddblock410}. Write \(\mathfrak m\) for
the maximal ideal of \(H\), \(J=(p(N,N))\) for the ideal of
products and \(D=3([F:\QQ]+q)\) for the bound of
Theorem~\ref{thm:limitselmer407}. The coefficient prime of the
curve defines the height-one prime \((\varpi)\) in \(O'[[s]]\).
Indeed the evaluation modulo \(\varpi\) is the isomorphism
\(k'[[s]]\to k'[[z]]\). Let \(\mathcal O\) be the completion
of \(O'[[s]]_{(\varpi)}\). This is a complete discrete valuation
ring with uniformizer \(\varpi\) and residue field \(K\).
Its map to \(H\) is induced by localization and completion.
It is injective because \(H\) has a characteristic-zero
prime over the stable prime in Corollary~\ref{cor:selmergenerators410}.
In particular \(H/\varpi H\) contains \(K\) as a coefficient
field. Put
\[
 J_2=(\varpi)+\mathfrak m^2+J,\qquad
 V=(\mathfrak m/J_2)^\vee ,
\]
where the dual is over \(K\). Thus \(H/J_2\) is the
split square-zero algebra \(K\oplus V^\vee\). In a module
over this quotient expressions such as \(h-\bar h\) use
this splitting, with \(\bar h\) the residue of \(h\).

Let \(E\) be the diagonal-block matrix algebra and \(M\) the
off-diagonal matrices for \(r_1\oplus r_2\). Write \(H^1_o(X)\)
for polarized ordinary cohomology of the limiting group with
coefficients in \(X\) over \(K\), including the prescribed
old inertia conditions. Cocycles and local witnesses have
bounded denominators. This is the fraction-field version of
Theorem~\ref{thm:limitselmer407}. Choose a basis
\(\lambda_1,\dots,\lambda_b\) of \((N/\mathfrak m N)^\vee\)
which maps to the chosen basis of \(H^1_o(M)\) under the
isomorphism in Corollary~\ref{cor:selmergenerators410}.
Choose \(n_j\in N\) lifting the dual basis. Reduction of the
matrix representation in Proposition~\ref{prop:orthogonalblocks410}
in those odd coordinates gives the actual cocycles
\(\nu_j\) in their chosen classes. We may use these representatives
in the calculations of Sections~\ref{sec:selmerlength368}
and~\ref{sec:selmerlift369} over \(K\). To verify this, write
the difference of representatives as a coboundary with
skew-adjoint off-diagonal witness \(Z_j\). Replacing the
universal first-order matrix basis by
\(1+\sum_j e_j Z_j\), or its inverse according to the
coboundary convention, changes the representatives as required.
It preserves the form because \(e_i e_j=0\). Conjugation
identifies the two coefficient modules \(W\) and also the
two mixed modules \(U\); it is the identity on their
diagonal quotients and transports their local flags. The
latter flags are intrinsic because their simple inertia
eigenvalues distinguish the flagged lines. Clearing a
finite set of denominators gives the same fraction-field
cohomology and lifting bounds.

\begin{lemma}[Tangents preserving the two blocks]\leavevmode
\label{lem:eventangent411}
There is a linear injection \(a:V\to H^1_o(E)\).
\end{lemma}
\begin{proof}
For \(v\in V\) evaluate \(h\in H\) in the dual numbers as
\(\bar h+\epsilon v(h-\bar h)\), and send \(N\) to zero.
This is a ring map from \(H^\sharp\) since it kills \(J\).
The group matrices now give a polarized block diagonal lift
of \(r\). Its first-order cocycle represents \(a(v)\).
The ordinary word and simple-root projectors give the
ordinary flag, as in Corollary~\ref{cor:selmergenerators410}.
Projecting a local witness onto the diagonal blocks preserves
upper triangularity in the original block-adapted flags.
The old inertia words give the old condition.

Here is the continuity argument, also used for mixed
evaluations below. On the complete presentation ring before
localization the coefficient maps to a square-zero
\(K\)-algebra fixing \(O'[[s]]\) have bounded denominators
and are continuous into its finite-dimensional \(z\)-adic
topology by Lemma~\ref{lem:parameterdepth409}. That
evaluation argument only uses the finite power series
presentation and applies to either kind of curve. Choose
a finite set of generators of the initial Cayley--Hamilton
algebra as a module. The group map into that module is
continuous for the original adic topology. After any fixed
matrix evaluation at those generators, continuity follows
from continuity of the scalar maps and the finite number
of their coordinates. The same argument applies when the
coefficients of group matrices are written in a fixed
basis over a finite-dimensional test algebra.

If \(a(v)=0\), strict conjugation gives the constant lift
of \(r\). The determinant values are constant. Simple-root
projectors at an ordinary inertia element recover each
local character as the trace of the projector times the
local matrix. The first auxiliary matrix word recovers
each ordered Frobenius root from the corresponding distinct
curve eigenvalue, and the second forces its diamond
coordinate to be zero since inertia on the curve is trivial.
Thus the evaluation is constant on every generating datum
of the original presentation. Its continuity there gives
equality on the full ring. Localization introduces no new
choice. Diamond completion also introduces none because
evaluation kills the square of the diamond ideal; the same
holds at the final completion for the square of the curve
maximal ideal. We obtain \(v=0\). Linearity follows by
taking coefficients in a square-zero vector space and
applying linear functionals to these coefficients.
\end{proof}

Form \(U=M\theta\oplus\bigoplus_j E e'_j\) and
\(W=E\oplus\bigoplus_j M e_j\) with the actions in
Section~\ref{sec:selmerlift369}. In particular their action
on the CM group is conjugation by
\(R(g)=(1+\sum_j e_j\nu_j(g))r(g)\); the extra products
are \(e_i e'_j=\delta_{ij}\theta\), all other products
between added basis elements being zero. Define
\[
 \begin{split}
 \mathcal U &=\{(v_j)\in V^b:(a(v_j))_j\text{ is in the
  image of }H^1_o(U)\longrightarrow H^1_o(E^b)\},\\
 V_0&=\{v\in V:(0,\ldots,v,\ldots,0)\in\mathcal U
                        \text{ in every position}\}.
 \end{split}
\]
The image condition depends only on classes. If a projected
cocycle differs from the desired one by a coboundary,
subtract a lifted coboundary. Such witnesses lift in the
coefficient vector spaces, and averaging against the
involution makes the lift polarized since two is a unit.

\begin{proposition}[Two dimension inequalities]\leavevmode
\label{prop:pairedrank411}
The spaces just defined satisfy
\[
             \operatorname{codim}_{V^b}\mathcal U\le b,
                    \qquad \dim_K V_0+b^2\le D .
\]
\end{proposition}
\begin{proof}
Apply part 2 of Theorem~\ref{thm:limitselmer407} to any
\(b+1\) tuples. There are finitely many cocycles and local
witnesses here, so clearing one denominator puts them in
the integral lattices. Over \(K\) both this denominator and
the factor \(z^C\) are invertible. Primitivity gives a
nonzero linear dependence modulo \(\mathcal U\). This is
the first inequality.

For the second, quotient \(W\) by all its off-diagonal
summands except the \(j\)-th, and call the result \(W_j\).
It is the preimage of the \(j\)-th copy of \(E\) in \(U\),
with \(M e_j\) identified with \(M\theta\).
Indeed the action on its off-diagonal coefficient acquires
exactly the commutator of \(\nu_j(g)\) with the conjugated
diagonal coefficient. This also identifies the upper
triangular conditions there, by using the local change of
basis for \(\nu_j\) in Section~\ref{sec:selmerlift369}.
More explicitly, after a lift has been chosen with zero
other diagonal cocycles, the projected local witnesses for
those other coordinates are invariant in \(E/B_u(E)\).
This quotient has zero local invariants since an ordinary
inertia element has distinct eigenvalues on the flagged
lines. Those projected witnesses vanish. Representatives
in the upper triangular subspace can be subtracted after
lifting to the upper part of \(U\), whose projection onto
the upper part of \(E^b\) is surjective. The resulting
witnesses lie in the stated preimage. The same argument
at old inertia kernels is immediate because those kernels
act trivially.

Consequently, for \(v\in V_0\) a representative of \(a(v)\)
lifts to an ordinary cocycle in each \(W_j\). Arrange that
the projections of these lifts coincide as cocycles, by
the preceding coboundary adjustment. The modules \(W_j\)
have fiber product \(W\) over their common diagonal
quotient \(E\). The local upper triangular conditions
have the same fiber product description. In block-adapted
coordinates the upper parts consist of the common upper
diagonal coefficient and the separate upper off-diagonal
coefficients, transported by the respective local changes
of basis. At an ordinary place the projected witnesses
for the common diagonal cocycle agree in \(E/B_u(E)\)
by the same vanishing of invariants. They can therefore
be represented by witnesses with a common diagonal part,
and assembled in \(W\). This gives a lift in \(H^1_o(W)\)
of every class in \(a(V_0)\).

The inclusion \(M^{\oplus b}\to W\) injects its ordinary
cohomology into \(H^1_o(W)\). Indeed if an off-diagonal
cocycle is a polarized coboundary in \(W\), the diagonal
part of its witness commutes with \(r\) and is skew-adjoint.
A commuting element of \(E\) is scalar on each absolutely
irreducible block, and such a scalar is skew-adjoint only
when zero. Thus the witness was already off-diagonal.
Upper triangular witnesses for a class in the submodule
are witnesses in \(W\), using its transported flags.
This supplies a subspace of dimension \(b^2\) killed in
\(H^1_o(E)\), by part 3 of Theorem~\ref{thm:limitselmer407}.
The lifted subspace maps onto the \(\dim V_0\) independent
classes of the lemma. Part 1 of the same theorem bounds
\(\dim H^1_o(W)\) by \(D\), proving the second inequality.
\end{proof}

\begin{lemma}[Evaluating products with an extension class]\leavevmode
\label{lem:pairedproduct411}
For \(\varphi\in (N/J_2 N)^\vee\) define
\[
             v_j(h)=\varphi\big((h-\bar h)n_j\big)
                  \quad(h\in H/J_2).
\]
On \(\mathfrak m/J_2\) these are elements of \(V\), and the
tuple \((v_j)\) belongs to \(\mathcal U\).
\end{lemma}
\begin{proof}
Use the commutative test algebra just defined with basis
\(1,e_j,e'_j,\theta\). The following formulas determine a
homomorphism into it from \(H^\sharp\):
\[
 h\longmapsto\bar h+\sum_j v_j(h)e'_j,\qquad
 n\longmapsto\sum_j\lambda_j(\bar n)e_j+\varphi(n)\theta.
\]
Here \(\bar n\) denotes reduction modulo \(\mathfrak m N\)
and \(\varphi\) is read modulo \(J_2 N\). Products of two
scalars give the first formula because \((\mathfrak m/J_2)^2=0\).
Products of two elements of \(N\) have image zero on
both sides. For a mixed product, multiplication by
\(h-\bar h\) on \(N/J_2 N\) factors through
\(N/\mathfrak m N\). Hence
\[
 \varphi(hn)=\bar h\varphi(n)
                  +\sum_j v_j(h)\lambda_j(\bar n),
\]
as required. The formulas respect addition and the unit.

Apply this map entrywise to the group matrices. Modulo the
square-zero ideal generated by the \(e'_j\) and \(\theta\)
we obtain exactly \(R(g)\). The difference from \(R(g)\)
multiplied on the right by \(R(g)^{-1}\) is an
\(U\)-valued cocycle: its \(e'_j\) entries are block diagonal
and its \(\theta\) entries off-diagonal. Its projection to
\(E^b\) consists of the actual cocycles of the \(v_j\) in
Lemma~\ref{lem:eventangent411}. Its polarization follows
from the matrix form, and it kills the old inertia kernels.
It is bounded and continuous by the finite module argument
of the lemma, since the scalar evaluation on \(H\) has only
first-order coefficients.

For the local condition, use the ordinary words in \(C\).
The eigenprojectors for a local element with simple roots
lift the ordered basis of the ordinary flag for \(R\)
across the square-zero ideal; projecting the basis vectors
and applying Nakayama gives a basis. The words with one
arbitrary letter then triangularize every local group matrix
with its specified diagonal. Thus the cocycle has a local
witness in the matrices with entries in this ideal, in
coordinates for the flag of \(R\). A witness can be taken
in \(U\). To see this, apply the involution which conjugates
by \(\mathrm{diag}(1,-1,-1)\) in the block basis and negates
the variables \(e_j,\theta\) while fixing \(e'_j\).
Its invariant part on the square-zero ideal of matrices
is exactly \(U\). It fixes \(R\) and preserves its ordinary
flag, since the original flag respects the two blocks and
the local change of basis for \(R\) has off-diagonal
\(e_j\)-terms. Averaging the witness gives the required
invariant witness. This proves membership in \(\mathcal U\).
\end{proof}

\begin{theorem}[A square matrix of relations from ordinary cohomology]\leavevmode
\label{thm:oddrelations411}
One has \(b^2\le D\). There are \(D\) elements \(x_k\in\mathfrak m\)
spanning \(\mathfrak m/J_2\) modulo \(J_2\); designate the
first \(b^2\) of them as \(x_{ij}\). For some \(w_{ij}\in J_2\)
the following equalities hold in \(N\):
\[
                 \sum_{j=1}^b(x_{ij}-w_{ij}) n_j=0
                         \qquad(1\le i\le b).
\]
The \(x_k\) are allowed to be dependent in \(H\).
\end{theorem}
\begin{proof}
By Proposition~\ref{prop:pairedrank411}, the annihilator of
\(\mathcal U\) in \((V^b)^\vee\) is spanned by \(b\) forms,
padding with zeros if necessary. Write their entries as
\(\ell_{ij}\in V^\vee\), so that
\(\sum_j\ell_{ij}(v_j)=0\) for tuples in \(\mathcal U\).
The common kernel of the \(b^2\) entries is precisely
\(V_0\): a vector belongs to it if and only if each single
insertion is killed by every row. Their span in \(V^\vee\)
therefore has a complement of dimension \(\dim V_0\le D-b^2\).
Use the entries followed by a spanning list for this
complement, padded as needed, to obtain \(D\) linear forms
spanning \(V^\vee=\mathfrak m/J_2\). Lift them to the
claimed \(x_k\).

For every functional \(\varphi\) of the last lemma we have
\[
  \varphi\!\left(\sum_j x_{ij}n_j\right)
          =\sum_j\ell_{ij}(v_j)=0 .
\]
Functionals separate vectors in the finite-dimensional
space \(N/J_2N\). The sum thus lies in \(J_2 N\), and the
\(n_j\) generate \(N\) by Nakayama. Writing the sum in
these generators with coefficients \(w_{ij}\in J_2\)
proves the asserted equalities. This is the presentation
step in \srcloc{Development/AutomorphicArithmetic/AutomorphicVarieties/PatchingModules/SelmerScaling/AlCSmul}{12}.
\end{proof}


\subsection{Hecke support through the curves with two constituents}
\label{sec:reduciblepropagation412}

The square matrix in Theorem~\ref{thm:oddrelations411}
annihilates the vector of generators of \(N\). We first
determine the component on which that vector can be nonzero.
The argument takes place over the same coefficient DVR
\(\mathcal O\) and uses a projective space of kernel lines.

\begin{lemma}[The incidence locus after completion]\leavevmode
\label{lem:incidence412}
Let \(b\ge1\), \(u\ge0\) and let \(S\) be the formal
power series ring over a complete DVR in the entries of a
\(b\)-by-\(b\) matrix \(A\), a column \(Y\) of length \(b\),
and \(u\) further variables. The radical of the ideal
\((AY,\det A)\) is prime. Every prime of \(S\) containing
the entries of \(AY\) has height at least \(b\).
\end{lemma}
\begin{proof}
Begin with polynomial coordinates over the DVR. Over the
affine space of these coordinates consider the projective
incidence scheme obtained by adjoining a line \(L'\) in a
free module of rank \(b\) and requiring
\[
                       A L'=0,\qquad Y\in L' .
\]
It is closed in the product with projective space. Over
projective space the column lies in the tautological
line and the matrix is an arbitrary map out of the
quotient by that line. Thus this scheme is the total
space of a vector bundle, with the \(u\) additional free
coordinates, and is regular.

Localize the polynomial base at the origin of its closed
fiber. The incidence scheme remains regular and the
projection is proper. It stays regular after completing
this local base. Indeed at each point of the closed fiber
the local rings before and after this base change have
isomorphic completions. Both schemes are of finite type
over their Noetherian local bases; every closed point
of the scheme after base change lies in the closed fiber
by properness. Regularity there implies regularity at
every generization. The base change at issue is precisely
the one to \(S\).
\href{https://stacks.math.columbia.edu/tag/0BG4}
{The Stacks Project, Lemmas 54.11.1 and 54.11.2.}

The completed incidence scheme is integral. It is regular
and Noetherian, so its irreducible components are disjoint
open and closed subsets. Each nonempty component meets
the closed fiber, since its image under the proper map is
nonempty and closed in the spectrum of a local ring. The
closed fiber is a single projective space over the residue
field and is connected, so there is just one component.
At any prime of \(S\) the fiber has a point precisely
when \(AY=0\) and \(\det A=0\) in its residue field.
If \(Y\ne0\) use its span, and if \(Y=0\) use any line
in the kernel of the singular matrix. Conversely such a
line gives both equations. The proper image of the integral
incidence scheme is therefore exactly \(V(AY,\det A)\),
which is irreducible.

For the height bound first reduce the polynomial coordinates
modulo the uniformizer. Away from \(Y=0\), invert a nonzero
coordinate of \(Y\) and eliminate one entry of each row of
\(A\). This locus has dimension \(b^2+u\). The locus
\(Y=0\) also has dimension \(b^2+u\). Thus the \(b\)
entries of \(AY\) cut out codimension \(b\) in affine
space over the residue field. In its regular local ring
at the origin these equations form a regular sequence:
the homogeneous components of the zero set pass through
the origin, and \(b\) equations of height \(b\) in a
Cohen--Macaulay local ring form part of a system of
parameters. Adjoining the uniformizer first gives a regular
sequence upstairs. Permutability in the local ring and
flatness of completion leave the \(b\) entries a regular
sequence in \(S\). Any prime containing them has height
at least \(b\), as asserted.
\end{proof}

\begin{proposition}[A component through every prime in the support of \(N\)]\leavevmode
\label{prop:localcomponent412}
Let \(Z\) be a prime in \(H\) whose geometric determinant is
irreducible and such that \(\dim H/Z\ge D+1\). Then
\(Z\subseteq Q\) for every \(Q\in\operatorname{Supp}_H N\).
\end{proposition}
\begin{proof}
Use the \(D\) scalars \(x_k\), the \(b\) odd generators
\(n_j\), and the corrections \(w_{ij}\) of
Theorem~\ref{thm:oddrelations411}. Since \(N^2=J\subseteq
\mathfrak m\), the finite \(H\)-algebra \(H^\sharp=H\oplus N\)
is local with maximal ideal \(\mathfrak m\oplus N\). Its
square lies in \(\mathfrak m H^\sharp\), and it contains
\(\mathfrak m H^\sharp\) before squaring. The two adic
topologies coincide and it is complete. Map
\[
       S=\mathcal O[[X_1,\ldots,X_D,Y_1,\ldots,Y_b]]
                 \longrightarrow H^\sharp
\]
by \(X_k\mapsto x_k\), \(Y_j\mapsto n_j\). The map is
surjective. Indeed modulo the square of its maximal ideal
the target is generated, in positive degree, by the images
of \(\varpi,X_k,Y_j\): the scalars \(x_k\) span modulo
\(J_2\), \(J\) consists of sums of products of odd elements,
and the \(n_j\) generate \(N\). Both residue fields are
\(K\). Nakayama gives generation of the maximal ideal.
Lifting the residue and then correcting at each successive
degree by homogeneous polynomials proves surjectivity by
completeness.

Lift each \(w_{ij}\in J_2\) to
\(W_{ij}\in(\varpi)+\mathfrak m_S^2\), since elements of
\(J\) map into the square upstairs in \(H^\sharp\)
and the surjection maps maximal ideals onto maximal ideals.
For the first \(b^2\) coordinates set
\(A_{ij}=X_{ij}-W_{ij}\). With the remaining \(X_k\)
and the \(Y_j\) these are again formal coordinates over
\(\mathcal O\). In fact the corresponding substitution
fixes \(\varpi\), and on the cotangent space it changes
each \(X_{ij}\) only by a multiple of \(\varpi\).
Its linear part is invertible. It induces an isomorphism
on the associated graded polynomial ring of the regular
local ring \(S\), hence on \(S\) by successive approximation
and separatedness. In the new coordinates the entries of
\(AY\) vanish in \(H^\sharp\). The matrix \(A\) has an
image \(A_H=(x_{ij}-w_{ij})\) in \(M_b(H)\). Multiplying
\(A_H(n_j)_j=0\) by its adjugate shows that
\(\det A_H\in\operatorname{Ann}_H N\).

The extension \(H\subset H^\sharp\) is finite integral
and injective. Choose a prime over \(Z\) in \(H^\sharp\)
and let \(\widetilde Z\) be its inverse image in \(S\).
The quotient has dimension \(\dim H/Z\) by integral
extension. Since \(\dim S=D+b+1\), the height of
\(\widetilde Z\) is at most \(b\). Some \(Y_j\) has nonzero
image at this prime. Otherwise it would contain \(N\)
and all its products, contrary to \(J\not\subset Z\)
by Theorem~\ref{thm:oddsupport410}. Thus \(AY=0\) gives
\(\det A\in\widetilde Z\). By Lemma~\ref{lem:incidence412}
the prime radical of \((AY,\det A)\) has height at least
\(b\), and is contained in \(\widetilde Z\). The two
primes must coincide.

For \(Q\in\operatorname{Supp}_H N\) take a prime of
\(H^\sharp\) over \(Q\) and its inverse image
\(\widetilde Q\subset S\). It contains \(AY\) by the
relations and \(\det A\) by the annihilator calculation.
It consequently contains \(\widetilde Z\). Descending
through the surjection and contracting to \(H\) gives
\(Z\subseteq Q\).
\end{proof}

For a prime \(\mathfrak q\) in the original ordinary
determinant ring \(R_0\), take the finite continuous
Cayley--Hamilton algebra for the base group \(\Gamma_0^0\)
and impose the ordinary, old inertia, and adjoint matrix
words, as in Lemma~\ref{lem:finitech410} without the new
places. Denote this algebra by \(C_0\). Its geometric
fiber means extension to an algebraic closure of
\(\kappa(\mathfrak q)\). An irreducible prime is one whose
geometric determinant is irreducible. We also consider
reducible primes at which this algebra has a nonzero
Jacobson radical in the geometric fiber. These are the
two cases tested by component propagation.

\begin{theorem}[Ordinary Hecke support through a two-constituent curve]\leavevmode
\label{thm:reduciblepropagation412}
Let \(\mathfrak p\subset R_0\) be any of the selected
characteristic-\(\ell\) curves having a character and a
plane. If \(\mathfrak q\subseteq\mathfrak p\) is irreducible,
or is reducible with nonsemisimple geometric fiber of
\(C_0\), then \((M_K)_{\mathfrak q}\ne0\). Equivalently,
the kernel of the natural good Hecke map \(R_0\to B_K\)
is contained in \(\mathfrak q\).
\end{theorem}
\begin{proof}
Let \(\mathcal P_H\) be the patched module localized at
\(\mathfrak P\) and completed over \(H\). Its depth is
at least \(D+1\). The proof of
Lemma~\ref{lem:parameterdepth409} uses the \(3q\) regular
diamond variables, the coefficient uniformizer, the \(3[F:\QQ]\)
translated weight variables and the nonzero ordered fiber.
These facts hold for both kinds of selected curves by
Theorem~\ref{thm:modulepatch408}, and the regular sequence
persists under faithfully flat completion.

We first produce a sufficiently large irreducible prime in
its support. By Theorem~\ref{thm:stableprime371} and the
classical systems in the proof of
Theorem~\ref{thm:modulepatch408}, there is an irreducible
prime \(\mathfrak q_1\subset\mathfrak p\) in the support
of \(M_K\). Indeed the same good Hecke annihilator
vanishes at every point of the dense family defining
\(\mathfrak q_1\). Choose a prime over \(\mathfrak q_1\)
in \(P_{\mathfrak x}\) by flat local descent as in
Corollary~\ref{cor:selmergenerators410}. Its pullback to
\(T_{\mathfrak P}\) is in the support of the patched
module. To see this, first reduce modulo the diamonds;
the ordered fiber of that module there is precisely
\(P_{\mathfrak x}\otimes_{R_0}M_K\), whose residue fiber
is a nonzero field extension of the fiber at
\(\mathfrak q_1\). Lift this prime to \(H\) by faithful
flatness. Its residue field extends the preceding one, so
it is still in support and still has irreducible determinant.
Choose a minimal prime \(Z\) of the support below this
lifted prime. Irreducibility holds at \(Z\) since its
locus is the complement of \(V(J)\). A minimal prime of
the support of a finite module is associated to the module,
and at any associated prime its quotient ring has dimension
at least the depth of the module. Hence \(\dim H/Z\ge D+1\).
Proposition~\ref{prop:localcomponent412} now gives
\(\operatorname{Supp}_H N\subseteq\operatorname{Supp}_H\mathcal P_H\),
since support is closed under specialization.
\href{https://stacks.math.columbia.edu/tag/00LE}
{The Stacks Project, Section 10.72, \emph{Depth}.}

Given \(\mathfrak q\) of the theorem, choose a prime over it
in \(P_{\mathfrak x}\), take its preimage in \(T_{\mathfrak P}\)
and lift to a prime \(Q\) of \(H\). These choices use the
same flat local maps as above. At \(Q\) all diamond
variables vanish. If \(\mathfrak q\) is irreducible then
so is \(Q\), and it belongs to \(\operatorname{Supp}_H N\)
by Theorem~\ref{thm:oddsupport410}. In the other case the
geometric fiber of \(C\) there surjects onto the geometric
fiber of \(C_0\) after extending the residue fields to a
common algebraically closed field. In fact pull the group
matrices of \(C_0\) back through \(h:\Gamma_\infty\to
\Gamma_0\). They satisfy the original determinant and its
matrix words. The new inertia is one, and the ordered
Frobenius polynomial on \(P\) is the characteristic
polynomial of the base determinant. Thus the first auxiliary
word is zero by Cayley--Hamilton, and the second is zero
since all diamonds vanish. The same holds after adjunction.
These facts give the asserted map by the base change property
of the initial continuous Cayley--Hamilton algebra in
Lemma~\ref{lem:finitech410}. Equivalently one extends the
finite integral base group algebra quotient along the
ordered base map and then specializes; its generating group
elements lift through \(h\), proving surjectivity.

A nonzero radical over an algebraically closed field remains
a nonzero nilpotent ideal after further field extension.
A quotient of a semisimple finite-dimensional algebra is
semisimple, so the source geometric fiber is nonsemisimple
as well. The determinant stays reducible under this
extension. Theorem~\ref{thm:oddsupport410} again gives
\(Q\in\operatorname{Supp}_H N\). In both cases
\(\mathcal P_H\) has nonzero residue fiber at \(Q\).
Since the diamonds vanish there, that fiber is also the
residue fiber of the ordered base module, extended to
\(\kappa(Q)\). Hence the original fiber of \(M_K\) at
\(\mathfrak q\) is nonzero. The equivalent ideal
containment follows because the good Hecke algebra is
the image of \(R_0\) on the finite module \(M_K\).
This proves the single-curve support implication for the
selected fields; the corresponding inclusion is formulated
in \srcloc{Development/AutomorphicArithmetic/AutomorphicVarieties/WeilGroupCharacters/UniversalOrdinaryDeterminantRing}{288}.
\end{proof}

Together with Theorem~\ref{thm:irrcurvesupport409} this
settles local propagation at the curves selected from the
ordinary family. The conclusions concern prime support.
The bound on lengths in Proposition~\ref{prop:res23}
requires control of integral multiplicities at the original
weights. The signed projected periods must also be compared
on those same fields and intervals.


\subsection{The exact annihilator on the irreducible component}
\label{sec:incidenceannihilator413}

In the presentation of Proposition~\ref{prop:localcomponent412}
write \(Y_j\) for the coordinate mapping to the off-diagonal
generator \(n_j\), and \(\mathfrak a\) for the ideal of the
auxiliary diamond variables in the patched ring. The first
coordinates describe a projective incidence calculation
before the specialization by \(\mathfrak a\). We retain
the equations of that calculation to determine which
elements act as zero on \(N\).

\begin{lemma}[The two components of the matrix equation]\leavevmode
\label{lem:reducedincidence413}
In the formal power series ring \(S\) of
Lemma~\ref{lem:incidence412}, let
\[
         L=((AY)_1,\ldots,(AY)_b),\qquad
         P_0=(Y_1,\ldots,Y_b),\qquad
         P_1=\sqrt{L+(\det A)} .
\]
Then \(S/L\) is reduced and \(L=P_0\cap P_1\). Both primes
have height \(b\).
\end{lemma}
\begin{proof}
The ideal \(P_0\) is prime of height \(b\).
The prime \(P_1\) was constructed in
Lemma~\ref{lem:incidence412}. It contains \(\det A\), which
is nonzero modulo \(P_0\). Also \(Y_1\notin P_1\). Indeed
one may reduce coefficients to the residue field, send
every matrix entry and all free extra variables to zero,
and send \(Y_1\) to the parameter of a formal power series
ring over that field, with the other \(Y_j\) zero. The
kernel of this map is prime and contains \(L+(\det A)\).
For any prime containing \(L\), either all \(Y_j\) vanish
there or \(\det A\) does, by the adjugate identity.
The two incomparable primes displayed in the statement are
therefore exactly the minimal primes. The height theorem
for \(b\) generators together with Lemma~\ref{lem:incidence412}
gives height \(b\) for both.

At \(P_0\), the matrix is invertible and \(L S_{P_0}=P_0 S_{P_0}\).
At \(P_1\), consider the \(b\) derivations
\(\partial/\partial A_{i1}\), defined term by term on
formal series and extended by the quotient rule to
\(S_{P_1}\). They send a square of the local maximal ideal
into that ideal by the product rule. On the generators
\((AY)_j\) their matrix is \(Y_1\) times the identity.
Thus these generators are independent in the cotangent
space at \(P_1\). To check this, lift coefficients of a
linear relation there to the local ring and apply the
derivations; the terms differentiating the coefficients
vanish in the residue field because \(AY=0\) there.
Since \(S_{P_1}\) is regular local of dimension \(b\), its
cotangent space has dimension \(b\). By Nakayama the entries
generate its maximal ideal. Hence \(S/L\) is reduced at
both minimal primes.

The entries of \(AY\) form a regular sequence by
Lemma~\ref{lem:incidence412}. The quotient is a complete
intersection and remains Cohen--Macaulay after localization.
Its associated primes are therefore minimal, since at
an associated prime depth is zero. If its nilradical
were nonzero this finite submodule would have an associated
prime, also associated to \(S/L\), where its localization
would be nonzero. That contradicts the generic reducedness
just proved. The intersection formula follows.
\end{proof}

Choose the presentation \(S\to H^\sharp\) of
Proposition~\ref{prop:localcomponent412} equivariantly
for the involution fixing \(H\) and sending \(n\) to
\(-n\) on \(N\). In the original coordinates the map is
equivariant when the \(X_k\) are fixed and the \(Y_j\)
are negated. Each correction \(w_{ij}\) is even, so
average any of its lifts in \((\varpi)+\mathfrak m_S^2\)
with its transform. The corrected matrix entries
\(A_{ij}=X_{ij}-W_{ij}\) are again fixed by the involution.
Thus \(P_0,P_1,L\) in the new coordinates are invariant.
Let \(D_1=S/P_1\), and denote its even and odd submodules
under this involution by \(D_1^+\) and \(D_1^-\).
These eigenspaces split because two is invertible.

\begin{theorem}[The off-diagonal module on the incidence component]\leavevmode
\label{thm:incidenceannihilator413}
At every selected two-constituent curve there is a unique
prime \(Z\) in the completed patched ring \(H\) with
irreducible geometric determinant and \(\dim H/Z\ge D+1\).
For the choices of coordinates above, the natural maps give
\[
     H/Z\simeq D_1^+,\qquad N\simeq D_1^-,\qquad
     \operatorname{Ann}_H N=Z .
\]
In particular \(N\) is a torsion-free module of rank one
over \(H/Z\). Its product ideal \(J\) maps injectively into
this domain, with image generated by the products of the
images of the \(Y_iY_j\).
\end{theorem}
\begin{proof}
The proof of Theorem~\ref{thm:reduciblepropagation412}
constructs such a prime \(Z\) in the support of the patched
module. Write \(K_0=\ker(S\to H^\sharp)\). For any prime
with the stated conditions, the proof of
Proposition~\ref{prop:localcomponent412} shows that the
inverse image in \(S\) of a prime over it in \(H^\sharp\)
is exactly \(P_1\). Its contraction to \(H\) is therefore
determined uniquely. These conclusions also give
\[
                       L\subseteq K_0\subseteq P_1 .
\]
Consequently \(H^\sharp\) has an equivariant surjection to
the domain \(D_1\), and its kernel on \(H\) is \(Z\).
Taking even and odd eigenspaces is exact by the projectors
\((1\pm\sigma)/2\), with \(\sigma\) the involution. It
follows that the induced maps \(H/Z\to D_1^+\) and
\(N\to D_1^-\) are surjective and the first is injective.

For injectivity on \(N\), every odd element of \(S\)
lies in \(P_0\), since \(S/P_0\) is entirely even.
An odd element mapping to zero in \(D_1\) therefore lies
in \(P_0\cap P_1=L\), hence is already zero in \(H^\sharp\).
Every odd element of \(H^\sharp\) has an odd lift by
averaging. This proves the second isomorphism. For the
product ideal, an element of \(J\) lifts to the ideal
\((Y_iY_j:1\le i,j\le b)\) in \(S\): write it as a sum of
products of the \(H\)-generators \(n_i,n_j\) with scalars
in \(H\), and lift those scalars. The lift belongs to
\(P_0\), so the same intersection argument gives injectivity
on \(J\). Its image is the indicated ideal of the invariant
ring since all even coefficients lift to \(H\).

The nonzero module \(D_1^-\) is torsion-free over
\(D_1^+\) because it lies in the domain \(D_1\), and
\(Y_1\) has nonzero image in it. Its annihilator over \(H\)
is precisely \(Z\). It has rank one. Indeed division by a
fixed nonzero odd element takes every odd element to an
invariant fraction. Such fractions are in
\(\operatorname{Frac}(D_1^+)\): for an invariant fraction
\(a/c\) multiply numerator and denominator by \(\sigma(c)\).
Both become invariant and the denominator is nonzero. This
bounds the rank above by one and the nonzero element gives
equality.
\end{proof}

\begin{corollary}[An annihilator acting integrally on both off-diagonal blocks]\leavevmode
\label{cor:unpowered413}
Every element of the annihilator in \(H\) of the completed
patched module \(\mathcal P_H\) kills \(N\) and its product
ideal \(J\). More precisely \(ZJ=0\) and \(Z\cap J=0\).
\end{corollary}
\begin{proof}
The prime \(Z\) was obtained in the support of
\(\mathcal P_H\), so the annihilator of this finite module
is contained in \(Z\). The theorem says \(ZN=0\). For an
element \(h\in Z\) and any \(n,n'\in N\), associativity in
\(H^\sharp\) gives \(h(nn')=(hn)n'=0\). This kills all of
\(J\). Injectivity of \(J\to H/Z\) proves the intersection
assertion.
\end{proof}

The corollary gives an equality of annihilators for the
off-diagonal module over the completed patched ring and
an inclusion for the patched Hecke module. Passage to
the original fixed weights and evaluation of the signed
integral periods remain necessary to obtain the two
inequalities in Section~\ref{sec:missing}.


\subsection{The unramified fiber and the integral extension module}
\label{sec:unramifiedcorner414}

Continue at a selected curve with a character and a plane.
Let \(\mathfrak a=(Y_{v,i})\) in the patched ring \(H\);
the new inertia operators are \(1+Y_{v,i}\). The symbols
\(Y_1,\ldots,Y_b\) in the incidence presentation still
denote the generators mapping to \(N\). We now set all
the new inertia operators to one. Our first task is to
identify the resulting ring, including its nilpotents.

\begin{lemma}[The ordinary ring after removing the auxiliary inertia]\leavevmode
\label{lem:diamondfiber414}
The map to the ordered base induces an isomorphism of
complete local rings
\[
                H/\mathfrak a H\ \simeq\ H_0,
                     \qquad H_0=\widehat{P_{\mathfrak x}} .
\]
Here \(P\) includes the coefficient extension and parameter
of Section~\ref{sec:irrcurvesupport409}.
\end{lemma}
\begin{proof}
Use the original complete presentation ring over \(O'[[s]]\)
before discriminant localization and diamond completion.
Impose the scalar relations defining \(T\) and all
\(Y_{v,i}=0\) there, and denote the quotient by \(B\).
The scalar relations are the trace and character expressions
in group elements on the fixed presentation, so have lifts
in this ring. Let \(\mathfrak t\) be the curve prime of
\(B\) and \(S_0=B_{\mathfrak t}\). The discriminant \(f\)
is a unit in \(S_0\). Adic completion modulo its defining
ideal agrees with the quotient before completion in a
Noetherian ring. Hence localizing \(T/\mathfrak a\) at the
curve gives exactly \(S_0\). The original ordered base data
give a surjection \(B\to P\), taking determinants from
\(\Gamma_0^0\) and choosing ordered roots. Surjectivity
follows as in Section~\ref{sec:irrcurvesupport409} from
generation by traces, characters, \(s\), and the roots.

Over \(B\) form the finite Cayley--Hamilton algebra with the
matrix words and their adjoints as in Lemma~\ref{lem:finitech410}.
Call it \(C_B\). Extend scalars first to \(S_0\). At a new
place the differences of the \(U_i=U_{v,i}\) are units
and \(\prod_i(\phi-U_i)=0\) in this algebra. Polynomial
interpolation therefore gives orthogonal idempotents
\[
           e_i=\prod_{j\ne i}(\phi-U_j)/(U_i-U_j),
                \qquad \sum_i e_i=1,\qquad \phi e_i=U_i e_i .
\]
The second matrix word gives \(e_i(\tau-1)=0\) for each
\(i\). Adding proves \(\tau=1\). Also \(e_i\) has determinant
degree one on its corner, since its specialization in the
semisimple curve representation is the projector onto a
simple root. The spectrum of the local ring \(S_0\) is
connected, so the corner degree equals this residual degree
by the corner lemma for determinants. Multiplicativity on
orthogonal corners yields
\[
                           D(X-\phi)=\prod_i(X-U_i)
                                  \quad\text{in }S_0[X].
\]
Adjunction also kills the conjugate inertia difference.
The multiplier at the auxiliary inertia is one, as it
comes from the unramified base. \cite{VIII-VI-chenevier364}

For descent through a closed normal subgroup it is important
to retain the original topology. Take
\(B'=B/\ker(B\to S_0)\) and let \(C'\) be the image of
\(C_B\) in \(C_B\otimes_B S_0\). These are quotients of
the original complete finite-residue ring and its finite
module by closed submodules. Thus \(C'\) has its compact
quotient topology and the continuous group map still
takes values in its units. It kills the new inertia and
its conjugates by the preceding computation. By continuity
it kills their closed normal closure, which is the full
kernel over \(\Gamma_0^0\) by
Theorem~\ref{thm:limitkernel407}. The determinant descends
through \(C'\) over \(B'\): elements of the removed module
kernel have zero trace against every element after mapping
to \(S_0\), hence already over \(B'\), which embeds there.
The trace argument with Newton's identities proves descent
of the polynomial law. The remaining local characters,
old inertia and ordinary words give the original ordinary
problem over \(B'\), with the prescribed coefficients and
\(s\). The ordered polynomial identity holds over \(B'\)
as well by its injection into \(S_0\). Universality of the
base ring and of its monic splitting algebra gives a
continuous map \(P\to B'\); the reductions of the roots
pick the indicated local factor. Its image contains all
the group traces, local characters and root coordinates,
because the group has descended and the Frobenius
markings at the base are the specified \(s_v\). The image
is compact and closed, hence the map is surjective.

To check the two directions after localization, put
\(P'=P/\ker(P\to P_{\mathfrak x})\). The map \(B\to P\)
induces \(B'\to P'\), since elements outside \(\mathfrak t\)
map outside \(\mathfrak x\). The composite \(P\to B'\to P'\)
is the original quotient: the continuous maps agree on all
the generators just listed. The prime in \(B'\) defined by
the curve contracts to \(\mathfrak x\) in \(P\) by the
same evaluations. Thus the surjection \(P\to B'\) gives
a surjection \(P_{\mathfrak x}\to S_0\) whose composite
with the original \(S_0\to P_{\mathfrak x}\) is the
identity on \(P_{\mathfrak x}\). These are isomorphisms.
Completing at the curve commutes with the quotient by
the finitely generated diamond ideal. This proves the lemma.
\end{proof}

\begin{proposition}[Base change of the ordinary Galois algebra]\leavevmode
\label{prop:ordinarycorner414}
With \(C_0\) the ordinary base group algebra quotient of
Section~\ref{sec:reduciblepropagation412}, one has
\[
              C/\mathfrak a C\ \simeq\
                     C_0\otimes_{R_0}H_0 .
\]
Let \(N_0\) denote the corner between the chosen primitive
idempotents in its two blocks. Then
\[
                  N_0\simeq N/\mathfrak a N
                     \simeq D_1^-/\mathfrak a D_1^- .
\]
The action in the last quotient is through
\(H\to D_1^+\) of Theorem~\ref{thm:incidenceannihilator413}.
\end{proposition}
\begin{proof}
In the notation of the lemma \(C_B\otimes_B H_0=C/\mathfrak a C\),
since both are formed by scalar extension of the same
initial finite Cayley--Hamilton algebra and quotient by
the scalar extensions of the same matrix relations.
There is a map to \(C_0\otimes_{R_0}H_0\) by pullback from
the base group: there the new inertia is one, the first
auxiliary word is Cayley--Hamilton and the second is zero.
For the reverse map use the continuous representation in
\(C'\) constructed over \(B'\). Its determinant is the
base determinant by the map \(R_0\to B'\); its ordinary
matrix words, old inertia and adjoints are zero. Thus the
universal continuous Cayley--Hamilton group algebra on
\(\Gamma_0^0\), with these matrix relations imposed,
maps to \(C'\). One can form this map over the finite
quotients of the compact algebra and take their inverse
limit. After localizing at the curve and extending to
\(H_0\), its target is \(C_B\otimes_B H_0\); the original
module kernel defining \(C'\) vanishes on localization.
The two maps over \(H_0\) are mutually inverse on all scalar
and group generators. These span the finite algebras,
proving the identity. Taking the two primitive corners
commutes with the scalar extension, since multiplication
by the left and right idempotents is a module projector.
The remaining assertion follows from the identification of
\(N\) in Theorem~\ref{thm:incidenceannihilator413}.
\end{proof}

\begin{theorem}[The matrix and the length at a reducible specialization]\leavevmode
\label{thm:cornermatrix414}
Let \(B_1\) be an \(H_0\)-algebra in which the product
ideal \(J\) has zero image. Let \(A_{B_1}\) be the
\(b\)-by-\(b\) matrix with entries the images of
\(x_{ij}-w_{ij}\) from Theorem~\ref{thm:oddrelations411}.
There is a presentation
\[
 B_1^b \xrightarrow{\ A_{B_1}^{\mathsf T}\ } B_1^b
     \longrightarrow N_0\otimes_{H_0} B_1\longrightarrow 0 .
\]
In particular, if \(B_1\) is a DVR with normalized valuation
\(v\) and the matrix has full rank, the length of this
module is \(v(\det A_{B_1})\). If its rank over the
fraction field is \(r<b\), its free rank is \(b-r\) and
its torsion length is the minimum valuation of a nonzero
\(r\)-by-\(r\) minor, with value zero for \(r=0\).
\end{theorem}
\begin{proof}
Use the formal coordinates of Proposition~\ref{prop:localcomponent412}
chosen equivariantly as in Section~\ref{sec:incidenceannihilator413}.
Put \(S'=\mathcal O[[A_{ij},u_1,\ldots,u_{D-b^2}]]\)
for the ring of coordinates excluding the odd variables.
Then \(S=S'[[Y_1,\ldots,Y_b]]\); the involution negates
every \(Y_j\) and fixes \(S'\). Write \(S^-\) for its odd
part. By Lemma~\ref{lem:reducedincidence413}, the kernel
of \(S^-\to N\simeq D_1^-\) is precisely the odd part
of \(L=(AY)\). Indeed an odd element belongs to
\((Y_1,\ldots,Y_b)=P_0\), and \(P_0\cap P_1=L\).
Furthermore, every odd power series of \(Y\)-degree at
least three belongs to \(\sum_{i,j}Y_iY_j S^-\). This
follows by partitioning its monomials among the finitely
many possible first two factors. The images of those
series in \(N\) are exactly \(JN\), since the products
\(n_i n_j\) generate \(J\) and every element of \(N\)
has an odd lift. The ring \(S'\to H/J\) is surjective:
a scalar in \(H\) has an even lift in \(S\), and all positive
\(Y\)-degrees of that lift are divisible by products
\(Y_iY_j\) with even coefficients.

It follows that computing \(N/JN\) over \(S'\) requires
only the terms of \(Y\)-degree one. They form the free
module on \(Y_1,\ldots,Y_b\). The relations in this module
are precisely the rows of \(A\): each generator of \(L\)
is a row times \(Y\), homogeneous of degree one over
\(S'\), and in such a product only the constant coefficient
of its multiplier contributes to degree one. Conversely
every \(S'\)-linear combination of the rows is an actual
relation in \(L\). This proves the presentation over
\(S'\). The map \(H\to B_1\) factors through \(H/J\).
The identifications above therefore give
\[
 N_0\otimes_{H_0}B_1=(N/JN)\otimes_{H/J}B_1
                   =(N/JN)\otimes_{S'}B_1 .
\]
Right exactness of tensor product yields the claimed matrix
over \(B_1\), with the indicated transpose because the
relations were rows. Over a DVR its Smith invariant factors
give the free rank and its torsion length. The product of
the nonzero invariant factors generates the ideal of the
maximal nonzero minors, which gives exactly the valuation
formulas stated.
\end{proof}

Here the length is computed for the ordinary off-diagonal
corner at a specified specialization. Applying this formula
to the strict cohomology groups over the fields \(E_x\)
requires an integral comparison with that corner and
uniform estimates for its minors at the prescribed arithmetic
weights. The period comparison in Section~\ref{sec:missing}
also remains open.

\clearpage
\section{References}\label{nav:bibliography-VIII}
\begingroup\def\section*#1{}\begin{thebibliography}{99}\small
\bibitem{VIII-VI-VI-nt119} J. Newton and J. A. Thorne, \href{https://arxiv.org/pdf/1912.11265}{\emph{Adjoint Selmer groups of automorphic Galois representations of unitary type}}, J. Eur. Math. Soc. \textbf{25} (2023), 1919--1967, Theorem~A and Section~2.5.
\bibitem{VIII-VI-chenevier364} G. Chenevier, \emph{The \(p\)-adic analytic space of pseudocharacters of a profinite group and pseudorepresentations over arbitrary rings}, in \emph{Automorphic Forms and Galois Representations}, Vol.~1, Cambridge University Press, 2014, 221--285, Introduction and Theorem A. \href{https://gaetan.chenevier.perso.math.cnrs.fr/articles/determinants.pdf}{Original paper}.
\bibitem{VIII-VI-VI-pan40} L. Pan, \emph{The Fontaine--Mazur conjecture in the residually reducible case}, J. Amer. Math. Soc. \textbf{35} (2022), 1031--1169. \href{https://arxiv.org/abs/1901.07166}{Author's paper}.
\bibitem{VIII-VI-VI-carayol117} H. Carayol, \href{https://www.numdam.org/articles/10.24033/asens.1512/}{\emph{Sur les repr\'esentations \(\ell\)-adiques associ\'ees aux formes modulaires de Hilbert}}, Ann. Sci. \'Ecole Norm. Sup. (4) \textbf{19} (1986), 409--468.
\bibitem{VIII-VI-VI-saito40} T. Saito, \emph{Modular forms and \(p\)-adic Hodge theory}, Invent. Math. \textbf{129} (1997), 607--620. See also \href{https://www.ms.u-tokyo.ac.jp/~t-saito/pp/h.pdf}{\emph{Hilbert modular forms and \(p\)-adic Hodge theory}}, Compos. Math. \textbf{145} (2009), 1081--1113.
\bibitem{VIII-VI-VI-deligne40} P. Deligne, \emph{La conjecture de Weil. I}, Publ. Math. IH\'ES \textbf{43} (1974), 273--307. \href{https://www.numdam.org/item/PMIHES_1974__43__273_0/}{Purity over finite fields}.
\bibitem{VIII-VI-V-milne98} J. S. Milne, \emph{Arithmetic Duality Theorems}, 2nd ed., BookSurge, 2006. \href{https://www.jmilne.org/math/Books/ADTnot.pdf}{Part I, Sections 2 and 4}.
\bibitem{VIII-VI-I-ac123} J. Arthur and L. Clozel, \href{https://www.claymath.org/library/cw/arthur/pdf/30.pdf}{\emph{Simple Algebras, Base Change, and the Advanced Theory of the Trace Formula}}, Annals of Mathematics Studies 120, Princeton University Press, 1989, Chapter III.
\bibitem{VIII-VI-V-shin} S. W. Shin, \emph{Galois representations arising from some compact Shimura varieties}, Ann. of Math. \textbf{173} (2011), 1645--1741. \href{https://annals.math.princeton.edu/2011/173-3/p09}{Theorem 1.2}.
\bibitem{VIII-VI-lgp361} T. Barnet-Lamb, T. Gee, D. Geraghty and R. Taylor, \emph{Local-global compatibility for \(l=p\), II}, Ann. Sci. \'Ec. Norm. Sup\'er. (4) \textbf{47} (2014), 165--179, Theorem A. \href{https://www.numdam.org/articles/10.24033/asens.2212/}{Published paper}.
\bibitem{VIII-VI-V-franke} J. Franke, \emph{Harmonic analysis in weighted \(L_2\)-spaces}, Ann. Sci. Ecole Norm. Sup. \textbf{31} (1998), 181--279. \href{https://www.numdam.org/articles/10.1016/s0012-9593(98)80015-3/}{Paper}.
\bibitem{VIII-VI-II-fw13} B. Ferrero and L. C. Washington, \emph{The Iwasawa invariant \(\mu_p\) vanishes for abelian number fields}, Ann. of Math. \textbf{109} (1979), 377--395. \href{https://doi.org/10.2307/1971116}{Original paper}.
\bibitem{VIII-VI-II-iwasawa15} K. Iwasawa, \emph{On \(\ZZ_\ell\)-extensions of algebraic number fields}, Ann. of Math. \textbf{98} (1973), 246--326.
\bibitem{VIII-VI-geraghty363} D. Geraghty, \emph{Modularity lifting theorems for ordinary Galois representations}, Math. Ann. \textbf{373} (2019), 1341--1427, Section 2. \href{https://doi.org/10.1007/s00208-018-1742-4}{Published paper}. Proposition and lemma numbers refer to the \href{https://citeseerx.ist.psu.edu/document?doi=8faa2d675c94ae78d0ac8ebb3f9d4f69ede42d44\&repid=rep1\&type=pdf}{original manuscript}.
\end{thebibliography}\endgroup
\end{document}
